\chapter{开弦与 Dirichlet 膜}\label{chap:3}

\SourceParagraph{ch03-intro-001}{119}
本章从二维共形场论，简称 CFT，的角度开始讨论 Dirichlet 膜，并在不借助现代表示论的范围内尽可能深入。

\SourceParagraph{ch03-intro-002}{119}
作为准备，\BookSectionRef{sec:3.1}{3.1}将回顾量子力学与各种上同调理论的联系，包括 de Rham 上同调、Dolbeault 上同调，以及它们如何嵌入 Hodge 理论等。对于许多目的，从 CFT 得到的结构，可以视为这些已经十分清楚的理论的形变。

\SourceParagraph{ch03-intro-003}{119}
关于 CFT 的讨论从\BookSectionRef{sec:3.2}{3.2}开始，先简要介绍它在物理上如何定义，即将其视为二维量子场论的一个特殊情形。大部分讨论偏重概念，不过，我们将较详细地讨论自由场论，以便在\BookSectionRef{sec:3.2.3.6}{3.2.3.6}中给出 T 对偶的标准论证。

\SourceParagraph{ch03-intro-004}{119}
随后，在\BookSectionRef{sec:3.3}{3.3}中简要介绍超共形场论及其拓扑扭曲，相关内容也在\BookBibText{ref242}{MS1}的第12章与第13章讨论。这包括“A 模型”与“B 模型”的物理定义、手征环，以及 $\mathcal{N}=2$ 超共形代数的结构。由于这些定义并不以靶空间几何为基础，物理允许提出超出代数几何、微分几何等标准数学框架的猜想。不过，要与数学建立联系，就必须假设，这些模型还具有几何定义，也就是非线性 $\sigma$ 模型。\BookSectionRef{sec:3.4}{3.4}将详细回顾这一点。

\SourceParagraph{ch03-intro-005}{119}
接着讨论边界条件与开弦。仍然先在\BookSectionRef{sec:3.5}{3.5}中作一般的物理讨论，再在\BookSectionRef{sec:3.6}{3.6}中将注意力限制于作了拓扑扭曲的 A、B 模型中的边界条件。最后，解释它们与\BookChapterRef{1}中校准子流形的关系，并发展足够支撑\BookChapterRef{5}中更一般讨论的形式体系，包括与全纯向量丛及一点的结构层关联的边界条件。

\SourceParagraph{ch03-intro-006}{119}
除了\BookBibText{ref242}{MS1}之外，关于 $(2,2)$ 超共形场论的其他标准参考文献，包括文献\BookCite{ref192,ref119}以及文献\BookCite{ref394}的第19章。面向数学家的超对称导论，可参见文献\BookCite{ref153}。

\section{拓扑量子力学与上同调理论}\label{sec:3.1}

\SourceParagraph{ch03-3.1-001}{120}
局部量子场论在物理中的标准定义，与\BookChapterRef{2}中采用的 TFT 定义看起来很不相同；不过，我们可以借助后者来说明物理定义的由来。因此，再次将 QFT 定义为从一个几何范畴，即带边界流形的某个范畴，到复向量空间与线性映射所成范畴的函子。不过，现在要求流形带有 Riemann 度量。于是，几何范畴中的对象是带度量的闭 $(d-1)$ 维流形，态射是带度量的 $d$ 维流形，它们给出对象之间的配边。当然，相应的线性对象与态射，也会由这些度量数据来参数化。

\SourceParagraph{ch03-3.1-002}{120}
在进入场论以前，先简要回顾 $d=1$ 的情形。这就是超对称量子力学与 Hodge 理论之间众所周知的联系。此时，一个对象是零维流形，换言之，是有限点集。将对应于一个点的复向量空间记为 $\mathcal{H}$。最简单的态射是区间 $[0,t]$，它为每个 $t$ 给出 $\mathcal{H}$ 上的一个线性算符。缝合的一致性要求这些算符构成半群，写作
\begin{equation}\label{eq:3.1}
\exp(-tH)\colon\mathcal{H}\longrightarrow\mathcal{H}.
\end{equation}
通常可以令 $\mathcal{H}$ 为 Hilbert 空间，并要求这是一个强连续半群，其每个算符都是有界、自伴的；我们也作这样的选择。于是，可以定义生成这个半群的 Hamilton 算符 $H$，它是一个下有界的自伴算符，通常并非有界算符。

\SourceParagraph{ch03-3.1-003}{120}
为了得到 $d>1$ 时那些有趣配边的类似物，可以允许任意图作为态射。度量数据，就是为每条边指定一个正实数长度。这些图由区间与顶点粘合而成；基本顶点取为三价顶点，高价顶点则用乘法的迭代解释，相应三价分解中人工引入的内边取零长度极限。这里选取的三价顶点给出交换、结合的乘法；在无限维情形下，乘法定义在适当的稠密代数 $\mathcal{A}\subset\mathcal{H}$ 上：
\begin{equation}\label{eq:3.2}
V\colon\mathcal{A}\times\mathcal{A}\longrightarrow\mathcal{A}.
\end{equation}
这样，就得到一个带半群作用的代数。最后，为了得到 Frobenius 型结构，选择迹 $\theta\colon\mathcal{A}\to\mathbb{C}$，使它与内积相容：
\[
(a,b)=\theta(a^*b).
\]
在有限维情形下，可以直接取 $\mathcal{A}=\mathcal{H}$，得到通常的 Frobenius 代数。
\begin{figure}[H]
\centering
\BookImage{01_Numbered_Figures/Chapter_03/Ch03_Fig01_PDFp120_Bookp110.pdf}
\caption{量子力学与图。}\label{fig:3.1}
\end{figure}

\SourceParagraph{ch03-3.1-004}{121}
\begin{example*}
我们主要关心的例子，是带度量流形 $(M,g)$ 上的复值函数代数。先在紧无边界流形上取 $\mathcal{A}=C^\infty(M,\mathbb{C})$，其 Hilbert 完备化为 $\mathcal{H}=L^2(M,\mathbb{C})$；\BookEqRef{eq:3.2}{3.2}中的乘法，是光滑函数的逐点乘法，迹则是关于体积形式 $\operatorname{Vol}_g$ 的积分。最后，令\BookEqRef{eq:3.1}{3.1}为 $M$ 上热方程的演化算符：
\[
-\frac{\partial}{\partial t}\bigl(\exp(-tH)f\bigr)
=\Delta\exp(-tH)f,
\]
其中，$\Delta$ 是标量 Laplace 算符，其坐标表达式为
\[
\Delta=-\frac{1}{\sqrt{\det g}}\partial_i\bigl(\sqrt{\det g}\,g^{ij}\partial_j\bigr).
\]
显然，取 $H=\Delta$ 就满足这个方程。由于 Schrödinger 方程可以通过将 $t$ 取为纯虚数从同一方程得到，这通常称为“虚时间量子力学”\footnote{或者，类比量子场论，称为“Euclid 时间量子力学”。}。在非紧情形下，需要为乘法、积分和 Laplace 算符选定适当的定义域及自伴实现。这一构造并不要求 $M$ 满足 Ricci 平坦性等曲率条件；这类限制会在考察 $d\ge2$ 时出现。

\SourceParagraph{ch03-3.1-005}{121}
这个理论的物理可观测量，是 Laplace 算符的谱，以及其他算符的矩阵元。其他算符中最重要的，是与一个函数相乘的算符；换言之，给定 $f\in C^\infty(M,\mathbb{C})$，定义
\[
O_f\colon\mathcal{H}\longrightarrow\mathcal{H},\qquad\psi\longmapsto f\psi.
\]
在上述紧情形下，这些都是有界算符，显然组成一个交换、结合的代数，本质上与\BookEqRef{eq:3.2}{3.2}定义的代数相同。
\end{example*}

\SourceParagraph{ch03-3.1-006}{121}
\begin{example*}
另一个例子，是沿向量场 $v\in\Gamma(M,TM)$ 的微分算符，定义在光滑函数上：
\[
\Pi[v]=-\mathrm{i}v^I\frac{\partial}{\partial X^I}.
\]
在物理中，当 $v$ 是生成等距变换的 Killing 向量场时，这个算符称为“动量算符”。它给出一个“守恒荷”，意思是，由于
\[
[H,\Pi(v)]=0,
\]
这两个算符可以在相容的谱分解中同时对角化，所以动量本征值，或者简称动量，不随时间改变。
\end{example*}

\SourceParagraph{ch03-3.1-007}{121}
这些数据虽然与 $(M,g)$ 的某些较简单的不变量有联系，但这种联系十分微妙而复杂，例如可参见文献\BookCite{ref041}。为了得到与拓扑更直接相关的理论，我们来考察超对称量子力学。

\subsection{超对称量子力学}\label{sec:3.1.1}

\SourceParagraph{ch03-3.1.1-001}{121}
\begin{definition*}\label{def:3.supersymmetric-qm}
超对称量子力学，是具有 $N$ 个线性无关的 Hermitian 算符 $Q_I$ 的理论，这些算符满足
\begin{equation}\label{eq:3.3}
\{Q_I,Q_J\}=\delta_{IJ}H,
\end{equation}
其中，$H$ 是 Hamilton 算符。$N>1$ 的情形，有时称为“扩展超对称”。
\end{definition*}

\SourceParagraph{ch03-3.1.1-002}{122}
最基本的 $N=1$ 超对称量子力学，简称 SQM，就是 Dirac 方程的理论。取
\begin{equation}\label{eq:3.4}
\mathcal{H}\equiv\Gamma(M,S),
\end{equation}
即 $M$ 上旋量丛 $S$ 的截面空间；这里先写光滑截面，作为 Hilbert 空间时取其 $L^2$ 完备化，并假设已选定所需的旋结构。在每个点 $x\in M$ 上，纤维 $S_x$ 承载 Clifford 代数 $\operatorname{Cl}(T_xM)$ 的作用。这个代数的 $\dim M$ 个实生成元称为“费米子”，它们张成 $T_xM$。$\operatorname{Cl}(T_xM)$ 的定义关系为
\[
\psi^2=(\psi,\psi)\,1,\qquad\psi\in T_xM.
\]
在 $T_xM$ 中选取一组具体基 $e_i$，将对应于线性坐标的 Clifford 生成元记为 $\psi^i$，就可以把定义关系认作典范反对易关系：
\begin{equation}\label{eq:3.5}
\{\psi^i,\psi^j\}=2g^{ij}|_x.
\end{equation}
可以取超荷 $Q$ 为 Dirac 算符：
\[
Q=\not{D}=-\mathrm{i}\psi^kD_k,
\]
其中，$D_k$ 是协变导数。按照\BookEqRef{eq:3.3}{3.3}的归一化，取 $H=2\not{D}^{\,2}$ 就满足超对称关系；它是旋量截面上的 Dirac 型 Laplace 算符的两倍。

\SourceParagraph{ch03-3.1.1-003}{122}
另一个重要算符记为 $(-1)^F$。在偶维情形下，它由以下性质确定，至多相差一个整体符号：其本征值为 $\pm1$，并且与所有 $\psi^i$ 反对易，
\begin{equation}\label{eq:3.6}
(-1)^F\psi^i=-\psi^i(-1)^F.
\end{equation}
当 $\dim M$ 为偶数时，它的 $\pm1$ 本征空间将旋量丛 $S$ 分解成两个不可约的半旋量表示，
\begin{equation}\label{eq:3.7}
S\cong S_+\oplus S_-,
\end{equation}
Dirac 算符 $Q$ 将二者相互映射。在物理中，这些不可约旋量表示称为“手征”表示，而 $\pm1$ 本征值称为手征性。在奇维情形下，如需这样的分次算符，须使用具有额外 $\mathbb{Z}/2$ 分次的旋量模，例如将模加倍。

\subsubsection{\texorpdfstring{$N=2$}{N=2} 超对称}\label{sec:3.1.1.1}

\SourceParagraph{ch03-3.1.1.1-001}{122}
一个朴素的尝试，是取 $N$ 组费米子 $\psi_I^i$，让它们作用于若干旋量丛的张量积，借此构造具有 $N$ 扩展超对称的理论。但一般来说，这种尝试会失败，因为 $[D_i,D_j]\ne0$，不同的超荷因而不再反对易。不过，在 $N=2$ 这一特殊情形下，对上述构造作适当修改，就可以得到一个可行的理论。

\SourceParagraph{ch03-3.1.1.1-002}{122}
在 $N=2$ 的情形下，可以用一个复超荷 $Q=Q_1+\mathrm{i}Q_2$ 来表示超对称代数：
\begin{equation}\label{eq:3.8}
\{Q,Q^\dagger\}=2H,\qquad
\{Q,Q\}=\{Q^\dagger,Q^\dagger\}=0.
\end{equation}

\SourceParagraph{ch03-3.1.1.1-003}{123}
Clifford 代数的关系为
\[
\{\psi_1^i,\psi_1^j\}=\{\psi_2^i,\psi_2^j\}=2g^{ij},
\qquad\{\psi_1^i,\psi_2^j\}=0.
\]
引入复费米子会很有用：
\[
\mathrm{d}x^i\equiv\frac12(\psi_1^i+\mathrm{i}\psi_2^i),
\qquad
\iota_j\equiv\frac12g_{jk}(\psi_1^k-\mathrm{i}\psi_2^k).
\]
它们满足微分形式及其伴随算符的代数关系：
\begin{equation}\label{eq:3.9}
\{\mathrm{d}x^i,\iota_j\}=\delta^i_j,
\qquad\{\mathrm{d}x^i,\mathrm{d}x^j\}=\{\iota_i,\iota_j\}=0.
\end{equation}
因此，在这一情形下，可以取向量空间 $\mathcal{H}$ 为 $M$ 上复值微分形式的空间 $\Omega^*(M,\mathbb{C})$，同样在需要时取 $L^2$ 完备化。这个空间按形式次数具有 $\mathbb{Z}$ 分次，在物理中称为“费米子数”。这一分次由算符
\begin{equation}\label{eq:3.10}
F=\mathrm{d}x^i\iota_i
\end{equation}
的作用给出。注意，这与\BookEqRef{eq:3.6}{3.6}一致。

\SourceParagraph{ch03-3.1.1.1-004}{123}
现在，将算符 $Q$ 与 $Q^\dagger$ 认作
\[
Q=\mathrm{d}\equiv\mathrm{d}x^i\frac{\partial}{\partial x^i}
\]
以及
\[
Q^\dagger=\mathrm{d}^\dagger=(-1)^{mF+1}*\mathrm{d}*,
\]
其中，$*$ 是 Hodge 星算符，$m=\dim M$，前置的次数算符 $F$ 作用于输出形式。它们满足\BookEqRef{eq:3.8}{3.8}，此时 Hamilton 算符为微分形式上的 Hodge Laplace 算符的一半：
\begin{equation}\label{eq:3.11}
2H=\mathrm{d}\mathrm{d}^\dagger+\mathrm{d}^\dagger\mathrm{d}.
\end{equation}

\SourceParagraph{ch03-3.1.1.1-005}{123}
在物理文献中，类比其他场合，例如弦的协变量子化，算符 $Q$ 经常称为“BRST 算符”。有时也会把 $Q^\dagger$ 称为“B 鬼场”，
\[
B\equiv Q^\dagger.
\]
无论采用什么称呼，关键都在于
\[
Q^2=0,
\]
因而可以用 $Q$ 定义上同调理论：
\[
H_Q^*(M,\mathbb{C})\equiv\frac{\operatorname{Ker}Q}{\operatorname{Im}Q}.
\]
在当前情形中，$Q=\mathrm{d}$；由此可见，de Rham 复形在物理上来自 $N=2$ 超对称量子力学。

\SourceParagraph{ch03-3.1.1.1-006}{123}
还有一个标准论证，说明代数结构在上同调上给出分次乘法：若 $Qa=Qb=0$，且 $a,b$ 为齐次元素，$|a|=p$、$|x|=p-1$、$|y|=|b|-1$，则由分次 Leibniz 法则，
\[
\begin{aligned}
(a+Qx)(b+Qy)=ab+Q\Bigl(&(-1)^p ay+xb\\
&+\tfrac12(-1)^p(Qx)y+\tfrac12xQy\Bigr).
\end{aligned}
\]

\SourceParagraph{ch03-3.1.1.1-007}{123}
将 $N=2$ 超对称量子力学限制到 $Q$ 上同调所得的理论，有时称为“拓扑量子力学”；类比下文的场论情形，甚至也称为“拓扑扭曲后的理论”。这一名称并非因为理论编码了 de Rham 上同调，而是因为到线性范畴的 QFT 函子不依赖于图上选定的度量，因此在\BookChapterRef{2}的意义下是“拓扑的”。这是由于 Hamilton 算符\BookEqRef{eq:3.11}{3.11}在上同调上的诱导作用恒为零，所以时间演化\BookEqRef{eq:3.1}{3.1}是平凡的。

\SourceParagraph{ch03-3.1.1.1-008}{124}
虽然仅用 $M$ 上微分形式的代数，就能定义 $Q$，但要定义 $B$ 和 $H$，则需要作出额外选择，在这里就是一个度量。在这个意义上，超对称量子力学不仅包含 de Rham 上同调，还包含它在 Hodge 理论中的实现。按照 Hodge 理论中的标准结果\BookCite{ref197}，对于紧定向无边界流形 $M$，每个 $Q$ 上同调类都具有一个由选定度量确定的典范代表，即满足
\[
Qa=Ba=0
\]
的调和形式 $a$。这就建立了 $H$ 的零本征空间与 $Q$ 上同调之间的同构。这个结果从物理观点看很重要：它表明，在 $N=2$ SQM 中，基态空间只依赖于 $M$ 的拓扑，而不依赖于其 Riemann 度量。高维场论中也有类似结果\BookCite{ref464}。

\subsection{泛函积分方法}\label{sec:3.1.2}

\SourceParagraph{ch03-3.1.2-001}{124}
到目前为止，我们只是将 de Rham 理论与 Hodge 理论改写成更接近物理的语言。Witten 在20世纪80年代初首先认识到，只要从物理中借用更多思想，就能得到处理许多数学问题的一种强有力的新方法。

\SourceParagraph{ch03-3.1.2-002}{124}
这些发展中的大部分，都以 QFT 的泛函积分定义为关键。\BookBibText{ref242}{MS1}对这种定义作了启发性的介绍，并讨论了它的许多应用。这里仅针对量子力学作简要回顾。一个较近的严格讨论，见文献\BookCite{ref089}。

\SourceParagraph{ch03-3.1.2-003}{124}
这种方法能够处理的最简单问题，来自几何范畴中由闭环给出的态射。因为它没有边界，QFT 函子将它映为一个复数，称为“配分函数”，或者“单圈配分函数”，记为 $Z(t)$。给定\BookEqRef{eq:3.1}{3.1}的具体表示，在演化算符为迹类的情形下，可以将其计算为
\begin{equation}\label{eq:3.12}
Z(t)=\operatorname{Tr}\mathrm{e}^{-tH}
\equiv\sum_{i\in I}\langle i|\mathrm{e}^{-tH}|i\rangle,
\end{equation}
其中，求和遍历一组标准正交基的指标集 $I$。

\SourceParagraph{ch03-3.1.2-004}{124}
利用半群性质，将时间演化算符分解成 $n$ 个算符，每个算符作用的时间为 $t/n$，其中 $n$ 为正整数：
\[
Z(t)=\sum_{i_1,\ldots,i_n\in I}
\langle i_1|\mathrm{e}^{-tH/n}|i_2\rangle
\langle i_2|\mathrm{e}^{-tH/n}|i_3\rangle\cdots
\langle i_n|\mathrm{e}^{-tH/n}|i_1\rangle.
\]

\SourceParagraph{ch03-3.1.2-005}{125}
随后考察 $n\to\infty$ 的极限。可以预期，在这一极限中，核
\[
\langle i|\mathrm{e}^{-tH/n}|j\rangle
\]
会集中到对角线上；在适当的光滑定义域上，短时间展开为
\[
\mathrm{e}^{-tH/n}=1-\frac{t}{n}H+\text{高阶项}.
\]
对于热方程，这种短时间核的渐近性质可以证明，并能写出其具体形式。于是，配分函数可以写成具有已知被积函数的多重积分。对于靶空间为 $M$ 的量子力学，得到
\begin{equation}\label{eq:3.13}
Z(t)=\int_{M^n}\left(\prod_{i=1}^n\mathrm{d}x_i\right)
\mathrm{e}^{-S_n[x_i]},
\end{equation}
其中，$x_i\in M$，$\mathrm{d}x_i$ 表示相应体积测度；被积函数写成函数 $S_n[x_i]$ 的指数形式，这个函数称为“作用量”，并包含各短时间核的归一化因子。

\SourceParagraph{ch03-3.1.2-006}{125}
最后，可以尝试将 $n\to\infty$ 的极限看作这个积分的连续形式，即对连续“路径” $S^1\to M$ 作积分。这当然是数学上难以处理的一步，因为极限测度的性质并不完全显然。一般来说，它依赖于具体的 Hamilton 算符，因而也依赖于作用量泛函 $S$。

\SourceParagraph{ch03-3.1.2-007}{125}
在 $d=1$ 的情形下，这一步已经得到很好的理解。在当前情形中，$H$ 是 Riemann 流形上的 Laplace 算符；对于这一情形及其超对称推广，上述极限导向 Wiener 测度及其推广，其支撑上的典型路径连续而几乎处处不可微。借助这种解释，可以从路径积分得到严格的结论。下文将讨论 $d\ge2$ 的情形。

\SourceParagraph{ch03-3.1.2-008}{125}
沿着这些论证，在普通量子力学中，与上文 $H=\Delta$ 的归一化相应的形式作用量为
\[
S=\frac14\int_0^t\mathrm{d}s\,
\left\|\frac{\mathrm{d}x(s)}{\mathrm{d}s}\right\|^2,
\]
其中，$s$ 参数化 $S^1$，范数由 $M$ 上的度量定义。因此，采用上述短时间核极限和 Wiener 测度的解释，就得到\BookEqRef{eq:3.12}{3.12}的积分表示。

\SourceParagraph{ch03-3.1.2-009}{125}
在任意维数中，也可以使用收敛到带度量 $d$ 维流形的一列单纯逼近，作出类似的形式论证。这导向如下通常形式的泛函积分：
\begin{equation}\label{eq:3.14}
Z=\int[\mathrm{D}\phi]\,\mathrm{e}^{-S[\phi]},
\end{equation}
其中，“场” $\phi\colon\Sigma\to M$ 是从 $d$ 维流形 $\Sigma$ 到靶空间 $M$ 的映射。作用量具有如下形式，总体常数取决于归一化约定：
\begin{equation}\label{eq:3.15}
S=\int_\Sigma|\mathrm{d}\phi|^2,
\end{equation}
其中，$\phi^i$ 是 $M$ 上的局部坐标，而范数现在依赖于 $\Sigma$ 与 $M$ 两者的度量。下文讨论 $d=2$ 时，会给出更具体的表达式。

\subsubsection{泛函积分与缝合}\label{sec:3.1.2.1}

\begin{proof}[缝合性质的形式推导]
\SourceParagraph{ch03-3.1.2.1-001}{126}
原则上，泛函积分为 QFT 函子提供了一种直接定义，其中\BookSectionRef{sec:2.1}{2.1}的缝合定理显然成立。现在解释这一十分有力的观点；不过，要记住，目前尚缺乏一般而严格的泛函积分定义，无法将下面的断言完全精确化，这大大限制了它在数学上的实用性。

\SourceParagraph{ch03-3.1.2.1-002}{126}
作用量泛函 $S$ 的关键性质，是局部性：它是 $\Sigma$ 上某个函数的积分，这个函数由一点 $\sigma\in\Sigma$ 处的场 $\phi(\sigma)$ 及其有限阶导数构造而成。形式测度 $[\mathrm{D}\phi]$ 也具有类似的独立性：在离散化中，未加入作用量权重的各点场变量采用乘积测度，因此不同点的场乘积的联合积分相应分解。

\SourceParagraph{ch03-3.1.2.1-003}{126}
为了看清这为何蕴含缝合定理，先简要说明泛函积分如何定义 QFT 函子。在物理中，这称为转到典范表述。回顾一下，给定几何范畴中的一个态射 $\Sigma$，即带边界的 $d$ 维流形，以及与边界 $\partial\Sigma$ 关联的 Hilbert 空间 $\mathcal{H}$ 中的元素 $v$，函子应当给出一个数。显然，一个积分原则上会给出一个数；还需要解释的，是如何将 $v$ 的选择纳入其中。

\SourceParagraph{ch03-3.1.2.1-004}{126}
这可以通过选择边界条件来实现。对于闭流形 $\Sigma$，我们对“所有”场 $\phi\colon\Sigma\to M$ 作泛函积分\BookEqRef{eq:3.14}{3.14}；而对于带边界流形，则需要指定场在边界上的行为。一种做法，是在映射 $\phi\colon\partial\Sigma\to M$ 的空间上指定测度\footnote{更精确地说，为了进行下文的构造，这应当是测度的“平方根”，即“半密度”。} $\mu[\phi]$。预先看来，选择很多：从支撑在一个取值为 $M$ 中某个固定点的常值映射上的原子测度，到支撑在所有映射上的“自由”测度，都可以考虑。无论如何，这里存在一项选择。

\SourceParagraph{ch03-3.1.2.1-005}{126}
随后，定义条件路径积分测度 $[\mathrm{D}\phi\mid\mu]$，使其在边界上与 $\mu$ 一致，从而定义带边界的泛函积分：
\[
\int[\mathrm{D}\phi\mid\mu]\,F(\phi|_{\partial\Sigma})
=\int\mathrm{d}\mu[\phi]\,F(\phi).
\]
于是，
\begin{equation}\label{eq:3.16}
U_\Sigma[\mu]=\int[\mathrm{D}\phi\mid\mu]\,\mathrm{e}^{-S[\phi]}.
\end{equation}
如果能够将 Hilbert 空间 $\mathcal{H}$ 认作测度 $\mu$ 的某个线性空间，那么带边界泛函积分\BookEqRef{eq:3.16}{3.16}就以粘合性质显然成立的形式，给出 QFT 函子的态射 $U_\Sigma$。

\SourceParagraph{ch03-3.1.2.1-006}{127}
考察沿子流形 $\Gamma\subset\partial\Sigma_i$，$i=1,2$，的粘合分解 $\Sigma=\Sigma_1\cup_\Gamma\Sigma_2$。我们要证明，由积分\BookEqRef{eq:3.16}{3.16}产生的对应于 $\Sigma$ 的线性映射，可以通过将对应于 $\Sigma_i$ 的线性映射，沿共同边界所对应的分量作缩并来得到。在测度 $\mu[\phi]$ 的空间中，选取对应于标准正交基 $|i\rangle\in\mathcal{H}$ 的基 $\mu_i$，$i\in I$；在反向定向的共同边界上，取与之配对的对偶基 $\mu^i$。于是，
\begin{equation}\label{eq:3.17}
U_\Sigma=\sum_{i\in I}
U_{\Sigma_1}(\ldots,\mu_i)\,U_{\Sigma_2}(\ldots,\mu^i).
\end{equation}

\SourceParagraph{ch03-3.1.2.1-007}{127}
作用量泛函 $S$ 的局部性，意味着权重 $\exp(-S)$ 分解成两个因子的乘积：一个是与 $\Sigma_1$ 关联的权重，另一个是与 $\Sigma_2$ 关联的权重。局部的乘积测度构造还意味着，固定共同边界 $\Gamma$ 上的数据后，$\Sigma$ 上的泛函测度也分解为乘积测度。剩下的一步，是验证测度满足\BookEqRef{eq:3.17}{3.17}；换言之，将 $\Sigma$ 上的泛函测度沿基 $\mu_i$ 展开为
\[
[\mathrm{D}\phi]|_\Sigma
=\sum_{i\in I}[\mathrm{D}\phi\mid\mu_i]_{\Sigma_1}
[\mathrm{D}\phi\mid\mu^i]_{\Sigma_2}.
\]
如果对任意分解 $\Sigma=\Sigma_1\cup_\Gamma\Sigma_2$ 都成立，那么，同一个 $\Sigma$ 的不同分解给出相同结果，就蕴含缝合定理。

\SourceParagraph{ch03-3.1.2.1-008}{127}
对于量子力学，将这些论证精确化并不困难。把\BookEqRef{eq:3.1}{3.1}的核看作测度之间的配对，并取 $t\to0$ 的极限，就得到 $\mathcal{H}\cong L^2(M,\mathbb{C})$ 上的标准内积。不过，目前在 $d\ge2$ 时，类似构造只有在极特殊的情形，也就是可精确求解的理论中，才能严格完成。若能在更一般的情形下做到这一点，将是一个重大进展。
\end{proof}

\subsection{指标定理}\label{sec:3.1.3}

\SourceParagraph{ch03-3.1.3-001}{127}
下面重新推导并证明 Dirac 算符的 Atiyah--Singer 指标定理\BookCite{ref472,ref043}，这是 Witten 关于超对称量子力学的观点的一个极好例子。

\begin{proof}[超对称推导]
\SourceParagraph{ch03-3.1.3-002}{127}
量子力学中有一个标准结果：演化算符\BookEqRef{eq:3.1}{3.1}，因而配分函数，具有小 $t$ 展开。对于紧无边界的 $d$ 维流形，若所作用的丛的秩为 $r$，$H$ 的主部为 $a\Delta$，则
\[
Z(t)=\frac{r\operatorname{Vol}M}{(4\pi at)^{d/2}}
\bigl(1+O(t)\bigr).
\]
上文的标量热方程中，$r=a=1$；按照\BookEqRef{eq:3.3}{3.3}的旋量归一化，$r=\operatorname{rank}S$、$a=2$。为了从泛函积分看出这一点，注意到，在 $t\to0$ 的极限中，\BookEqRef{eq:3.13}{3.13}启发式地化为 $M$ 上的一个积分。当然，必须先取离散划分数 $n\to\infty$ 的极限，所以这还过于朴素；不过，可以证明，正确处理这一极限带来的修正，归结为泛函积分的鞍点逼近的首阶修正，换言之，是在环路空间的切空间上作 Gaussian 积分，从而得到 $t^{-d/2}$ 因子。

\SourceParagraph{ch03-3.1.3-003}{127}
这已经很有意思；如果还能找到一个与配分函数不同、不依赖于 $t$ 的量，就更有意义了。典型的例子，是紧偶维旋流形上 Dirac 算符的指标，它可以实现为 $N=1$ 超对称量子力学中的修改后的配分函数，
\begin{equation}\label{eq:3.18}
I(t)=\operatorname{Tr}\bigl(\mathrm{e}^{-tH}(-1)^F\bigr),
\end{equation}
其中，$(-1)^F$ 由\BookEqRef{eq:3.6}{3.6}定义。要看清这一点，首先注意，$t$ 的依赖来自 $\mathcal{H}$ 中 $H>0$ 的子空间。但是，由于 $H=2Q^2$，这个子空间中 $H$ 的任一本征向量 $v$，都会与另一个本征向量 $Qv$ 配对；二者具有相同的本征值，而 $(-1)^F$ 本征值相反，所以在\BookEqRef{eq:3.18}{3.18}中相互抵消。详细解释见\BookBibText{ref242}{MS1}第10.2节。最后，$(-1)^F$ 与 Dirac 算符反对易，使两种手征性的旋量带上相反的权重，从而
\[
I(t)=\dim\operatorname{ker}\not{D}_+
-\dim\operatorname{ker}\not{D}_-
\equiv\operatorname{ind}\not{D}_+.
\]
因此，$I(t)$ 必须等于其小 $t$ 渐近展开中不依赖于 $t$ 的首项。基本上沿用前面为 $Z(t)$ 概述的论证，可以通过 Gaussian 泛函积分来计算它，得到将\BookEqRef{eq:3.18}{3.18}写成一个局部密度的积分的表达式；这个密度由度量和联络的曲率构造而成，也就是 Atiyah--Singer 指标公式。

\SourceParagraph{ch03-3.1.3-004}{128}
通向同一结果的另一种路径积分方法，是局部化论证。相关介绍见文献\BookCite{ref047}及\BookBibText{ref242}{MS1}第9章。它从超对称量子力学的路径积分出发。这个积分有几种写法：既可以用从 $1|1$ 维超空间到 $M$ 的映射来写，也可以用分量来写。我们采用第二种方法，在 $M$ 上的局部坐标 $X^i$ 之外，引入沿映射 $X\colon\Sigma\to M$ 取值于切丛的费米场 $\psi^i$。作用量为
\[
S=\int_\Sigma\mathrm{d}t\,
\Bigl(\|\partial_tX\|^2+g_{ij}\psi^j(D_t\psi)^i\Bigr),
\]
这里省略了不影响指标的共同整体归一化常数。$D_t$ 是 $TM$ 上的协变导数，具体为
\[
(D_t\psi)^i=\partial_t\psi^i
+(\partial_tX^k)\Gamma^i_{jk}\psi^j,
\]
其中，$\Gamma^i_{jk}$ 是与度量 $g$ 相容的 Levi--Civita 联络。

\SourceParagraph{ch03-3.1.3-005}{128}
泛函积分的超对称性，来自作用量在如下奇的无穷小变量变换下的不变性：
\[
\delta X^i=\psi^i,\qquad\delta\psi^i=-\partial_tX^i.
\]
接着，可以论证，具有这种奇对称性的积分，局部化到该对称性的固定配置，也就是“保持超对称”的配置上。于是，泛函积分化为对这些超对称配置的积分，每个配置都带有一个 Gaussian 因子，或者称为“单圈”因子。

\SourceParagraph{ch03-3.1.3-006}{128}
在指标定理的证明中，这些超对称配置的玻色部分，就是满足 $\partial_tX^i=0$ 的常值环路，因此，上述论证非常直接地给出指标公式。更进一步，如果出现的超对称配置的模空间是有限维的，由于 Gaussian 泛函积分容易处理，这个论证就能相当直接地推广到 $d>1$。关于局部化在镜像对称中的应用，仍请参见\BookBibText{ref242}{MS1}。
\end{proof}

\subsection{Dolbeault 上同调与扩展超对称}\label{sec:3.1.4}

\SourceParagraph{ch03-3.1.4-001}{129}
得到扩展 $N=2$ 超对称的另一种方法，是对靶流形 $M$ 的几何施加限制。从 $N=1$ SQM 出发，考虑如下拟设：
\[
Q_2=-\mathrm{i}J^j_k\psi^kD_j,
\]
其中，$J^j_k$ 是张量场。考察当 $I=1,2$ 时，关系\BookEqRef{eq:3.3}{3.3}对 $J$ 施加的限制，可以得到两个代数条件，
\[
g^{ik}J^j_k+g^{jk}J^i_k=0
\]
以及
\[
J^k_iJ^j_k=-\delta^j_i,
\]
还要求 $J$ 可积，即满足微分条件
\[
N^i_{jk}(J)\equiv
J^l_j\partial_lJ^i_k-J^l_k\partial_lJ^i_j
-J^i_l\partial_jJ^l_k+J^i_l\partial_kJ^l_j=0.
\]
上述代数条件说明，$J$ 是与度量相容的概复结构，而 $N(J)=0$ 等价于 $J$ 的可积性。对于这里使用 Levi--Civita 联络的无挠构造，还要求 $\nabla J=0$，也就是复结构平行；因此，$M$ 为 Kähler 流形时，就有这一额外超对称。这个 SQM 的态空间仍然是 $\Gamma(M,S)$，但也可以用微分形式写成
\[
\mathcal{H}=\Omega^{0,*}(M,\mathcal{K}^{1/2}),
\]
其中，$\mathcal{K}=\bigwedge^n(T^{1,0}M)^*$ 是 $M$ 的典范线丛，$n=\dim_\mathbb{C}M$；$\mathcal{K}^{1/2}$ 是所选旋结构对应的全纯平方根。

\SourceParagraph{ch03-3.1.4-002}{129}
另一种理解方法，是利用复化切丛的分解
\[
T_\mathbb{C}M=TM\otimes\mathbb{C}=T^{1,0}M\oplus T^{0,1}M,
\]
将 Clifford 代数 $\operatorname{Cl}(T_xM)$ 改写为\BookEqRef{eq:3.9}{3.9}的复类似物。与\BookEqRef{eq:3.5}{3.5}的归一化相容的识别为
\[
\frac{\psi^{\bar j}}{\sqrt2}\longmapsto\mathrm{d}\bar z^{\bar j},
\qquad
\frac{g_{i\bar j}\psi^i}{\sqrt2}\longmapsto\iota_{\bar j}.
\]
这里，$\mathrm{d}\bar z^{\bar j}$ 表示外乘算符。于是，在微分形式表述中，复超荷 $Q=Q_1+\mathrm{i}Q_2$ 成为
\[
Q=-2\mathrm{i}\psi^{\bar j}D_{\bar j}
=-2\sqrt2\,\mathrm{i}\,\mathrm{d}\bar z^{\bar j}D_{\bar j},
\]
其中，$D_{\bar j}$ 是全纯线丛 $\mathcal{K}^{1/2}$ 上的反全纯协变导数。由这个公式和 $[D_{\bar j},D_{\bar k}]=0$，立刻得到 $Q^2=0$。$Q$ 上同调就是 Dolbeault 上同调，
\[
H_Q(M)\cong H^{0,*}_{\bar\partial}(M,\mathcal{K}^{1/2}),
\]
而按照\BookEqRef{eq:3.8}{3.8}，$H=\tfrac12\{Q,Q^\dagger\}$ 仍是一个 Laplace 型算符。也可以定义 $\mathbb{Z}$ 分次，
\begin{equation}\label{eq:3.19}
R=\mathrm{d}\bar z^{\bar j}\iota_{\bar j},
\end{equation}
它就是反全纯形式的次数。

\SourceParagraph{ch03-3.1.4-003}{130}
可以将线丛 $\mathcal{K}^{1/2}$ 换成任意全纯线丛 $\mathcal{L}$。例如，可以取 $\mathcal{L}$ 为平凡线丛；这时，$Q$ 成为作用于 $(0,p)$ 型形式的通常 Dolbeault 算符的常数倍，而 $H$ 成为微分形式上的通常 Laplace 型算符。还可以进一步放宽 $M$ 为 Kähler 流形的条件，只要求 $M$ 是带 Hermitian 度量的复流形。为了在 $(0,p)$ 型形式上实现 $N=2$ 超对称代数，必要时再张量一个全纯线丛 $\mathcal{L}$，可以选择超荷的相位，并作相容的常数归一化：$Q_I\mapsto cQ_I$ 时，$H\mapsto c^2H$，其中 $c$ 为实常数。这样，取 $Q=\bar\partial$、$B=Q^\dagger$。按照\BookEqRef{eq:3.8}{3.8}，对应的 Hamilton 算符为 Dolbeault Laplace 算符的一半，
\[
H=\tfrac12\Delta_{\bar\partial},\qquad
\Delta_{\bar\partial}=\bar\partial\bar\partial^\dagger
+\bar\partial^\dagger\bar\partial.
\]
即使 $\mathcal{L}$ 平凡，它一般也不同于形式上的标准 de Rham Laplace 算符。在 Kähler 情形下，二者的关系为 $\Delta_{\mathrm{d}}=2\Delta_{\bar\partial}$\BookCite{ref197,ref249}。

\SourceParagraph{ch03-3.1.4-004}{130}
上述 $N=2$ 超对称代数的构造可以推广：若靶空间具有 $k$ 个线性无关、与度量相容且平行的复结构 $J_a$，并且满足
\[
J_aJ_b+J_bJ_a=-2\delta_{ab}\operatorname{id},
\]
就能得到 $k+1$ 个超对称。最有趣的额外情形，是 $M$ 为超 Kähler 流形，此时得到 $N=4$ 超对称。相关讨论见文献\BookCite{ref239}。

\subsubsection{扩展超对称的记号}\label{sec:3.1.4.1}

\SourceParagraph{ch03-3.1.4.1-001}{130}
我们已经讨论了两种具有 $N=2$ 超对称的系统。第一种以任意 Riemann 流形 $M$ 为基础，使用取值于 $T_\mathbb{C}M$ 的复费米子；第二种以 Kähler 流形为基础，使用取值于 $TM$ 的实费米子。为了区分 $N=2$ 超对称代数的这两种实现，分别称它们为 $N=(1,1)$ 模型与 $N=(2,0)$ 模型。这一术语来自二维场论，在那里，超荷可以具有正或负的旋量手征性：$N=(1,1)$ 模型可以通过维数约化，从每种手征性各有一个实超荷的二维场论得到；类似地，$N=(2,0)$ 模型来自具有两个相同手征性的实超荷的二维场论。类似地，具有 $N=4$ 超对称的量子力学模型，可以由具有 $N=(2,2)$ 超对称的二维模型约化得到，此时，在这里考察的无挠情形下，流形 $M$ 必须为 Kähler 流形，费米子取值于 $T_\mathbb{C}M$；也可以由具有 $N=(4,0)$ 超对称的二维模型约化得到，此时 $M$ 必须为超 Kähler 流形，费米子取值于 $TM$。

\subsubsection{R 对称}\label{sec:3.1.4.2}

\SourceParagraph{ch03-3.1.4.2-001}{130}
$N$ 扩展超对称代数\BookEqRef{eq:3.3}{3.3}，可以自然地与群 $\operatorname{SO}(N)$ 在超荷 $Q_I$ 上的线性作用，组成半直积。如果这个作用能够提升为 Hilbert 空间 $\mathcal{H}$ 上的作用，并保持量子力学的所有其他结构，例如代数\BookEqRef{eq:3.2}{3.2}，就称它为“理论的对称”。

\SourceParagraph{ch03-3.1.4.2-002}{131}
\BookStatementNumber{1}
\begin{definition}\label{def:3.1}
R 对称群，是由理论的对称组成的子群 $G_R\subseteq\operatorname{SO}(N)$。
\end{definition}

\SourceParagraph{ch03-3.1.4.2-003}{131}
\begin{example*}
$N=(1,1)$ 模型具有 $U(1)$ R 对称，其生成元为费米子数 $F$。$N=(2,0)$ 模型也具有 $U(1)$ R 对称，其生成元 $R$ 在玻色场 $\phi^i$ 上作用平凡，而在费米子 $\psi$ 上，通过复结构张量作用：
\[
R\cdot\psi=\mathrm{i}J\psi.
\]
$N=(2,2)$ 模型具有 $U(1)\times U(1)$ R 对称，而 $N=(4,0)$ 模型具有 $\operatorname{SU}(2)$ R 对称。
\end{example*}

\subsection{丛值上同调}\label{sec:3.1.5}

\SourceParagraph{ch03-3.1.5-001}{131}
通过选取一个结构群为 $G$、带联络 $\nabla$ 的向量丛 $V$，取
\[
\mathcal{H}=\Omega^*(M,V),
\]
并将所有导数适当地协变化，就可以推广 $N=2$ 模型。在一个局部标架 $e_I$ 中，可以用取值于 $V\otimes V^*$ 的联络一形式 $A$，来写协变微分 $D$：
\[
D=\mathrm{d}+A.
\]

\SourceParagraph{ch03-3.1.5-002}{131}
一般来说，超对称代数中会出现曲率项，
\[
Q^2\sim F.
\]
为了得到上同调理论以及正确的超对称代数，这一项必须为零。在 $N=(1,1)$ SQM 中，靶空间 $M$ 可以是任意 Riemann 流形，而 $Q=D$；此时 $D^2=F\wedge$，所以必须在 $V$ 上取平坦联络。所得的 $Q$ 上同调，同构于平坦向量丛 $V$ 的扭曲 de Rham 上同调；关于这一扭曲上同调的进一步讨论，见\BookSectionRef{sec:6.2.2}{6.2.2}。在 $N=(2,0)$ SQM 中，$M$ 是复流形，曲率项与 $F^{0,2}$ 成比例，因此，可以取满足 $F^{0,2}=0$ 的联络。所得的 $Q$ 上同调，是丛值 Dolbeault 上同调，
\[
H^{0,*}_{\bar\partial}(M,V).
\]
在上述紧无边界的情形下，仍然可以将 $Q$ 上同调与零能态空间识别，这对应于选取上同调类的调和代表。

\SourceParagraph{ch03-3.1.5-003}{131}
用物理语言说，这个 SQM 是具有“色”的粒子的量子力学。下面描述它的泛函积分定义，因为同一套形式体系，也用于定义开弦端点与 Dirichlet 膜上联络的耦合。考虑作为 $\operatorname{Hom}(\mathcal{H},\mathcal{H})$ 中线性算符的时间演化算符\BookEqRef{eq:3.1}{3.1}。在紧情形下，正时间的这一算符是 Hilbert--Schmidt 算符，它可对应于完备张量积 $\mathcal{H}\widehat\otimes\mathcal{H}^*$ 中的元素。相应的泛函积分核，在“初始”边界条件 $(\phi(0),e)$，$e\in V_{\phi(0)}$，与“末端”的对偶边界条件 $(\phi(t),e')$，$e'\in V_{\phi(t)}^*$，上取复数值；这里写出的是色自由度的边界标签。它由与\BookSectionRef{sec:3.1.2}{3.1.2}中讨论的类似的路径积分给出，但现在带上沿路径的 $\nabla$ 的平行移动作为权重，这是一个取矩阵值的路径泛函。在完整的超对称路径积分中，联络还通过曲率与本底费米子耦合；这里着重写出平行移动因子。将平行移动写成路径有序指数，是标准的做法；按照上面的 $D=\mathrm{d}+A$ 约定，这个联络因子写为
\[
\int[\mathrm{d}\phi]\,\mathrm{e}^{-S_0[\phi]}
\left(\mathcal{P}\exp\left(-\int\phi^*A\right)\right)^I_{\ J}.
\]

\SourceParagraph{ch03-3.1.5-004}{132}
有人可能质疑这个公式，因为这样的泛函积分看起来并非局部的。可以通过引入额外的量子力学自由度 $\Sigma_\lambda$，并用这些自由度的局部路径积分表示平行移动，来解决这一问题。为此，需要一个量子化后产生有限维 Hilbert 空间的系统，类似于费米子的量子化。事实上，最简单的方法，就是把一般情形嵌入一个费米子系统。为这一实现选择 $G$ 到某个 $\operatorname{SO}(N_c)$ 的正交嵌入，必要时扩大表示，使该嵌入具有所需的旋提升，且所要的 $G$ 表示出现在相应旋量模中。引入一个具有 $N_c$ 个生成元 $\lambda_I$ 的新 Grassmann 代数。量子化这个系统，就得到费米算符 $\lambda_I$，它们作用在 $\operatorname{Spin}(N_c)$ 的旋量表示中。将这个正交表示中的联络矩阵记为 $A_i^{IJ}$；在与前面一致的 Euclid 时间约定下，可以写出局部作用量
\[
S_\lambda=\frac12\int\mathrm{d}s\,
\left(\lambda_I\partial_s\lambda_I
+\partial_s\phi^iA_i^{IJ}(\phi)\lambda_I\lambda_J\right).
\]
第一项是动能项，第二项是在旋量表示中起作用的平行移动的无穷小形式。随后，可以将这个表示分解成 $G$ 表示，并限制到所需的那个表示；细节见文献\BookCite{ref334}。

\section{二维 QFT、CFT 与 TFT}\label{sec:3.2}

\SourceParagraph{ch03-3.2-intro-001}{132}
刚才回顾的所有 $d=1$ 理论，都有 $d=2$ 的推广。我们仍然从范畴框架出发，再逐渐转向更标准的物理方法。

\SourceParagraph{ch03-3.2-intro-002}{132}
首先，需要为闭理论选择一组简单的对象与态射。与通向泛函积分的论证一样，我们希望将态射看作由某些基本态射接合而成。例如，给定一个 $(d-1)$ 维流形 $Y$，自然可以考察带乘积度量的态射 $Y\times[0,t]$，其中 $t\in\mathbb{R}_+$。类似于\BookEqRef{eq:3.1}{3.1}，这在 Hilbert 空间 $\mathcal{H}_Y$ 上生成半群作用，同样称为 Euclid 时间演化。还可以将其写成一个自伴算符的指数；这个算符就是 Hamilton 算符 $H$，它隐含地依赖于 $Y$。

\SourceParagraph{ch03-3.2-intro-003}{132}
现在可以陈述 QFT 中的几条重要物理公理。严格地说，这里讨论的是幺正 QFT，因为物理学家也研究更一般的 QFT，其中“Hilbert 空间” $\mathcal{H}$ 带不定度量\footnote{例如，在弦理论的标准协变世界面量子化中，态空间带不定度量。这有两个来源：首先，靶时空 $M$ 带不定度量；其次，局部对称性的规范固定也会产生这一现象。}。不过，本书只讨论幺正 QFT，并在以下公理中选取具有唯一真空的情形；有限维谱子空间的条件针对紧空间截面。

\SourceParagraph{ch03-3.2-axiom-001}{133}
\begin{letteraxiom}{1}\label{axiom:3.1}
$\mathcal{H}_Y$ 是 Hilbert 空间。
\end{letteraxiom}

\SourceParagraph{ch03-3.2-axiom-002}{133}
\begin{letteraxiom}{2}\label{axiom:3.2}
$H$ 的谱下有界。$H$ 的本征值通常称为“能量”。
\end{letteraxiom}

\SourceParagraph{ch03-3.2-axiom-003}{133}
\begin{letteraxiom}{3}\label{axiom:3.3}
$H$ 的最低能量本征空间是一维的，称为“真空”或“基态”，通常记为 $|0\rangle$\footnotemark。
\end{letteraxiom}
\footnotetext{如果 $(d-1)$ 维流形是带平坦度量的 $\mathbb{R}^{d-1}$，习惯上将 Hamilton 算符 $H$ 定义为使真空能量为零。一般来说，真空能量以非平凡方式依赖于空间切片的度量，这称为 Casimir 效应。}

\SourceParagraph{ch03-3.2-axiom-004}{133}
\begin{letteraxiom}{4}\label{axiom:3.4}
对于任意 $E\in\mathbb{R}$，$\mathcal{H}_Y$ 中由能量不超过 $E$ 的态组成的谱子空间，是有限维的。
\end{letteraxiom}

\subsection{二维情形}\label{sec:3.2.1}

\SourceParagraph{ch03-3.2.1-001}{133}
现在取 $Y\cong S^1$。它的度量由一个实数，即周长 $\ell$，来参数化。将对应的复向量空间记为 $\mathcal{H}_\ell$。

\SourceParagraph{ch03-3.2.1-002}{133}
时间演化态射的半群，现在由带 $Y\times[0,t]$ 上乘积度量的圆环给出。当然，这些态射并不足以生成所有带度量的配边。

\SourceParagraph{ch03-3.2.1-003}{133}
考虑两个更大的族，就能得到一组足够的生成元。第一族态射，在拓扑上为 $Y\times[0,1]$，但配备更一般的度量 $g$，使“入”边界 $S^1$ 的周长为 $\ell$，“出”边界 $S^1$ 的周长为 $\ell'$。将这个态射记为 $D_g$。
\begin{figure}[H]
\centering
\BookImage{01_Numbered_Figures/Chapter_03/Ch03_Fig02_PDFp133_Bookp123.pdf}
\caption{二维 QFT 的态射。}\label{fig:3.2}
\end{figure}

\SourceParagraph{ch03-3.2.1-004}{133}
第二族从第一族得到：在一个标记点 $\mathrm{pt}$ 周围，挖去一个半径为 $r$ 的小圆盘。将所得态射记为 $O_g$，并将新产生的边界 $S^1$ 视为入边界。它的度量周长记为 $\ell_r$；在局部平坦圆盘的情形下，$\ell_r=2\pi r$。

\SourceParagraph{ch03-3.2.1-005}{133}
也许会认为，为了构造一般的 $\Sigma$，还需要另一族态射，也就是带一般度量的圆盘。稍后将论证，它们可以通过取 $D_g$ 的一个极限来得到：使一个边界缩成穿孔，并在该边界上放置一个确定的态，即真空。

\SourceParagraph{ch03-3.2.1-006}{134}
任何 $O_g$ 型态射，都可以用于定义乘法，
\begin{equation}\label{eq:3.20}
V_g\colon\mathcal{H}_\ell\times\mathcal{H}_{\ell_r}
\longrightarrow\mathcal{H}_{\ell'}.
\end{equation}
更具体地说，给定 $\mathcal{H}_\ell$ 与 $\mathcal{H}_{\ell_r}$ 中的态，用第一个态确定圆环入边界 $S^1$ 上的边界条件，用第二个态确定挖去圆盘后产生的 $S^1$ 上的边界条件。出边界上的态则属于 $\mathcal{H}_{\ell'}$。

\SourceParagraph{ch03-3.2.1-007}{134}
固定 $\phi\in\mathcal{H}_{\ell_r}$ 后，也可以将其看作定义了一个线性算符，
\begin{equation}\label{eq:3.21}
O_g[\phi]\colon\mathcal{H}_\ell\longrightarrow\mathcal{H}_{\ell'}.
\end{equation}
从 $\mathcal{H}_{\ell_r}$ 到这个线性算符空间的映射，称为态—算符对应。它的像，是局部算符组成的子空间。

\SourceParagraph{ch03-3.2.1-008}{134}
当然，$O_g$ 中的度量 $g$ 有许多可能的选择。由于可以附加态射 $D_{g'}$，使度量“增长”，自然会想到，将取 $g$ 的“零体积极限”所得的算符，作为局部算符的基本定义。更精确地说，可以在\BookEqRef{eq:3.20}{3.20}中取 $\ell=\ell'$，从 $Y\times[0,t]$ 上的乘积度量出发，在一个固定点 $\mathrm{pt}$，例如 $(x,t/2)$，处挖去一个圆盘，得到 $O_g$；然后，取 $t\to0$ 与 $r\to0$ 的极限，以定义 $O_{\mathrm{local}}[\phi]$。这一构造还需相应地识别不同 $\ell_r$ 的态，并在必要时对插入作归一化。
\begin{figure}[H]
\centering
\BookImage{01_Numbered_Figures/Chapter_03/Ch03_Fig03_PDFp134_Bookp124.pdf}
\caption{局部算符作为几何态射的极限。}\label{fig:3.3}
\end{figure}

\subsubsection{与量子力学的比较}\label{sec:3.2.1.1}

\SourceParagraph{ch03-3.2.1.1-001}{134}
给定这样的几何范畴，也许会设想，总能从中提取出\BookSectionRef{sec:3.1}{3.1}所定义的量子力学。做法是取 $Y\cong S^1$，使其周长 $\ell\to0$，并考察与“近似一维”的曲面 $\Sigma$ 关联的态射。例如，时间演化态射是乘积 $Y\times[0,t]$，其中，在 $\ell\to0$ 时保持 $t$ 不变。

\SourceParagraph{ch03-3.2.1.1-002}{134}
不深入细节，在适当的低能极限中，Hilbert 空间 $\mathcal{H}$ 的大部分态“被提升到无穷大能量”。换言之，作必要的真空能量平移并选定相应尺度后，时间演化会抑制这些高能态，投影到某个低能子空间 $\mathcal{H}'\subset\mathcal{H}$；这个子空间上仍可能保留非平凡的有效时间演化。对于剩下的空间，最简单的猜想，是它具有某个有限维流形 $M$ 上的 $L^2(M,\mathbb{C})$ 的形式。

\SourceParagraph{ch03-3.2.1.1-003}{134}
如果这个极限确实是一个量子力学理论，下一步就是用三价顶点\BookEqRef{eq:3.2}{3.2}定义一个交换、结合的代数，以此识别 $M$。随后，定义这个代数的谱，并证明它就是流形 $M$。能够做到这一点，正是以下观点的核心：量子力学并非一种“根本上全新”的数学结构，而是思考函数空间、流形上的度量等已知结构的另一种方法\footnote{当然，这种说法很不符合历史。这些数学结构的发展，曾深受关于量子力学的思考的影响。}。

\SourceParagraph{ch03-3.2.1.1-004}{135}
不过，如果尝试取\BookEqRef{eq:3.21}{3.21}的极限，以得到一个三价顶点，就会发现，所得算符 $O_{\mathrm{local}}[\phi]$ 是无界的，而且它们的乘积几乎总有奇性。因此，一般无法从这些局部算符的点乘积，直接得到量子力学中那种自然的交换、结合代数。

\SourceParagraph{ch03-3.2.1.1-005}{135}
这一障碍不仅是技术上的，而且触及 QFT 的基本结构。它正是我们迄今尚未找到类似 $C^0(M)\leftrightarrow M$ 那样简单关系的主要原因；对于紧 Hausdorff 空间 $M$，这里的 $C^0(M)$ 是连续函数组成的交换 $C^*$ 代数，这一关系正是前面量子力学讨论的基础。

\subsubsection{局部算符与算符乘积展开}\label{sec:3.2.1.2}

\SourceParagraph{ch03-3.2.1.2-001}{135}
要克服这一困难，需要对形式体系作相当多的发展。这里不作深入讨论，但至少要定义算符乘积展开，简称 OPE。对于某些 CFT，可以用顶点代数的形式体系，使这一构造严格化。

\SourceParagraph{ch03-3.2.1.2-002}{135}
如果接受粘合公理，那么，处理局部算符的一种更简单的方法，是将每个局部算符视为与 $\Sigma$ 上的一个穿孔关联，而不必先把曲面 $\Sigma$ 切成 $D_g$ 与 $O_g$，再取极限。因此，局部算符由一点 $p\in\Sigma$ 来参数化。此外，定义这种算符时，通常还要在 $p$ 的邻域选取坐标 $x$，于是，将算符写为 $\phi(x)$。

\SourceParagraph{ch03-3.2.1.2-003}{135}
随后，将局部算符乘积的关联函数，作为基本可观测量：
\begin{equation}\label{eq:3.22}
F(x_1,x_2,\ldots,x_n)
\equiv\bigl\langle\phi_1(x_1)\phi_2(x_2)\cdots\phi_n(x_n)\bigr\rangle_\Sigma.
\end{equation}

\SourceParagraph{ch03-3.2.1.2-004}{135}
这些关联函数可以有多种定义。在范畴语言中，如果能将 $\Sigma$ 从 $D_g$ 与 $O_g$ 组装出来，就有
\begin{equation}\label{eq:3.23}
\begin{aligned}
F(x_1,x_2,\ldots,x_n)
={}&\phi_1\cdot D_{g_{1,2}}O_{\mathrm{local}}[\phi_2]D_{g_{2,3}}\\
&\cdots D_{g_{n-2,n-1}}O_{\mathrm{local}}[\phi_{n-1}]
D_{g_{n-1,n}}\cdot\phi_n.
\end{aligned}
\end{equation}
这里，左端的 $\phi_1$ 按边界定向参与对偶配对。我们将这样的量称为未归一化关联函数。在泛函积分表述中，这些函数是关于泛函测度的平均：
\begin{equation}\label{eq:3.24}
F(x_1,x_2,\ldots,x_n)
=\int[\mathrm{d}\phi]\,\mathrm{e}^{-S[\phi]}
\phi_1(x_1)\phi_2(x_2)\cdots\phi_n(x_n).
\end{equation}

\SourceParagraph{ch03-3.2.1.2-005}{136}
更精确地说，“关联函数”一词更经常用于归一化的期望值，
\begin{equation}\label{eq:3.25}
F(x_1,x_2,\ldots,x_n)
=\frac1Z\int[\mathrm{d}\phi]\,\mathrm{e}^{-S[\phi]}
\phi_1(x_1)\phi_2(x_2)\cdots\phi_n(x_n),
\end{equation}
其中，$Z$ 由\BookEqRef{eq:3.14}{3.14}定义。对于具有正 Euclid 权重的玻色情形，$Z^{-1}[\mathrm{d}\phi]\mathrm{e}^{-S[\phi]}$ 是概率测度；在更一般的情形下，仍可以用上述归一化的泛函期望值。无论如何，给定一组算符后，关联函数都定义在 $\Sigma^n$ 中各个插入点两两不相重合的部分，即除去所有碰撞对角线之后的空间上。

\SourceParagraph{ch03-3.2.1.2-006}{136}
给定完整的一组关联函数，就能通过基于 Gel'fand--Naimark--Segal 构造的论证，重建 $\mathcal{H}$ 及其他范畴数据，因此，这也是 QFT 的一种同样有效的表述\footnote{为了做到这一点，这里讨论的 Euclid 关联函数必须满足 Osterwalder--Schrader 公理；关于这些公理及重建定理，见文献\BookCite{ref182}。}。

\SourceParagraph{ch03-3.2.1.2-007}{136}
\begin{definition*}\label{def:3.ope}
现在讨论局部算符的乘积。Wilson 首先提出，这种乘积可以写为局部算符的无穷级数展开，这就是算符乘积展开，或 OPE。在 CFT 中，沿固定方向考察其径向展开时，可写成
\begin{equation}\label{eq:3.26}
\phi_1(x_1)\phi_2(x_2)
=\sum_iC_{12,i}|x_1-x_2|^{\Delta_{12,i}}\phi_i(x_1),
\end{equation}
其中，求和包括相应的初级算符及后裔算符。在选定的理论中，系数 $C_{12,i}$ 与指数 $\Delta_{12,i}$ 只依赖于算符的选择；位置的径向依赖已经显式写出。一般的有自旋算符还带有相应的角向因子与指标缩并。对于非共形的 QFT，OPE 系数通常是位置差的更一般的函数。
\end{definition*}

\SourceParagraph{ch03-3.2.1.2-008}{136}
在满足前面有限能量谱条件的幺正 CFT 中，从尺度分次可以看出：一般会有一些 $\Delta_{12,i}<0$，它们反映了乘积的奇性，但发散项的个数是有限的。此外，最强的可能奇性，来自 $\phi_i=\phi_0$ 这一项，其中 $\phi_0$ 对应于真空 $|0\rangle$。真空在态—算符对应下的像，至多相差一个整体系数，就是恒等算符，因此也常写作 $\phi_0=1$。这一项出现时，$\Delta_{12,0}$ 由 $\phi_1$ 与 $\phi_2$ 的尺度权重，亦即相对于真空的相应 Hamilton 谱数据决定：若二者的尺度维数为 $\Delta_1,\Delta_2$，则 $\Delta_{12,0}=-\Delta_1-\Delta_2$。对于具有非零二点函数的一对共轭初级场，二者的权重相同；后裔场则不必满足这一相等条件。

\SourceParagraph{ch03-3.2.1.2-009}{136}
在\BookEqRef{eq:3.24}{3.24}中，各个算符的插入处于对称地位；交换费米算符时，需计入相应的分次符号。为了将\BookEqRef{eq:3.23}{3.23}写成具有这一性质的形式，可以将两端的“封帽”改为恒等算符 $\phi_0$，同时将原先位于封帽上的 $\phi_1$ 与 $\phi_n$ 保留为内部的局部插入。至多相差一个整体系数，这就等价于用圆盘态射，将两端分别封闭起来。

\SourceParagraph{ch03-3.2.1.2-010}{136}
反过来，给定泛函积分表述后，场的局部泛函会定义一组自然选定的局部算符，具体如下。令 $\Sigma$ 为带标准度量的圆盘，其边界是半径为 $r$ 的圆周。随后，对路径积分乘以在原点取值的局部泛函 $F$ 作为权重，就得到 $\mathcal{H}_{\ell_r}$ 中的一个元素，其中边界周长 $\ell_r=2\pi r$。这个元素线性地依赖于 $F$，所以，在某种意义上，这是态—算符对应的逆构造。

\SourceParagraph{ch03-3.2.1.2-011}{136}
粘合关系对 OPE 系数施加强有力的约束，通常称为“OPE 的结合律”；不过，由于位置也参与其中，这并不是通常的结合律。在下文讨论的某些简单 CFT 中，包括 $c\le1$ 的例子，指定相应的谱与算符内容后，这些约束可以实际求解，并在这样的特定情形中唯一确定理论。

\subsubsection{能动张量}\label{sec:3.2.1.3}

\SourceParagraph{ch03-3.2.1.3-001}{137}
\begin{definition*}\label{def:3.stress-tensor}
能动张量通常记为 $T$，是在这里考察的局部 QFT 中都存在的、最重要的局部算符。采用与后文\BookEqRef{eq:3.32}{3.32}中的 Virasoro 模式相容的 CFT 归一化，可以用泛函积分如下定义它；这里仍写未归一化关联函数，并固定其余插入：
\begin{equation}\label{eq:3.27}
\bigl\langle T^{\mu\nu}(x)\cdots\bigr\rangle
=-\frac{4\pi}{\sqrt{\det g(x)}}\,
\frac{\delta}{\delta g_{\mu\nu}(x)}
\int[\mathrm{D}\phi]_g\,\mathrm{e}^{-S[\phi,g]}\cdots.
\end{equation}
\end{definition*}
换言之，在关联函数中插入能动张量 $T_{\mu\nu}(x)$，就生成 $x$ 点处的一次无穷小度量形变；与其他插入相重合时，还需处理相应的接触项。另一方面，由定义\BookEqRef{eq:3.20}{3.20}及其与线性范畴的关系，$T_{\mu\nu}(x)$ 作为线性算符作用于 $\mathcal{H}_\ell$。如果知道这一作用，原则上就能为任意 $g$ 计算态射 $D_g$。

\SourceParagraph{ch03-3.2.1.3-002}{137}
\begin{example*}
考察 Euclid 时间演化\BookEqRef{eq:3.1}{3.1}，它由带乘积度量的态射 $Y\times[0,t]$ 定义。正如前面讨论的，它是 Hamilton 算符 $H$ 的指数。按照\BookEqRef{eq:3.27}{3.27}的归一化，这个 Hamilton 算符可以用能动张量定义为
\[
H=\frac1{2\pi}\int_Y\mathrm{d}\operatorname{Vol}_Y\,
T_{\mu\nu}v^\mu v^\nu,
\]
其中，$v^\mu$ 是时间方向，即 $[0,t]$ 因子方向的单位向量场。
\end{example*}

\SourceParagraph{ch03-3.2.1.3-003}{137}
因此，Hamilton 算符由能动张量决定。另一方面，能动张量是局部算符，所以受到更强的约束。

\subsection{二维共形场论}\label{sec:3.2.2}

\SourceParagraph{ch03-3.2.2-001}{137}
\begin{definition*}\label{def:3.cft}
现在可以给出 CFT 的定义：它是这样的 QFT，度量的共形重标度通过边界态的共轭变换起作用，并带上反常引出的普适标量因子。对于态射族 $D_g$，将出边界记为 $B_1$、入边界记为 $B_2$，可以将这一条件写为
\begin{equation}\label{eq:3.28}
D_{[\mathrm{e}^{h}g]}
=\mathrm{e}^{c\alpha[h]}\,
L^{-1}[h|_{B_1}]\,D_g\,L[h|_{B_2}].
\end{equation}
对于任意 $\Sigma$，也有类似的陈述，即对每个边界上的态作相应的共轭变换。
\end{definition*}

\SourceParagraph{ch03-3.2.2-002}{137}
这里，$L$ 是线性算符，只依赖于 $h$ 在圆环的某个边界上的限制。在固定背景度量后，远离边界的共形重标度的全部依赖，都由一个普适的、与具体 CFT 无关的泛函 $\alpha[h]\in\mathbb{R}$ 决定；它以整体乘法因子 $\mathrm{e}^{c\alpha[h]}$ 的形式出现。量 $c$ 称为“Virasoro 中心荷”，或者在本章中简称“中心荷”；稍后会给出更精确的定义。

\SourceParagraph{ch03-3.2.2-003}{137}
先考察整体重标度这一特殊情形，此时 $h$ 是常数。与量子力学中的讨论一样，将对数长度尺度作为参数后，对应的算符 $L[h]$ 构成半群，具有自伴生成元 $H_{\mathrm{dil}}$\footnote{它与下文\BookEqRef{eq:3.32}{3.32}中的 Virasoro 生成元的关系为 $H_{\mathrm{dil}}=L_0+\overline{L}_0$。圆周长度为 $\ell$ 时，乘积圆柱上的时间演化 Hamilton 算符为 $H_\ell=(2\pi/\ell)\bigl(L_0+\overline{L}_0-(c_L+c_R)/24\bigr)$；其中的常数项是圆柱的 Casimir 能量。}。由前面的下有界谱条件，可以在有限能量态等适当的共同稠密定义域上，将这一半群扩展为群作用；逆算符一般无界，不能仅由谱下有界就断言整个 Hilbert 空间上存在有界群作用。利用这些重标度，可以将任意周长对应的 $\mathcal{H}_\ell$ 与某个固定周长对应的空间识别，例如选取 $\ell=1$。以下作出这种识别，并直接使用较简单的记号 $\mathcal{H}\cong\mathcal{H}_1$。

\SourceParagraph{ch03-3.2.2-004}{138}
如何确定 $L[h]$？大致的做法如下。首先，对 $\Sigma$ 作单值化：换言之，找出从带边界曲面 $\Sigma$ 到一个常曲率曲面的共形微分同胚 $\phi$。随后，将 $\phi$ 限制到各个边界分支 $B_i$，得到 $\phi_i\in\operatorname{Diff}^+(S^1)\times\mathbb{R}_+$，其中，$\mathbb{R}_+$ 因子作整体重标度。接着，用圆周微分同胚群的 Lie 代数，即圆周向量场所生成的流，来表示每个 $\phi_i$；一般需要由随时间变化的向量场给出的路径有序指数，在确实可嵌入一个自治流的情形下，才可写成某个 Lie 代数元素 $l_i$ 的单一指数。最后，在 $\mathcal{H}$ 上寻找这个 Lie 代数的适当射影表示。

\SourceParagraph{ch03-3.2.2-005}{138}
实际进行这一构造，会遇到若干十分重要的微妙之处。正确理解它们，就能得到这些理论中大部分精确结果。首先，这里出现的圆周向量场 Lie 代数，实际嵌入两个彼此对易的代数的直和中；这两个代数分别作用于 $\mathcal{H}$ 中的“左行”与“右行”因子\footnote{也称为“全纯”与“反全纯”，或者“手征”与“反手征”。}。因此，在这里的离散谱情形下，可以将 $\mathcal{H}$ 写成这个直和代数的不可约表示的 Hilbert 直和：
\begin{equation}\label{eq:3.29}
\mathcal{H}=\widehat{\bigoplus_i}
\mathcal{H}_{L,i}\widehat\otimes\mathcal{H}_{R,i}.
\end{equation}
指标 $i$ 也计入可能的重数。这两个彼此对易的代数，分别是圆周向量场 Lie 代数的中心扩张，通常称为 Virasoro 代数，或简记为 $\operatorname{Vir}$。

\SourceParagraph{ch03-3.2.2-006}{138}
在讨论这一代数的表示论以前，先说明共形不变性如何使给定关联函数中的 OPE\BookEqRef{eq:3.26}{3.26}具有非零的局部收敛半径，直到碰到其他插入点。考察含有两个算符 $\phi_1$ 与 $\phi_2$ 的关联函数，并选取它们的位置，使某个圆周恰好包围二者，而不包围任何其他算符。在任意二维 QFT 中，都可以在这个圆周上定义 Hilbert 空间 $\mathcal{H}$；但在 CFT 中，还能将它重标度得任意小，使 $\mathcal{H}$ 中的态再次与局部算符对应。在关联函数中使用这一展开时，应保持其他插入位于包围圆周之外。

\SourceParagraph{ch03-3.2.2-007}{138}
迭代这一过程，就会发现，任何有限个局部算符的乘积所产生的态，都对应于相应的局部算符展开。在这个意义下，CFT 的态—算符对应是一个同构；有限能量态对应通常的局部场，而一般 Hilbert 态还需要相应的完备化解释。

\subsubsection{Virasoro 表示论的约束}\label{sec:3.2.2.1}

\SourceParagraph{ch03-3.2.2.1-001}{138}
考察圆周微分同胚群 $\operatorname{Diff}(S^1)$ 对 $S^1$ 上函数的自然作用，用 $\theta\in[0,2\pi)$ 参数化圆周。复化后，先考察由以下 Fourier 模式生成的向量场子代数：
\begin{equation}\label{eq:3.30}
l_n=\mathrm{i}\,\mathrm{e}^{\mathrm{i}n\theta}
\frac{\partial}{\partial\theta},\qquad n\in\mathbb{Z}.
\end{equation}
它们满足
\begin{equation}\label{eq:3.31}
[l_m,l_n]=(m-n)l_{m+n}.
\end{equation}
\begin{definition*}\label{def:3.vir}
Virasoro 代数是这一 Witt 代数的泛中心扩张。其生成元为 $L_n$，$n\in\mathbb{Z}$，以及一个中心元；在中心元作用为实数 $c$ 的表示中，关系为
\begin{equation}\label{eq:3.32}
[L_m,L_n]=(m-n)L_{m+n}
+\frac{c}{12}m(m^2-1)\delta_{m+n,0}\,1.
\end{equation}
参数 $c$ 仍称为 Virasoro 中心荷。
\end{definition*}

\SourceParagraph{ch03-3.2.2.1-002}{139}
容易说明，对于本书讨论的左右中心荷相等的幺正 CFT，非平凡情形必须具有非零的中心扩张。下面概述这一论证；我们并不需要它的全部细节，但它可以为后面关于 $N=2$ 超共形代数的类似而重要的论证作准备。首先，幺正性与其他 QFT 公理要求 Virasoro 表示作用在 Hilbert 空间上，并且 $L_{-n}=L_n^\dagger$。特别地，$L_0$ 自伴；在这里采用的离散有限能量谱条件下，可以将其对角化。前面的\BookRef{axiom:3.2}{公理2}要求正能量表示，考察其不可约组成部分，就得到最高权表示。
\begin{definition*}\label{def:3.highest-weight}
在这种组成部分中，$L_0$ 的最小本征值记为 $h$，相应的本征向量称为“最高权态”\footnotemark。将它记为 $|h\rangle$，则
\begin{equation}\label{eq:3.33}
L_0|h\rangle=h|h\rangle,
\end{equation}
并将其归一化为 $\langle h|h\rangle=1$。
\end{definition*}
\footnotetext{这里名称与本征值大小的倒置，是标准约定。}

\begin{proof}[幺正性给出的约束]
\SourceParagraph{ch03-3.2.2.1-003}{139}
作为准备，考察
\[
\langle h|[L_1,L_{-1}]|h\rangle
=\langle h|2L_0|h\rangle=2h.
\]
由于 $|h\rangle$ 是最高权态，可以证明 $L_1|h\rangle=0$，否则就会出现更低的 $L_0$ 本征值。因此，上式也等于
\[
\langle h|L_1L_{-1}|h\rangle
=\|L_{-1}|h\rangle\|^2,
\]
这里用了 $L_1=L_{-1}^\dagger$。最后，由于这是 Hilbert 空间中的范数，得到 $h\ge0$，而等号只能在 $L_{-1}|h\rangle=0$ 时成立。这样，不仅与\BookRef{axiom:3.2}{公理2}的下有界谱要求相容，还得到关于 $h=0$ 的真空的一些信息。事实上，$L_{-1}|0\rangle=0$ 可以与真空的平移不变性联系起来，后者是另一条物理公理。

\SourceParagraph{ch03-3.2.2.1-004}{139}
证明非平凡情形下 $c>0$ 的论证相同，只需考察
\[
\langle0|[L_2,L_{-2}]|0\rangle
=\frac c2
=\|L_{-2}|0\rangle\|^2\ge0.
\]
等号只能在 $L_{-2}|0\rangle=0$ 时成立。在这里讨论的幺正 CFT 假设下，还可以证明，等号意味着所有 $L_n|0\rangle=0$，从而 CFT 完全平凡。沿着这些论证继续发展，就得到不可约最高权 Virasoro 表示的结构由 $h$ 与 $c$ 这两个数决定。
\end{proof}

\SourceParagraph{ch03-3.2.2.1-005}{139}
用局部算符来改述上述讨论，会比用态来表述更方便。取 $\Sigma$ 为无限圆柱 $\mathbb{R}\times S^1$；也可以通过共形变换，使用带复坐标 $z$ 的穿孔复平面 $\mathbb{C}^*$。在 CFT 中，平面能动张量的分量 $T_{zz}$，可以用 Virasoro 生成元表示为
\[
T_{zz}\equiv T(z)=\sum_{n\in\mathbb{Z}}L_nz^{-n-2}.
\]
分量 $T_{\bar z\bar z}$ 是反全纯的，类似地，可以用第二个 Virasoro 代数的生成元 $\overline L_n$ 表示：
\[
T_{\bar z\bar z}\equiv\overline T(\bar z)
=\sum_{n\in\mathbb{Z}}\overline L_n\bar z^{-n-2}.
\]

\SourceParagraph{ch03-3.2.2.1-006}{140}
混合分量 $T_{z\bar z}=T_{\bar zz}$ 是一个标量数乘恒等算符，即物理中所谓的 $c$ 数；对平坦度量，它在离开接触点时为零。对应于 $T(z)$ 的态，是 $L_{-2}|0\rangle$。

\subsubsection{初级场}\label{sec:3.2.2.2}

\SourceParagraph{ch03-3.2.2.2-001}{140}
\begin{definition*}\label{def:3.primary-field}
对应于最高权态 $|h\rangle$ 的局部算符，称为共形初级场，并满足
\[
T(z)O_h(w,\bar w)
=\frac{h}{(z-w)^2}O_h(w,\bar w)
+\frac1{z-w}\partial_wO_h(w,\bar w)+\cdots,
\]
其中，省略号表示在 $z\to w$ 时保持有限的项。
\end{definition*}
这是算符乘积展开的一个特殊例子，说明，当插入点彼此接近时，两个局部算符的乘积通常没有良好的极限。从这一公式可以推出，与态 $(L_{-1})^n|h\rangle$ 对应的局部算符，具有 $\partial_w^nO_h$ 的形式；一般的 $L_{-n}|h\rangle$ 则是 Virasoro 后裔态，并非都能写成这样的普通高阶导数。

\SourceParagraph{ch03-3.2.2.2-002}{140}
CFT 的共形对称代数，是彼此对易的“左”与“右”Virasoro 代数的直和，分别记为 $L_L$ 与 $L_R$。它们的不可约最高权张量积表示，由 $(h_L,c_L,h_R,c_R)$ 的取值确定。在这里讨论的所有理论中，$c_L=c_R$。此外，对于降到拓扑理论的态，在下文\BookSectionRef{sec:3.3.4.1}{3.3.4.1}的 Ramond 基态表述中，左右权重分别为 $c_L/24$ 与 $c_R/24$，因而有 $h_L=h_R$。通过谱流与这些态对应的 NS 算符，要与这一表述区分：其未扭曲的左右权重不必相等。

\subsubsection{缝合与因子分解}\label{sec:3.2.2.3}

\SourceParagraph{ch03-3.2.2.3-001}{140}
对于无边界的 $\Sigma$，几何函子赋予它一个数 $Z[\Sigma]$，称为配分函数。虽然 $Z[\Sigma]$ 依赖于 $\Sigma$ 上的度量 $g$，但它在共形变换 $g\to\alpha g$ 下的变化，由\BookEqRef{eq:3.28}{3.28}决定，其中 $\alpha\in C^\infty(\Sigma,\mathbb{R}_+)$；另一方面，它以非平凡方式依赖于 $\Sigma$ 的复结构。复结构由有限个参数刻画，因此，在固定共形因子的处理约定后，配分函数可以看作模空间上的函数。

\SourceParagraph{ch03-3.2.2.3-002}{140}
“因子分解”提供了一个重建图景：人们期望，所有亏格曲面的配分函数组成的一组数据，原则上足以确定 CFT。其基础，是 Riemann 曲面的复结构模空间在加入退化曲面之后的边界，与曲面退化成带穿孔的较低亏格曲面的极限有关。具体的重建须保持退化展开与缝合配对的兼容性，并识别其中的态及其可能的重数。

\SourceParagraph{ch03-3.2.2.3-003}{140}
\begin{proof}[长颈的缝合展开]
考察 $Z[\Sigma]$ 在这样一个边界附近的行为：$\Sigma$ 逐渐产生一条长颈，最终断成两个曲面，每个曲面各带一个穿孔。这样的退化，可以用一个“长度—扭转”参数 $\tau$ 来参数化。若颈长为 $a$、周长为 $\ell$，并约定扭转在态上作用为 $\mathrm{e}^{\mathrm{i}\vartheta(L_0-\overline L_0)}$，则取
\[
\tau=\frac{2\pi a}{\ell}-\mathrm{i}\vartheta.
\]
因此，它的实部是颈长与周长之比的 $2\pi$ 倍，虚部则记录旋转角，符号由上述扭转方向的约定确定。

\SourceParagraph{ch03-3.2.2.3-004}{140}
利用\BookRef{axiom:3.4}{公理4}及缝合公理，$Z[\Sigma]$ 具有如下展开：
\[
Z[\Sigma]
=\sum_i C_i\,
\mathrm{e}^{-\tau(h_{L,i}-c_L/24)
-\bar\tau(h_{R,i}-c_R/24)}.
\]
将已知的圆柱真空能量因子提取出来，即定义
\[
\widehat Z
=\mathrm{e}^{-(c_L\tau+c_R\bar\tau)/24}Z[\Sigma],
\]
就得到
\[
\widehat Z=\sum_iC_i\mathrm{e}^{-\tau h_{L,i}-\bar\tau h_{R,i}}.
\]
各项系数 $C_i$，可以识别为在穿孔处插入权重为 $(h_{L,i},h_{R,i})$ 的算符所得关联函数的组合；简并的态通过缝合配对求和。
\end{proof}

\SourceParagraph{ch03-3.2.2.3-005}{140}
在这一重建图景中，期望通过同时取多次退化极限，从配分函数确定全部关联数据；如果这种恢复能够完成，就可以利用前面的态—算符对应，重建几何范畴上的 CFT 函子。若谱中有简并，退化系数给出的是相应配对下的组合，重建还须按这一配对处理态基的选择。

\subsubsection{按 Virasoro 中心荷 \texorpdfstring{$c$}{c} 分类}\label{sec:3.2.2.4}

\SourceParagraph{ch03-3.2.2.4-001}{141}
中心荷是 CFT 最重要的不变量。它类似于流形的维数 $d$；事实上，对于下文讨论的 $\sigma$ 模型，二者成比例。

\SourceParagraph{ch03-3.2.2.4-002}{141}
为了说明这一点，考察“态密度”的积分计数函数 $N(\lambda)$，它定义为 Hamilton 算符 $H$ 小于指定值 $\lambda\in\mathbb{R}$ 的本征值个数，并计入重数。它在大 $\lambda$ 时的渐近行为，受维数控制；对于紧无边界的 $d$ 维流形上的标量量子力学，若 $H=\Delta_g$，Weyl 定理给出
\begin{equation}\label{eq:3.34}
N(\lambda)\sim
\frac{\omega_d}{(2\pi)^d}\operatorname{Vol}_g(M)\,
\lambda^{d/2},\qquad d=\dim M,
\end{equation}
其中，$\omega_d$ 是 $\mathbb{R}^d$ 中单位球的体积；其主要幂次就是 $\lambda^{\dim M/2}$。在适用模变换与真空支配的手征计数中，取 $H=L_0$，则 Cardy 公式给出的主要指数为
\begin{equation}\label{eq:3.35}
N(h)=\exp\left(2\pi\sqrt{\frac{ch}{6}}+o(\sqrt h)\right).
\end{equation}
这里写的是主要指数的渐近式，前面的多项式因子包含在 $o(\sqrt h)$ 中。若模变换后的最低权重为 $h_{\min}\ne0$，应以有效中心荷 $c_{\mathrm{eff}}=c-24h_{\min}$ 替代 $c$。对于完整的左右理论，计数应使用 $L_0+\overline L_0$ 等总能量，不能把右行无限维因子的重数全部计入仅由 $L_0$ 限制的有限态数。

\SourceParagraph{ch03-3.2.2.4-003}{141}
对于任何有限维 $M$，正有效中心荷的 Cardy 增长都快于\BookEqRef{eq:3.34}{3.34}；但由于它仍是次指数增长，像配分函数\BookEqRef{eq:3.12}{3.12}这样的量，在相应有限态数的计数中，对于任意 $\operatorname{Re}t>0$ 都有定义。事实上，椭圆曲线上的完整 CFT 配分函数，在 $\operatorname{SL}(2,\mathbb{Z})$ 群作用下具有模性质；单独的手征特征标则可能彼此混合。这在物理上的原因，是它们来自以 $\Sigma$ 为椭圆曲线的泛函积分。

\SourceParagraph{ch03-3.2.2.4-004}{141}
Virasoro 代数的表示论已经得到很好的理解。那么，这对理解 CFT 能有多大帮助？可以从最高权表示的最简单图景，也就是“Verma 模”，看出基本情形。
\begin{definition*}\label{def:3.verma-module}
Verma 模从一个最高权态出发，施加负模式子代数的泛包络代数，并先不再对零范数后裔取商。利用代数关系，其元素都可以写成以下态的线性组合：
\begin{equation}\label{eq:3.36}
\prod_{i>0}L_{-i}^{N_i}|h\rangle.
\end{equation}
乘积按固定次序排列，例如从左到右令 $i$ 递减；$N_i$ 是非负整数，且只有有限个 $N_i$ 非零。
\end{definition*}

\SourceParagraph{ch03-3.2.2.4-005}{141}
如果这些态全都线性无关，配分函数\BookEqRef{eq:3.12}{3.12}对应的手征特征标，就由 Dedekind $\eta$ 函数给出。若 $q=\mathrm{e}^{2\pi\mathrm{i}\tau_{\mathrm{ell}}}$，其中 $\operatorname{Im}\tau_{\mathrm{ell}}>0$，则
\[
\operatorname{Tr}_{V_{h,c}}q^{L_0-c/24}
=q^{h-c/24}\prod_{n\ge1}(1-q^n)^{-1}
=\frac{q^{h-(c-1)/24}}{\eta(\tau_{\mathrm{ell}})}.
\]
因此，可以将其渐近态数与\BookEqRef{eq:3.35}{3.35}比较。如果增长率相当，就有希望将完整 Hilbert 空间分解成\BookEqRef{eq:3.29}{3.29}形式的有限和，表示论因而会施加强有力的约束；实际分解还要考虑零范数关系及可能的扩充手征代数。另一方面，如果\BookEqRef{eq:3.35}{3.35}的增长快得多，就不能指望作出这样的有限 Virasoro 分解。

\SourceParagraph{ch03-3.2.2.4-006}{141}
事实上，\BookEqRef{eq:3.36}{3.36}的渐近态数，对应于\BookEqRef{eq:3.35}{3.35}中的有效中心荷1，因此，一般理论分成两种情形。在这里的紧、幺正、离散谱的玻色 CFT 分类中，$c\le1$ 的表示论与物理论证，给出了完整的分类及相应的精确解：$c<1$ 时是“极小模型”，$c=1$ 时是“自由玻色子”及其轨形，包括特殊半径处的轨形理论。完整讨论见文献\BookCite{ref112}。

\subsubsection{\texorpdfstring{$c>1$}{c>1} CFT 的构造}\label{sec:3.2.2.5}

\SourceParagraph{ch03-3.2.2.5-001}{141}
此时，仅靠表示论，并不能得到很强的结果，需要借助其他定义。在镜像对称的研究中，通常使用三种定义。我们的主要方法，是非线性 $\sigma$ 模型；\BookSectionRef{sec:3.2.6}{3.2.6}及\BookSectionRef{sec:3.3.2}{3.3.2}将较详细地讨论它。

\SourceParagraph{ch03-3.2.2.5-002}{142}
第二种一般方法，是\BookBibText{ref242}{MS1}中经常使用的线性 $\sigma$ 模型，“Landau--Ginzburg 模型”是其中一个特殊情形。在物理上，这比非线性 $\sigma$ 模型的方法更有力；关于我们将要讨论的问题，许多早期结果都是用这一方法得到的。另一方面，运用它需要的物理技术，远远超出了本书能够容纳的范围。因此，除了\BookSectionRef{sec:3.2.6.3}{3.2.6.3}中的简要讨论，我们决定不依赖这一方法。

\SourceParagraph{ch03-3.2.2.5-003}{142}
最后，还可以留在代数方法的框架中，从 $c\le1$ 理论中选取适当的有理理论，考察它们的张量积及其他组合，得到“有理”CFT 这一子类中的模型。这里的有理性是相对于适当的扩充手征代数而言，并不一定意味着只需有限个 Virasoro 表示。虽然这是一类十分特殊的理论，却能对另外两种方法所采用的物理论证，提供独立的确认。\BookSectionRef{sec:5.6}{5.6}将较详细地讨论平坦空间的轨形；\BookSectionRef{sec:3.3.6}{3.3.6}则会简要介绍“Gepner 模型”的构造。

\subsection{自由玻色 CFT}\label{sec:3.2.3}

\SourceParagraph{ch03-3.2.3-intro-001}{142}
CFT 的最简单例子，也是每本教科书中的第一个例子，就是自由玻色子。现在描述这一理论：既把它作为一个具体例子，也因为其中许多结果，会在后面的讨论中发挥关键作用。

\SourceParagraph{ch03-3.2.3-intro-002}{142}
“自由玻色场”在形式上是随机映射 $\varphi\colon\Sigma\to M$，其中，$M$ 是带平坦 Riemann 度量 $g_{ij}$ 的流形；取 $g_{ij}$ 为常数实对称正定矩阵。严格的连续场实现中，坐标场一般是广义函数，而不是逐点定义的普通映射。在这种情形下，可以为路径积分\BookEqRef{eq:3.14}{3.14}及作用量\BookEqRef{eq:3.15}{3.15}给出精确的 Gaussian 场定义，见文献\BookCite{ref182}；这里作较为直观的物理描述。

\SourceParagraph{ch03-3.2.3-intro-003}{142}
在 $\Sigma$ 上选取复坐标 $z=x+\mathrm{i}y$，将作用量写为
\[
S=\frac1{2\pi}\int_\Sigma\mathrm{d}^2z\,
(\partial\varphi,\bar\partial\varphi),
\qquad\mathrm{d}^2z=2\,\mathrm{d}x\,\mathrm{d}y,
\]
其中，$\partial=\partial/\partial z$、$\bar\partial=\partial/\partial\bar z$，$(\, ,\,)$ 是 $TM$ 上实内积的复双线性延拓\footnote{等价地，
\[
S=\frac12\frac1{2\pi}\int\mathrm{d}x\,\mathrm{d}y\,
\bigl((\partial_x\varphi,\partial_x\varphi)
+(\partial_y\varphi,\partial_y\varphi)\bigr).
\]
前置的 $1/(2\pi)$ 是归一化约定，用来给 Green 函数\BookEqRef{eq:3.43}{3.43}选取一个简单的系数。}。完整的被积形式具有 $(1,1)$ 型，因此，作用量仅通过 $\Sigma$ 的复结构依赖于其度量；以它为基础的 QFT 因而是 CFT 的候选。当然，根据\BookEqRef{eq:3.28}{3.28}及前面关于非平凡理论中 $c>0$ 的论证，这种对共形因子的独立性，会在量子化时受到 Weyl 反常的修正。

\subsubsection{泛函积分表述}\label{sec:3.2.3.1}

\SourceParagraph{ch03-3.2.3.1-001}{142}
先从 $M\cong\mathbb{R}^d$ 出发。由于作用量是 $\varphi$ 的二次式，这个泛函积分是无穷维 Gaussian 积分。类比有限维 Gaussian 积分，就能理解它的基本特点。
\begin{example*}
取 $x\in\mathbb{R}^N$，$C$ 为实对称正定矩阵，则
\begin{equation}\label{eq:3.37}
Z[C]=\left(\frac{(2\pi)^N}{\det C}\right)^{1/2}
=\int_{\mathbb{R}^N}\mathrm{d}^Nx\,
\exp\left(-\frac12x^{\mathrm T}Cx\right).
\end{equation}
\end{example*}
对于场积分，相应的二次型算符由作用量的系数决定。将 $\Sigma$ 上的度量记为 $\gamma$，其非负标量 Laplace 算符记为 $\Delta_\Sigma$；在局部平坦坐标中，$\Delta_\Sigma=-4\partial\bar\partial$。取 $g$ 正交的靶空间坐标后，作用量的每个实分量可写为 $\tfrac12(\varphi,K\varphi)_{L^2}$，其中
\[
K=\frac{\Delta_\Sigma}{2\pi}.
\]
若每个 $L^2$ 正交模的测度取为 $\mathrm{d}a/\sqrt{2\pi}$，类比\BookEqRef{eq:3.37}{3.37}，就会希望将非零模的场积分写成 $(\det{}'K)^{-d/2}$。

\SourceParagraph{ch03-3.2.3.1-002}{143}
行列式在形式上，是完整本征值组的乘积：
\[
\det K\equiv\prod_i\lambda_i,\qquad
K\varphi_i=\lambda_i\varphi_i.
\]
当然，这是一个发散的乘积，不过，人们已经发展出多种适当的定义，使它实现\BookEqRef{eq:3.28}{3.28}中的变换律。一个很好的例子，是 $\zeta$ 函数正则化。这里不需要这些细节，相关讨论见文献\BookCite{ref144,ref182}。

\SourceParagraph{ch03-3.2.3.1-003}{143}
除了这些“紫外”发散，简称 UV 发散，还存在另一个问题：在闭连通的 $\Sigma$ 上，标量 Laplace 算符有零本征值，其本征函数是常值函数，因此，未去除零模的行列式为零。为了处理这个“零模”，将场分解为
\begin{equation}\label{eq:3.38}
\varphi(z)=\varphi_0+\widetilde\varphi(z),
\end{equation}
其中，$\varphi_0\in M$，且 $\int_\Sigma\widetilde\varphi\,\mathrm{d}\operatorname{Vol}_\gamma=0$。记 $\Sigma$ 的面积为 $A_\Sigma$。常值正交模为 $A_\Sigma^{-1/2}$，因而该模的系数是 $\sqrt{A_\Sigma}\varphi_0$。按照上面的逐模测度约定，测度分解可用任意泛函 $F$ 写为
\begin{equation}\label{eq:3.39}
\begin{aligned}
\int[\mathrm{D}\varphi]\,F[\varphi]
={}&\left(\frac{A_\Sigma}{2\pi}\right)^{d/2}
\int_M\mathrm{d}\operatorname{Vol}_g(\varphi_0)\\
&\times\int[\mathrm{D}\widetilde\varphi]\,
F[\varphi_0+\widetilde\varphi].
\end{aligned}
\end{equation}
因此，可以写出形式表达式
\[
Z=\left(\frac{A_\Sigma}{2\pi}\right)^{d/2}
(\det{}'K)^{-d/2}\operatorname{Vol}_g(M),
\qquad
\operatorname{Vol}_g(M)=\int_M\sqrt{\det g}\,\mathrm{d}^d\varphi_0,
\]
其中，带撇号的行列式只对非零本征值取乘积：
\[
\det{}'K\equiv\prod_{\lambda_i\ne0}\lambda_i.
\]
对于当前的非紧靶空间 $\mathbb{R}^d$，靶空间体积因子无穷大，须固定零模或取每单位靶空间体积的配分函数；不能将这个完整零模积分当作有限量。

\SourceParagraph{ch03-3.2.3.1-004}{143}
在分离上述体积因子并作正则化后，这一表达式给出 $\Sigma$ 上度量的实值泛函。利用\BookSectionRef{sec:3.2.2}{3.2.2}中的思想，或者在每个共形类中选取规定面积的常曲率代表，就可以将它化为 $\Sigma$ 的复结构模参数的函数。这个函数可以用自守函数显式写出，见文献\BookCite{ref012,ref011}。

\SourceParagraph{ch03-3.2.3.1-005}{143}
接着讨论\BookEqRef{eq:3.25}{3.25}所定义的归一化关联函数。它们是矩阵 Gaussian 积分\BookEqRef{eq:3.37}{3.37}中的期望值的无穷维类似物。
\begin{example*}
例如，取 $v_1,v_2\in(\mathbb{R}^N)^*$，有
\[
\langle(v_1\cdot x)(v_2\cdot x)\rangle
\equiv\frac1{Z[C]}\int_{\mathbb{R}^N}\mathrm{d}^Nx\,
\mathrm{e}^{-\frac12x^{\mathrm T}Cx}
(v_1\cdot x)(v_2\cdot x).
\]
这些期望值，可以很容易地通过对以下生成函数求导得到：
\begin{align}
Z[C,j]&\equiv\int_{\mathbb{R}^N}\mathrm{d}^Nx\,
\exp\left(-\frac12x^{\mathrm T}Cx+j\cdot x\right),
\label{eq:3.40}\\
&=Z[C]\exp\left(\frac12j^{\mathrm T}C^{-1}j\right).
\label{eq:3.41}
\end{align}
具体为
\begin{equation}\label{eq:3.42}
\langle(v_1\cdot x)(v_2\cdot x)\cdots(v_n\cdot x)\rangle
=\left.\frac1{Z[C]}
\left(\prod_{a=1}^n v_a\cdot\frac{\partial}{\partial j}\right)
Z[C,j]\right|_{j=0}.
\end{equation}
\end{example*}

\SourceParagraph{ch03-3.2.3.1-006}{144}
于是，任意 Gaussian 关联函数，都可以用“Green 函数”，也就是二次型算符 $K$ 的形式逆，来表示。分离闭曲面上的零模后，其方程为
\[
K_zG(z,z')=\delta_\gamma(z,z')-\frac1{A_\Sigma},
\qquad
\int_\Sigma G(z,z')\,\mathrm{d}\operatorname{Vol}_\gamma(z)=0,
\]
其中，$K_z$ 作用于第一个变量，$\delta_\gamma$ 关于体积测度归一化。在二维平面上，选择一个零模或红外归一化约定后，局部 Green 函数满足 $K_zG(z,z')=\delta^{(2)}(z-z')$。对于一个单位靶空间度量的实分量，直接计算得到
\begin{equation}\label{eq:3.43}
\langle\varphi(z_1)\varphi(z_2)\rangle
=G(z_1,z_2)=-\frac12\log|z_1-z_2|^2,
\end{equation}
这里允许一个由所选红外约定确定的加性常数；在一般紧 Riemann 曲面上，应使用相应的全局 Green 函数，而不是直接把平面上的对数当作全局表达式。

\SourceParagraph{ch03-3.2.3.1-007}{144}
这一形式体系，是微扰量子场论的起点，每本相关教科书都会发展它。
\begin{example*}
另外两个例子，是
\begin{equation}\label{eq:3.44}
\langle\partial\varphi(z_1)\partial\varphi(z_2)\rangle
=-\frac1{2(z_1-z_2)^2}
\end{equation}
以及
\begin{equation}\label{eq:3.45}
\begin{aligned}
&\langle\partial\varphi(z_1)\partial\varphi(z_2)
\partial\varphi(z_3)\partial\varphi(z_4)\rangle\\
&\quad=\frac14\left(
\frac1{(z_1-z_2)^2(z_3-z_4)^2}
+\frac1{(z_1-z_3)^2(z_2-z_4)^2}
+\frac1{(z_1-z_4)^2(z_2-z_3)^2}\right).
\end{aligned}
\end{equation}
\end{example*}

\subsubsection{顶点代数表述}\label{sec:3.2.3.2}

\SourceParagraph{ch03-3.2.3.2-001}{144}
利用典范量子化这一标准物理框架，可以导出\BookSectionRef{sec:3.2}{3.2}中的 Hilbert 空间 $\mathcal{H}$ 及 Hamilton 算符 $H$。关于这一方法，见文献\BookCite{ref112}或其他量子场论教科书。这里要做的，是定义 Heisenberg 顶点代数这一最简单的非平凡例子，再将它与刚才用泛函积分得到的结果比较。

\SourceParagraph{ch03-3.2.3.2-002}{144}
顶点代数的公理，见文献\BookCite{ref157,ref158,ref246}。对于前面自由玻色子的相应实现，与\BookEqRef{eq:3.44}{3.44}相容的归一化流是
\[
J(z)=\mathrm{i}\sqrt2\,\partial\varphi(z),
\qquad J(z)J(w)\sim\frac1{(z-w)^2}.
\]
这一 $U(1)$ 流代数，在数学文献中也称为 Heisenberg 代数\footnote{用物理中的术语说，这是无穷组彼此独立的 Heisenberg 振子代数，它们共享恒等中心元；也称为二维的典范对易关系代数，简称 CCR。}。
\begin{definition*}\label{def:3.heisenberg-current}
引入生成元 $\alpha_n$，$n\in\mathbb{Z}$，满足
\begin{equation}\label{eq:3.46}
[\alpha_m,\alpha_n]=m\delta_{m+n,0}\,1.
\end{equation}
注意，生成元 $\alpha_0$ 属于中心。
\end{definition*}

\SourceParagraph{ch03-3.2.3.2-003}{144}
这一代数的不可约最高权 Fock 表示，由一个参数 $p$ 确定。在幺正情形下，$p$ 为实数。引入最高权态 $|p\rangle$，由以下条件定义：
\[
\alpha_n|p\rangle=0\quad(n>0),\qquad
\alpha_0|p\rangle=p|p\rangle.
\]
$n>0$ 的算符 $\alpha_n$ 称为“湮灭算符”，因为它们湮灭最高权态；$n<0$ 的算符则称为“产生算符”，因为它们产生新的向量。

\SourceParagraph{ch03-3.2.3.2-004}{144}
随后，将这一表示定义为在 $|p\rangle$ 上施加泛包络代数，并对上述关系取商所得的线性空间，这就是“Fock 空间”。按照 $\alpha_n^\dagger=\alpha_{-n}$ 赋予它内积后，可以取 Hilbert 完备化。这一空间的一组基，由多重指标 $N$ 标记；它是只有有限个非零项的非负整数无穷序列：
\begin{equation}\label{eq:3.47}
\alpha_{-n}^{N_n}\alpha_{-n+1}^{N_{n-1}}
\cdots\alpha_{-1}^{N_1}|p\rangle.
\end{equation}
整数 $N_i$ 通常称为“占据数”。

\SourceParagraph{ch03-3.2.3.2-005}{145}
通过以下定义，可以与前面的自由玻色子定义建立联系：
\begin{equation}\label{eq:3.48}
\frac{\partial}{\partial z}\varphi(z)
=\partial\varphi(z)
=-\frac{\mathrm{i}}{\sqrt2}\,J(z)
=-\frac{\mathrm{i}}{\sqrt2}
\sum_{n\in\mathbb{Z}}z^{-n-1}\alpha_n.
\end{equation}
这里写的是实场的全纯部分。这一式子可以积分，以得到 $\varphi(z)$，但这里不需要这样做。

\SourceParagraph{ch03-3.2.3.2-006}{145}
\begin{example*}
用代数方法核实\BookEqRef{eq:3.44}{3.44}。
\begin{solution}
先取 $|z_1|>|z_2|$。在 $p=0$ 的 Fock 真空中，
\[
\begin{aligned}
\langle0|\partial\varphi(z_1)\partial\varphi(z_2)|0\rangle
&=-\frac12\sum_{n\ge1}n\,z_1^{-n-1}z_2^{n-1}\\
&=-\frac1{2(z_1-z_2)^2}.
\end{aligned}
\]
由此，得到与泛函积分一致的二点函数。
\end{solution}
\end{example*}

\subsubsection{能动张量与 Virasoro 代数}\label{sec:3.2.3.3}

\SourceParagraph{ch03-3.2.3.3-001}{145}
现在说明，如何用这一形式体系，为自由玻色子推导\BookEqRef{eq:3.32}{3.32}，并确定中心荷 $c$。类似但更长的计算，足以导出 $N=2$ 超共形代数，并验证\BookSectionRef{sec:3.3.2}{3.3.2}中使用的结构结论。

\SourceParagraph{ch03-3.2.3.3-002}{145}
将 $S$ 当作经典泛函，应用定义\BookEqRef{eq:3.27}{3.27}，就会发现，能动张量的各个分量，都是 $\partial\varphi$ 与 $\bar\partial\varphi$ 的二次式。其中一个分量纯粹全纯：
\[
T_{zz}(z)=-(\partial\varphi(z))^2.
\]
不过，从\BookEqRef{eq:3.44}{3.44}可以看出，这一表达式并不能直接在量子理论中定义，因为这样的局部算符乘积会发散。

\SourceParagraph{ch03-3.2.3.3-003}{145}
一个自然的处理办法，是减去发散部分。\BookEqRef{eq:3.44}{3.44}的形式，提示定义
\begin{equation}\label{eq:3.49}
\begin{aligned}
T(z)&=-\lim_{z'\to z}
\left(\partial\varphi(z)\partial\varphi(z')
+\frac1{2(z-z')^2}\right)\\
&=\frac12\lim_{z'\to z}
\left(J(z)J(z')-\frac1{(z-z')^2}\right).
\end{aligned}
\end{equation}
\begin{definition*}\label{def:3.normal-ordering}
与此相关的代数操作，是“正规排序”，用冒号表示，例如
\[
T(z)=-:\partial\varphi(z)\partial\varphi(z):
=\frac12:J(z)J(z):.
\]
其定义是：先取算符的乘积，作模式展开\BookEqRef{eq:3.48}{3.48}，再重新排列，使所有产生算符位于所有湮灭算符的左侧。
\end{definition*}

\SourceParagraph{ch03-3.2.3.3-004}{145}
在当前情形中，这两个定义等价。更一般地，可以通过取 OPE 的正则部分，并令两点重合，来定义算符的正规乘积：
\begin{equation}\label{eq:3.50}
\begin{aligned}
\partial\varphi(z_1)\partial\varphi(z_2)
&=-\frac1{2(z_1-z_2)^2}
+:\partial\varphi(z_1)\partial\varphi(z_2):\\
&=-\frac1{2(z_1-z_2)^2}
+:\partial\varphi(z)\partial\varphi(z):
+O(z_1-z_2),
\end{aligned}
\end{equation}
其中，第二行取 $z_2=z$，表示在 $z_1\to z$ 时展开正则部分。这一操作虽然有用，但舍去奇异项的代价，是所得局部正规乘积一般不满足普通的结合律。

\begin{proof}[模式关系与最低权重]
\SourceParagraph{ch03-3.2.3.3-005}{146}
用这些规定中的任何一种，都可以计算能动张量的模式展开，
\[
T(z)=\sum_{n\in\mathbb{Z}}L_nz^{-n-2},
\]
并从代数上验证\BookEqRef{eq:3.32}{3.32}，确定其中的常数 $c$。得到 $c=1$，且
\[
L_n=\frac12\sum_{m\in\mathbb{Z}}:\alpha_m\alpha_{n-m}:.
\]
特别地，
\[
L_0=\frac12\alpha_0^2
+\sum_{m\ge1}\alpha_{-m}\alpha_m.
\]
利用这一表达式，直接计算就能看出，在最高权态上，
\begin{equation}\label{eq:3.51}
L_0|p\rangle=\frac{p^2}{2}|p\rangle.
\end{equation}

\SourceParagraph{ch03-3.2.3.3-006}{146}
此外，$\alpha_0$ 属于中心，而产生算符只能提高 $h$，所以，对于每个本征态，
\begin{equation}\label{eq:3.52}
h=\frac{p^2}{2}+N\ge\frac{p^2}{2},
\qquad N\in\mathbb{Z}_{\ge0}.
\end{equation}
这里，模式能级是 $N=\sum_{m\ge1}mN_m$。
\end{proof}

\subsubsection{完整理论与分次}\label{sec:3.2.3.4}

\SourceParagraph{ch03-3.2.3.4-001}{146}
自由玻色子的完整理论，通过取这个代数的左行与右行副本的张量积来定义。最高权表示由两个“荷”，或者称为“动量”，$(p_L,p_R)$ 标记，记为
\[
V_{p_L,p_R}=\mathcal{F}_{p_L}^{L}
\widehat\otimes\mathcal{F}_{p_R}^{R}.
\]
在定义理论时，剩下的选择，是具体采用哪些不可约表示及其重数。对于离散的荷谱，可以写成 Hilbert 直和
\begin{equation}\label{eq:3.53}
\mathcal{H}
=\widehat{\bigoplus}_{p_L,p_R}
\mathbb{C}^{N_{p_L,p_R}}\widehat\otimes V_{p_L,p_R},
\end{equation}
其中，$N_{p_L,p_R}$ 是非负整数重数；连续荷谱则需用相应的 Hilbert 直积分来表述。

\SourceParagraph{ch03-3.2.3.4-002}{146}
荷 $(p_L,p_R)$ 在 $\mathcal{H}$ 上定义双分次，OPE 保持这一分次，这称为“荷守恒”。因此，在这一自由场理论中，\BookEqRef{eq:3.53}{3.53}中出现的荷组成的集合，记为 $Q\subset\mathbb{R}^2$，须在加法下封闭。

\SourceParagraph{ch03-3.2.3.4-003}{146}
\begin{example*}
最简单的例子，是靶空间 $M\cong\mathbb{R}$ 的理论。它取所有 $p_L=p_R\in\mathbb{R}$，每个荷的重数为1；此时，分解是直积分。更精确地，$\mathcal{H}$ 包含 $L^2(M,\mathbb{C})$ 因子，其 Fourier 分解给出这一连续分解：
\[
\mathcal{H}\cong
\int_{\mathbb{R}}^\oplus V_{k/\sqrt2,k/\sqrt2}\,\mathrm{d}k.
\]
这里，$k$ 是实场 $\varphi$ 的物理 Fourier 动量，而 $p=k/\sqrt2$ 是归一化流的荷；其余因子是左右振子的 Fock 空间。
\end{example*}

\SourceParagraph{ch03-3.2.3.4-004}{147}
荷守恒有一个简单的论证，使用由\BookEqRef{eq:3.48}{3.48}得到的关系
\[
\alpha_0=\frac1{2\pi\mathrm{i}}\oint\mathrm{d}z\,J(z)
=\frac{\sqrt2}{2\pi}\oint\mathrm{d}z\,\partial\varphi(z).
\]
因此，可以将任意一组局部算符产生的态的荷 $p_L$，写成归一化流 $J$ 沿包围这些插入的轮廓的积分。全纯性与留数公式随后给出加法性。虽然这里不作这一构造，还可以定义局部算符，通常称为“顶点算符”，使它们在不同 $(p_L,p_R)$ 的表示之间起交织作用。

\subsubsection{紧化玻色子}\label{sec:3.2.3.5}

\SourceParagraph{ch03-3.2.3.5-001}{147}
现在转向周长为 $2\pi R$ 的 $M\cong S^1$。为了定义泛函积分，需要修改分解\BookEqRef{eq:3.38}{3.38}，加入具有非平凡缠绕数的映射所对应的项。一般的 $S^1$ 值映射，可以分解成一个调和缠绕代表与平均值为零的实值映射；为了使分解唯一，需要固定调和代表的常值部分。调和映射到 $S^1$ 的微分，是 $H^1(\Sigma,\mathbb{R})$ 中周期量子化的元素。取 $H^1(\Sigma,\mathbb{Z})$ 一组整基的调和代表 $\omega_i$，令 $f_i\colon\Sigma\to\mathbb{R}/\mathbb{Z}$ 满足 $\mathrm{d}f_i=\omega_i$，并在一个选定基点处取 $f_i=0$。于是，可以写为
\begin{equation}\label{eq:3.54}
\varphi(z)=\varphi_0+\widetilde\varphi(z)
+2\pi R\sum_i m_i f_i(z)\pmod{2\pi R},
\qquad m\in\mathbb{Z}^{b_1(\Sigma)}.
\end{equation}
这里，$\varphi_0\in S^1$，且 $\int_\Sigma\widetilde\varphi\,\mathrm{d}\operatorname{Vol}_\gamma=0$。调和代表与实值涨落的微分，在 $L^2$ 内积下正交。

\SourceParagraph{ch03-3.2.3.5-002}{147}
泛函测度现在成为对这些缠绕扇区的求和。沿用\BookEqRef{eq:3.39}{3.39}的逐模归一化，得到
\begin{equation}\label{eq:3.55}
\begin{aligned}
\int[\mathrm{D}\varphi]\mathrm{e}^{-S[\varphi]}
={}&\left(\frac{A_\Sigma}{2\pi}\right)^{1/2}
\int_{S^1}\mathrm{d}\varphi_0
\int[\mathrm{D}\widetilde\varphi]\,
\mathrm{e}^{-S[\widetilde\varphi]}\\
&\times\sum_{m\in\mathbb{Z}^{b_1(\Sigma)}}
\exp\left(-\pi R^2\sum_{i,j}m_i m_j
\int_\Sigma\omega_i\wedge*\omega_j\right).
\end{aligned}
\end{equation}
计算这个和，就会发现，配分函数中的新因子，是由 $\Sigma$ 的周期数据确定的 theta 函数；细节见文献\BookCite{ref112}。

\SourceParagraph{ch03-3.2.3.5-003}{147}
\BookSectionRef{sec:3.2.3.2}{3.2.3.2}的讨论仍然适用，直到需要指定不可约表示之和\BookEqRef{eq:3.53}{3.53}为止。可以通过以下比较来确定这一分解：一方面，用上面的方法计算模参数为 $\tau$ 的环面配分函数，其中 $\operatorname{Im}\tau>0$；另一方面，将它与 Hilbert 空间上的迹比较，
\begin{equation}\label{eq:3.56}
Z(\tau)=\operatorname{Tr}_{\mathcal{H}}
\exp\left(2\pi\mathrm{i}\tau(L_{0,L}-1/24)
-2\pi\mathrm{i}\bar\tau(L_{0,R}-1/24)\right).
\end{equation}
这里，左右中心荷均为1。

\SourceParagraph{ch03-3.2.3.5-004}{147}
还有一种确定 $Z(\tau)$ 的代数方法，沿用\BookChapterRef{2}的思路：把缝合约束写成\BookEqRef{eq:3.53}{3.53}中的重数数据所满足的条件，再求这些条件的一般解。对于前面 $c\le1$ 的分类系列，这一方案能够作出完整处理；较高中心荷时，通常需要选择更具体的构造。缝合约束隐含的一个结果，是配分函数在 $\Sigma$ 的“大”微分同胚，即不能连续地与恒等映射相连的微分同胚，作用下保持不变。对于 $T^2$，保定向的映射类由 $\operatorname{SL}(2,\mathbb{Z})$ 分类，这一不变性通常称为模不变性。

\SourceParagraph{ch03-3.2.3.5-005}{148}
最简单的例子，是要求\BookEqRef{eq:3.56}{3.56}在 $\tau\to\tau+1$ 下不变。由于左右 Casimir 平移相同，这要求所有态都满足
\[
h_L-h_R\in\mathbb{Z}.
\]
由\BookEqRef{eq:3.51}{3.51}以及整数振子能级，得到
\begin{equation}\label{eq:3.57}
p_L^2-p_R^2\in2\mathbb{Z}.
\end{equation}
另一种表述，也可以推广到 $d$ 维：将\BookEqRef{eq:3.57}{3.57}的左端看作荷群 $Q$ 上符号型为 $(1,1)$ 的二次型。此时，方程说明 $Q$ 必须是偶格。进一步利用\BookEqRef{eq:3.56}{3.56}的显式表达式及 Poisson 求和，可以证明，$\tau\to-1/\tau$ 下的模不变性，要求这个满秩格自对偶。

\SourceParagraph{ch03-3.2.3.5-006}{148}
因此，$Q$ 是符号型为 $(1,1)$ 的偶自对偶格。可取参考格
\[
U=\mathbb{Z}\frac{(1,1)}{\sqrt2}
+\mathbb{Z}\frac{(1,-1)}{\sqrt2},
\]
其整基的 Gram 矩阵为 $\left(\begin{smallmatrix}0&1\\1&0\end{smallmatrix}\right)$。在实平面中的任意这样的格嵌入，可以从这一参考格，通过保持二次型的线性变换得到；在选定的连通分支中，它写为
\[
(p_L,p_R)\longmapsto
(p_L\cosh\alpha+p_R\sinh\alpha,\,
p_L\sinh\alpha+p_R\cosh\alpha),
\qquad\alpha\in\mathbb{R}.
\]
具体的圆周理论，其荷为
\[
(p_L,p_R)=\frac1{\sqrt2}
\left(\frac nR+wR,\frac nR-wR\right),
\qquad n,w\in\mathbb{Z}.
\]
它的范数为 $2nw$，正好满足偶性条件。在上述参考基的选择下，$R=\mathrm{e}^{-\alpha}$。因此，紧自由玻色子理论由一个参数 $\alpha$ 刻画，与定义泛函积分时引入的单个半径参数 $R$ 一致\footnote{$c=1$ 理论的完整分类也有讨论，见文献\BookCite{ref112}。其他理论通过对离散对称性取轨形获得，方法与后文讨论 $(2,2)$ 理论时所用的思路类似。}。

\enlargethispage{0.5pt}
\subsubsection{T 对偶：泛函表述}\label{sec:3.2.3.6}

\SourceParagraph{ch03-3.2.3.6-001}{148}
正如\BookSectionRef{sec:1.3}{1.3}中讨论的那样，CFT 中通向镜像对称的核心新特征，是 T 对偶。现在稍详细地讨论它：一方面，它在镜像对称中起着基础性作用；另一方面，它也是一种典型的“对偶”论证，而这类论证，已在更广泛的量子场论与超弦理论研究中变得十分重要。

\SourceParagraph{ch03-3.2.3.6-002}{148}
先讨论 $M\cong S^1$。所要说明的结论是
\begin{equation}\label{eq:3.58}
\operatorname{CFT}(S^1,R)\cong\operatorname{CFT}(S^1,1/R).
\end{equation}
换句话说，存在一个从 $\operatorname{CFT}(S^1,R)$ 的 Hilbert 空间到 $\operatorname{CFT}(S^1,1/R)$ 的 Hilbert 空间的幺正变换，将第一个理论的态射变为第二个理论的态射。

\SourceParagraph{ch03-3.2.3.6-003}{148}
特别地，这要求相应归一化下的配分函数相等。按照\BookSectionRef{sec:3.2.2.3}{3.2.2.3}中的思路，需要比较每个 $\Sigma$，即每一种亏格的 Riemann 曲面，上的振幅，并使这种对应与缝合相容。下面用泛函积分展示这一对应；在当前自由场模型中，下一小节还会直接给出流与态空间的对应。

\SourceParagraph{ch03-3.2.3.6-004}{148}
为了显式写出 $R$ 的依赖，取一个周长为 $2\pi$、度规固定的 $M\cong S^1$，将半径放进作用量的系数。这里的 $\varphi$ 是原先周期为 $2\pi R$ 的场除以 $R$ 所得的无量纲场，周期为 $2\pi$。作用量写成
\begin{equation}\label{eq:3.59}
S=\frac{R^2}{4\pi}\int_\Sigma
\mathrm{d}\varphi\wedge*\mathrm{d}\varphi.
\end{equation}

\begin{proof}[泛函积分推导]
\SourceParagraph{ch03-3.2.3.6-005}{149}
虽然主要目的是理解 $S^1$ 上的 T 对偶，但论证的核心，可以先通过考虑取值于 $\mathbb{R}$ 的场 $\varphi$ 来理解。随后，再回到 $S^1$ 的情形。

\SourceParagraph{ch03-3.2.3.6-006}{149}
首先，作用量为\BookEqRef{eq:3.59}{3.59}的泛函积分，可以改写成另一个泛函积分，同时对 $\varphi$ 与第二个场 $\Pi\in\Omega^1(\Sigma)$ 积分，作用量为
\begin{equation}\label{eq:3.60}
S=\frac{\pi}{R^2}\int_\Sigma\Pi\wedge*\Pi
+\mathrm{i}\int_\Sigma\Pi\wedge\mathrm{d}\varphi.
\end{equation}
由于新场 $\Pi$ 只以二次形式出现，可以按照\BookSectionRef{sec:3.2.3.1}{3.2.3.1}中\BookEqRef{eq:3.40}{3.40}的方法，精确地作出它的泛函积分。结果由鞍点 $\delta S/\delta\Pi=0$ 处的值给出\footnote{由于作用量不含 $\Pi$ 的导数，鞍点附近通常的单圈积分，在这里给出一个与场无关的 Gaussian 归一化因子。这个因子可以依赖于 $R$；比较不同亏格的振幅时，必须连同零模与测度的归一化一起保留。}。在 Euclidean 二维曲面的1形式上，$*^2=-1$，因此鞍点为
\begin{equation}\label{eq:3.61}
\frac{2\pi}{R^2}\Pi=\mathrm{i}*\mathrm{d}\varphi.
\end{equation}
将这一复鞍点代回，就重新得到\BookEqRef{eq:3.59}{3.59}。

\SourceParagraph{ch03-3.2.3.6-007}{149}
另一方面，也可以先对 $\varphi$ 积分，得到一个只对 $\Pi$ 积分的新泛函积分。先对 $(\Pi,\mathrm{d}\varphi)$ 项作分部积分。在闭曲面上，得到
\[
S=\frac{\pi}{R^2}\int_\Sigma\Pi\wedge*\Pi
+\mathrm{i}\int_\Sigma\varphi\,\mathrm{d}\Pi.
\]
随后，因为 $\varphi$ 只线性地出现在作用量中，可以作出它的泛函积分。形式上，这类似于有限维积分
\[
\int_{\mathbb{R}}\frac{\mathrm{d}\lambda}{2\pi}
\,\mathrm{e}^{-\mathrm{i}\lambda x}=\delta(x),
\]
其中，$\delta(x)$ 是 Dirac delta 函数。因此，对 $\varphi$ 作泛函积分，所得测度支撑在满足以下约束的场上：
\begin{equation}\label{eq:3.62}
\mathrm{d}\Pi=0.
\end{equation}

\SourceParagraph{ch03-3.2.3.6-008}{149}
随后，利用一个新的标量场，即“对偶玻色子” $\widehat\varphi$，局部地解出这一约束：
\begin{equation}\label{eq:3.63}
\Pi=\frac1{2\pi}\mathrm{d}\widehat\varphi.
\end{equation}
这种归一化的理由，很快就会说明。一般闭1形式还可有非零周期，因此，当前的 $\widehat\varphi$ 可以只在局部定义，其微分却在全局有定义；闭1形式不一定是单值实函数的微分。

\SourceParagraph{ch03-3.2.3.6-009}{149}
代回以后，得到
\begin{equation}\label{eq:3.64}
S=\frac1{4\pi R^2}\int_\Sigma
\mathrm{d}\widehat\varphi\wedge*\mathrm{d}\widehat\varphi.
\end{equation}
它与\BookEqRef{eq:3.59}{3.59}具有相同的形式，只需作替换 $R\to1/R$，这支持了\BookEqRef{eq:3.58}{3.58}的结论。

\SourceParagraph{ch03-3.2.3.6-010}{150}
要完成论证，必须推广到映射 $\varphi\colon\Sigma\to S^1$。利用\BookEqRef{eq:3.54}{3.54}，并采用这里周期为 $2\pi$ 的无量纲场，可以将它们写成
\[
\varphi(z)=\varphi_0+\widetilde\varphi(z)
+2\pi\sum_i m_i f_i(z)\pmod{2\pi},
\qquad m\in\mathbb{Z}^{b_1(\Sigma)}.
\]
这样，\BookEqRef{eq:3.60}{3.60}成为
\begin{equation}\label{eq:3.65}
\begin{aligned}
S={}&\frac{\pi}{R^2}\int_\Sigma\Pi\wedge*\Pi
+\mathrm{i}\int_\Sigma\Pi\wedge\mathrm{d}\widetilde\varphi\\
&+2\pi\mathrm{i}\sum_i m_i
\int_\Sigma\Pi\wedge\mathrm{d}f_i.
\end{aligned}
\end{equation}
由于 $\mathrm{d}\widetilde\varphi$ 是单值的，上面对 $\Pi$ 作泛函积分、重新得到\BookEqRef{eq:3.59}{3.59}的论证，在每个缠绕扇区中都仍然适用。

\SourceParagraph{ch03-3.2.3.6-011}{150}
不过，重复导出\BookEqRef{eq:3.64}{3.64}的步骤、求解\BookEqRef{eq:3.62}{3.62}时，在尚未施加周期条件之前，\BookEqRef{eq:3.63}{3.63}中的局部实值原函数，可能具有任意的实周期。而我们需要的，是对到 $S^1$ 的映射作泛函积分。换句话说，目前允许的周期太多。将它约束成正确积分的条件，来自泛函测度\BookEqRef{eq:3.55}{3.55}中对 $m_i$ 的额外求和：
\begin{equation}\label{eq:3.66}
\sum_{m\in\mathbb{Z}^{b_1(\Sigma)}}
\exp\left(-\mathrm{i}\sum_i m_i
\int_\Sigma\mathrm{d}\widehat\varphi\wedge\mathrm{d}f_i\right).
\end{equation}
注意，缠绕分解中的 $2\pi$，正好被\BookEqRef{eq:3.63}{3.63}中的 $1/(2\pi)$ 抵消。

\SourceParagraph{ch03-3.2.3.6-012}{150}
作出这个和，所得测度支撑在满足以下条件的场上：
\[
\frac1{2\pi}\int_\Sigma
\mathrm{d}\widehat\varphi\wedge\mathrm{d}f_i\in\mathbb{Z},
\qquad\text{对每个 }i.
\]
单值实函数 $\widehat\varphi$ 自动满足这一条件。对调和1形式的整基 $\mathrm{d}f_i$ 使用 Poincar\'e 对偶，可以得到一般解
\begin{equation}\label{eq:3.67}
\widehat\varphi=\widehat\varphi_0+
\widetilde{\widehat\varphi}+2\pi\sum_j n_j f_j(z)
\pmod{2\pi},\qquad
n\in\mathbb{Z}^{b_1(\Sigma)},
\end{equation}
其中各项的定义，与\BookEqRef{eq:3.54}{3.54}相同。因此，$\widehat\varphi$ 与 $\varphi$ 一样，是到 $S^1$ 的映射。

\SourceParagraph{ch03-3.2.3.6-013}{150}
这样，通过对 Gaussian 泛函积分\BookEqRef{eq:3.60}{3.60}作简单操作，即线性换元与鞍点处求值，就同时得到了\BookEqRef{eq:3.59}{3.59}和\BookEqRef{eq:3.64}{3.64}；两个作用量只相差替换 $R\to1/R$。还须一致地处理测度的归一化。沿用\BookEqRef{eq:3.55}{3.55}的测度，对亏格为 $g$ 的闭曲面，Poisson 求和给出
\[
Z_g(R)=R^{\chi(\Sigma)}Z_g(1/R),
\qquad\chi(\Sigma)=2-2g.
\]
因此，采用与缝合相容的 Euler 因子归一化后，
\[
\mathcal{Z}_g(R)=R^{-\chi(\Sigma)/2}Z_g(R)
=\mathcal{Z}_g(1/R),
\]
两个泛函积分的振幅就相等。特别地，在环面上 $\chi=0$，无需这一额外因子。

\SourceParagraph{ch03-3.2.3.6-014}{150}
同样的论证适用于任意紧闭曲面 $\Sigma$；有边界的情形，将在\BookSectionRef{sec:3.5.4}{3.5.4}讨论。考虑 $\Sigma$ 的复结构退化，并使用缝合公理，可以把这种对应传递到缝合所涉及的振幅。在这里的自由场理论中，结合下一小节对流、荷与振子态的显式对应，就得到两个 CFT 的同构，证明了\BookEqRef{eq:3.58}{3.58}。
\end{proof}

\subsubsection{顶点算符上的 T 对偶}\label{sec:3.2.3.7}

\SourceParagraph{ch03-3.2.3.7-001}{151}
结合\BookEqRef{eq:3.61}{3.61}与\BookEqRef{eq:3.63}{3.63}，可以得到 T 对偶理论中顶点代数之间的一个简单关系：
\[
\frac1{R^2}\mathrm{d}\widehat\varphi
=\mathrm{i}*\mathrm{d}\varphi.
\]
这里使用 Euclidean 约定。令原场与对偶场的物理坐标分别为 $X=R\varphi$ 与 $\widehat X=\widehat\varphi/R$，则 $\mathrm{d}\widehat X=\mathrm{i}*\mathrm{d}X$。由于 $*\mathrm{d}z=-\mathrm{i}\,\mathrm{d}z$，而 $*\mathrm{d}\bar z=\mathrm{i}\,\mathrm{d}\bar z$，左右行分量中的一个保持不变，另一个反号。若再复合一个交换左右行的世界面反射，可以写成
\begin{equation}\label{eq:3.68}
J_L\longmapsto J_R,\qquad
J_R\longmapsto-J_L,\qquad
V_{p_L,p_R}\longmapsto V_{-p_R,p_L}.
\end{equation}
这里，对态的荷的作用，由对流的交织关系确定；左右行的交换，同时交换了全纯与反全纯结构。

\SourceParagraph{ch03-3.2.3.7-002}{151}
对于保持世界面定向的 T 对偶，撤去这一左右行交换，就得到
\begin{equation}\label{eq:3.69}
J_L\longmapsto J_L,\qquad
J_R\longmapsto-J_R,\qquad
V_{p_L,p_R}\longmapsto V_{p_L,-p_R}.
\end{equation}
因此，T 对偶作用于荷格 $Q$，对应于 $O(1,1;\mathbb{Z})$ 中的一个非平凡变换。在前面的半径参数化中，它将 $R$ 变为 $1/R$，并交换动量数 $n$ 与缠绕数 $w$。左行振子保持不变，右行振子变号；选取相容的顶点算符相位以后，这一幺正对应保持能动张量与 OPE。它因而给出\BookEqRef{eq:3.58}{3.58}所要求的同构。

\Needspace{5\baselineskip}
\subsubsection{推广到 \texorpdfstring{$T^d$}{Tᵈ}}\label{sec:3.2.3.8}

\SourceParagraph{ch03-3.2.3.8-001}{151}
这一推广相当直接。从作用量
\[
S=\frac1{2\pi}\int_\Sigma\mathrm{d}^2z\,
(g_{ij}+B_{ij})\partial\varphi^i\bar\partial\varphi^j
\]
出发，除度规 $g_{ij}$ 以外，还可以加入反对称的二阶张量 $B_{ij}$。随后，按前面的方法处理。相容理论的荷 $(p_L,p_R)$，仍然位于符号型为 $(d,d)$ 的偶自对偶格中。在固定整格而允许其实嵌入变化时，这些数据的连续参数空间，是保持二次型的群 $O(d,d;\mathbb{R})$ 对 $O(d)\times O(d)$ 取商。它局部是一个实维数为 $d^2$ 的齐性空间，与矩阵 $g+B$ 的参数数目一致。最后，如果两个理论通过 $O(d,d;\mathbb{Z})$ 中的整格基变换相联系，它们就同构，因此模空间写为
\[
O(d,d;\mathbb{Z})\backslash O(d,d;\mathbb{R})
\big/\bigl(O(d)\times O(d)\bigr).
\]
这个整群包含由 $T^d$ 上换基诱导的 $\operatorname{GL}(d,\mathbb{Z})$ 子群；其中，保定向的换基构成 $\operatorname{SL}(d,\mathbb{Z})$。它也包含对任意一组坐标作出的 T 对偶变换。

\subsection{\texorpdfstring{$U(1)$}{U(1)} CFT 的因子分解}\label{sec:3.2.4}

\SourceParagraph{ch03-3.2.4-001}{151}
正如\BookSectionRef{sec:3.2.3.2}{3.2.3.2}中讨论的那样，从代数的观点看，自由玻色子的定义性特征，是 $U(1)$ 流代数\BookEqRef{eq:3.46}{3.46}。事实上，可以证明，CFT 中出现的 $U(1)$ 流代数，都由一个自由玻色子扇区描述。由于这是 A 模型与 B 模型中边界条件的一般分类的基础，下面说明其原理。

\SourceParagraph{ch03-3.2.4-002}{151}
考虑一个中心荷为 $c$ 的幺正 CFT $X$，其中包含一个 $U(1)$ 流 $J$，即具有以下 OPE 的算符：
\[
J(z_1)J(z_2)=\frac1{(z_1-z_2)^2}+\text{非奇异项}.
\]
基本例子是自由玻色子，其中 $J=\mathrm{i}\sqrt2\,\partial\varphi$；其 OPE 正是\BookEqRef{eq:3.50}{3.50}在这一归一化下的形式。不过，$U(1)$ 流远比自由玻色子普遍。事实上，在通常的局域幺正 CFT 中，连续的 $U(1)$ 自同构作用，产生相应的守恒流，至少有一个非零的手征流分量；通常有两个，即全纯与反全纯分量。后面，$N=2$ 超共形代数将成为主要例子。

\SourceParagraph{ch03-3.2.4-003}{152}
这样的 CFT 可以“分解”为两个部分：一个是中心荷为 $c=1$ 的自由玻色子，其能动张量为
\[
T_{U(1)}=\frac12:JJ:;
\]
另一个是商 CFT $X'$，其中心荷与能动张量分别为
\[
c'=c-1,\qquad T_{X'}=T_X-T_{U(1)}.
\]
粗略地说，因子分解意味着，将 $X$ 的 Hilbert 空间写成 $X'$ 部分与自由玻色子部分的张量积，并使两部分具有独立的 OPE。特别地，$J$ 与 $X'$ 的局部算符具有非奇异的 OPE，流代数在各个重数空间上作用平凡。完整的态空间还可能将不同荷扇区配对起来，因此，精确表述是下面的分解。

\SourceParagraph{ch03-3.2.4-004}{152}
更精确地，自由玻色子 Hilbert 空间分解成最高权表示 $V_{p_L,p_R}$，见\BookSectionRef{sec:3.2.3.4}{3.2.3.4}，这一分解提升到 $X$ 的 Hilbert 空间，成为
\begin{equation}\label{eq:3.70}
\mathcal{H}=\widehat{\bigoplus}_{p_L,p_R}
\mathcal{H}^{X'}_{p_L,p_R}\widehat\otimes V_{p_L,p_R}.
\end{equation}
这里，写的是离散荷谱；连续荷谱时，相应地改用直积分。$X$ 理论中的 OPE 与关联函数，也有类似的乘积关系，以 $X'$ 与自由玻色子的相应数据表示。

\Needspace{4\baselineskip}
\begin{proof}[代数验证]
\SourceParagraph{ch03-3.2.4-005}{152}
这一构造可以追溯到文献\BookCite{ref183}，并可从代数上作如下验证。首先，已经在\BookSectionRef{sec:3.2.3.3}{3.2.3.3}中核实，$T_{U(1)}$ 定义中心荷为1的 Virasoro 代数。其次，上面的条件足以计算 $T_{X'}$ 的 OPE，并表明，它定义中心荷为 $c'=c-1$ 的 Virasoro 代数。

\SourceParagraph{ch03-3.2.4-006}{152}
再作一些代数运算，就可以看出，这两个 Virasoro 代数彼此对易：
\[
\bigl[L_m^{U(1)},L_n^{X'}\bigr]=0.
\]
Virasoro 作用约束关联函数的位置依赖，而 $U(1)$ 流的 Ward 恒等式，固定自由玻色子部分。结合上面的荷扇区分解，就可以将 $X$ 理论的关联函数，表示为这两部分的关联数据的乘积之和；这给出所述的因子分解关系。
\end{proof}

\SourceParagraph{ch03-3.2.4-007}{152}
这一分解的一个结果是，当 $U(1)$ 因子的自同构与重数空间的荷配对及 OPE 相容时，它可以提升为 $X$ 的自同构。特别地，这个因子上的 T 对偶作用，可以按相同方式提升为完整理论之间的对偶；要使它成为同一理论的自同构，还须保持这些配对数据。

\SourceParagraph{ch03-3.2.4-008}{152}
另一个结果是，后面将讨论的边界条件，也可以按照类似的方式分解。对于保持相应流代数、并采用相容因子分解的边界条件，自由玻色子的边界条件分类，将为理论 $X$ 的边界条件分类提供其中一部分。

\SourceParagraph{ch03-3.2.4-009}{152}
下一节将引入 $N=2$ 超共形代数，并说明，它包含一个 $U(1)$ 流代数，因此这些结果适用。随后会看到，这个子扇区上的 T 对偶，就是镜像对称在 CFT 中的定义。

\subsection{形变理论}\label{sec:3.2.5}

\SourceParagraph{ch03-3.2.5-intro-001}{152}
原则上，局域 QFT 与 CFT 的形变理论，已经由 QFT 函子确定，或者由同一组数据的等价表述确定，例如算符代数，或对泛函积分的完整理解。使 QFT 发生形变的最简单方法，是在泛函积分的作用量中加入一个算符。于是，可以定义一族 QFT 的配分函数 $Z[g]$，其参数 $g^i$ 通常称为“耦合常数”，或者简称“耦合”：
\begin{equation}\label{eq:3.71}
Z[g]=\int[\mathrm{D}\varphi]\,
\mathrm{e}^{-S_0[\varphi]+\delta S},
\end{equation}
其中
\begin{equation}\label{eq:3.72}
\delta S=\sum_i g^i\int_\Sigma\mathrm{d}^dx\,O_i(x),
\end{equation}
而 $O_i(x)$ 是局部算符。这里按生成函数的符号约定，将 $\delta S$ 写在指数的正号项中。

\SourceParagraph{ch03-3.2.5-intro-002}{153}
理解\BookEqRef{eq:3.71}{3.71}的最简单方式，是把它看作未形变理论中关联函数的生成函数。换句话说，它在零耦合处的各阶导数，为有限多个积分算符 $\int O_i$ 的关联函数提供了简写。当然，这在以下意义上是“形式的”：并不要求这个 Taylor 级数收敛。要证明形变确实存在，就需要处理这一问题\footnote{虽然物理上的讨论较复杂，但在某些情形下，这一级数具有有限的正收敛半径是明确的。例如，\BookSectionRef{sec:3.2.3.5}{3.2.3.5}中改变紧化玻色子的半径 $R$ 的形变。}。

\SourceParagraph{ch03-3.2.5-intro-003}{153}
\BookEqRef{eq:3.71}{3.71}还在另一个同样重要的意义上是“形式的”：要给出精确定义，需要更多信息。正如\BookEqRef{eq:3.49}{3.49}中所见，重合点处的算符乘积通常会发散。要定义积分算符的关联函数，必须减去这些发散，也就是定义重整化的关联函数。在物理文献中，像\BookEqRef{eq:3.71}{3.71}这样的表达式，通常表示重整化关联函数的生成函数。

\Needspace{5\baselineskip}
\subsubsection{重整化理论}\label{sec:3.2.5.1}

\SourceParagraph{ch03-3.2.5.1-001}{153}
这里只介绍这一理论中几项最重要的结果。\BookSectionRef{sec:3.2.6}{3.2.6}将给出一个计算例子。

\SourceParagraph{ch03-3.2.5.1-002}{153}
首先，一个涉及度规共形缩放行为的简单论证表明，能够将一个 $d=2$ CFT 形变成另一个 CFT 的局部算符，必须具有\BookEqRef{eq:3.33}{3.33}中定义的标度维数 $h=1$\footnote{更一般地，总标度维数必须等于时空维数 $d$，这样算符才按密度变换。回忆，在 $d=2$ 时，总标度维数是 $L_0+\bar L_0$ 的本征值，即 $h_L+h_R$。还要求 $\operatorname{SO}(2)$ 转动的生成元 $L_0-\bar L_0$ 作用平凡，因此 $h_L=h_R\equiv h$。}。因此，根据\BookRef{axiom:3.4}{公理4}，CFT 空间的切空间是有限维的。

\SourceParagraph{ch03-3.2.5.1-003}{153}
如果试图作出一个 $h\ne1$ 的形变，又会怎样？要恰当地讨论它，需要引入重整化群，简称 RG。它控制着以下两者之间的联系：一边是支撑在点上的原始局部算符形变，因而决定极短距离尺度上的行为；另一边是一般距离尺度上的行为。有关 RG 的简要介绍，见 \BookBibText{ref242}{MS1} 第14章，特别是文献\BookCite{ref109}中的 Witten 讲义。

\SourceParagraph{ch03-3.2.5.1-004}{154}
这些论证最简单的结果，是将 $h<1$、$h=1$ 与 $h>1$ 分成三种情形。
\begin{definition*}\label{def:3.relevant-irrelevant-deformations}
$h>1$ 的形变称为“无关形变”。在固定点附近，它们的无量纲耦合在长距离下减小，通常不能单独给出以该固定点为紫外起点的可重整化连续形变；它们仍可作为有截断的有效场论中的相互作用。另一方面，$h<1$ 的算符称为“相关算符”。这样的形变破坏共形不变性，并在较长距离上变得更重要。
\end{definition*}
下面还会进一步解释，可以将这种对标度不变性的破坏，重新表述为 QFT 空间中的一条流，称为这一算符诱导的 RG 流。

\SourceParagraph{ch03-3.2.5.1-005}{154}
在二维局域幺正理论中，可以构造一个沿 RG 流向红外单调下降的 $c$ 函数；在共形固定点处，它等于相应 Virasoro 中心荷。因此，如果一条非平凡 RG 流连接两个共形固定点，传统上将其起点与终点分别称为紫外（UV）和红外（IR）CFT，就有
\begin{equation}\label{eq:3.73}
c_{\mathrm{UV}}>c_{\mathrm{IR}}.
\end{equation}
一般的相关形变，往往产生有质量隙的理论，其长距离极限中 $c_{\mathrm{IR}}=0$，传播自由度消失；但通过调节参数，即选择 UV 理论与形变，可以得到非平凡的 IR CFT。可以把这一过程，与将 Riemann 流形的度规缩放为 $g\to\lambda g$ 作一个粗略但有用的类比。在 $\lambda\to0$ 的极限中，紧流形趋于一个点；这可以用 Gromov--Hausdorff 拓扑精确表述，见\BookSectionRef{sec:7.3.6}{7.3.6}。不过，也能找到“塌缩”为非平凡低维流形的例子，例如带平坦度规的 $S^1\times\mathbb{R}$，在带基点的相应极限中塌缩为 $\mathbb{R}$。

\SourceParagraph{ch03-3.2.5.1-006}{154}
\begin{definition*}\label{def:3.marginal-deformations}
$h=1$ 的算符称为“边缘算符”。它生成的形变，至少在无穷小的一阶上，给出中心荷 $c$ 相同的新 CFT。
\end{definition*}
不过，前面所述的标度论证，只在耦合的一阶上成立；对算符乘积作重整化，会产生耦合的非线性修正。如果这些修正非零，而且不能通过重新定义吸收，那么标度不变性，以及共形不变性，就被破坏。此时，形变再次产生非平凡的 RG 流，$c$ 函数沿流向红外下降。如果耦合在长距离下增长，称为“边缘相关”；若耦合在长距离下减小，则称为“边缘无关”。能够保持共形不变性的边缘形变，称为“严格边缘形变”。

\SourceParagraph{ch03-3.2.5.1-007}{154}
RG 流的显式描述，由 beta 函数给出；它表示耦合随着所选重整化尺度而变化。将这个能量尺度记为 $\Lambda$，并将耦合取为无量纲变量，有
\begin{equation}\label{eq:3.74}
\Lambda\frac{\partial g^i}{\partial\Lambda}=\beta^i.
\end{equation}
因此，$\beta^i$ 是 QFT 空间上的向量场；这一空间，以\BookEqRef{eq:3.72}{3.72}中那样显式选定的耦合为参数。标度不变的 QFT，在这里的共形理论情形，是一个 RG 固定点，因而是 beta 函数的零点。在固定点附近，对于单个无量纲耦合 $g$，其线性项为
\[
\beta(g)=(2h-2)g+O(g^2).
\]
所以，对于小的正耦合，$h>1$ 时，线性 beta 函数为正，耦合向红外减小；$h<1$ 时，它为负，耦合向红外增长。当这样的流连接两个共形固定点时，终点满足\BookEqRef{eq:3.73}{3.73}。一般情况下，应根据流的线性化与非线性行为作判定，而不能仅凭某个 beta 分量的符号。

\SourceParagraph{ch03-3.2.5.1-008}{154}
在许多情形下，CFT 的形变理论有相应的几何表述。例如，有一个一般论证\BookCite{ref192}，表明 $(2,2)$ SCFT 中保持超共形结构的闭弦模参数形变没有阻碍。下面介绍的 sigma 模型构造，将说明，它包含一个数学结论作为极限情形：紧 Calabi--Yau 流形的复结构模空间没有阻碍。开弦形变的相应陈述则不成立，这对应于 Calabi--Yau 流形上的向量丛形变可能有阻碍的数学事实。\BookSectionRef{sec:3.6.2.3}{3.6.2.3}与\BookSectionRef{sec:3.6.2.4}{3.6.2.4}还会进一步讨论这一联系。

\subsection{非线性 sigma 模型}\label{sec:3.2.6}

\SourceParagraph{ch03-3.2.6-intro-001}{155}
非线性 sigma 模型，是 QFT 与 CFT 的一个非常普遍的来源。它是\BookSectionRef{sec:3.1}{3.1}中基本例子的 $d\ge2$ 维推广：那个例子描述一个量子力学粒子在靶空间 $M$ 中运动。

\SourceParagraph{ch03-3.2.6-intro-002}{155}
在 $d=2$ 时，可以把映射 $\varphi\colon\Sigma\to M$，看作描绘 $M$ 中一条弦的时空历程，因此，常把 $\Sigma$ 称为“弦的世界面”。考虑到后面共形不变性的作用，通常将 $\Sigma$ 看作复一维的曲面，而不仅仅是实二维的曲面。

\SourceParagraph{ch03-3.2.6-intro-003}{155}
现在，将靶流形记为 $X$。非线性 sigma 模型，是一个对所有映射 $\varphi\colon\Sigma\to X$ 作泛函积分来定义的场论，其作用量为
\begin{equation}\label{eq:3.75}
S[\varphi]=\frac1{2\pi}\int_\Sigma\mathrm{d}^2z\,
\widehat G_{ij}(\varphi)
\frac{\partial\varphi^i}{\partial z}
\frac{\partial\varphi^j}{\partial\bar z}.
\end{equation}
这里，$z$ 是 $\Sigma$ 上的复坐标，$\varphi^i$ 是 $X$ 的局部实坐标，而 $\widehat G$ 是 $T^*X\otimes T^*X$ 的一个张量，不必对称。也可以将它写成
\[
\widehat G=g+B,
\]
其中，$g_{ij}$ 是度规，$B_{ij}$ 是反对称张量，或者说2形式。这里沿用前面 $\alpha'=1$ 以及 $\mathrm{d}^2z=2\,\mathrm{d}x\,\mathrm{d}y$ 的归一化。

\SourceParagraph{ch03-3.2.6-intro-004}{155}
\begin{definition*}\label{def:3.sigma-model-equations}
按照定义，Euler--Lagrange 方程，也称“运动方程”，表示作用量在场的无穷小变化下，保持一阶不变。在这里，这些方程是
\begin{equation}\label{eq:3.76}
\begin{aligned}
0=\frac{\delta S}{\delta\varphi^k}\quad\Longleftrightarrow\quad
0={}&\partial\!\left(\widehat G_{kj}(\varphi)
\bar\partial\varphi^j\right)
+\bar\partial\!\left(\widehat G_{ik}(\varphi)
\partial\varphi^i\right)\\
&-\widehat G_{ij,k}(\varphi)
\partial\varphi^i\bar\partial\varphi^j.
\end{aligned}
\end{equation}
这些方程的解，称为“理论的经典解”。
\end{definition*}
显然，最简单的解，是 $\varphi$ 在 $\Sigma$ 上取常值。另一个简单的解族，满足假设 $\partial\varphi/\partial\operatorname{Im}z=0$。此时，$B$ 项消失，方程化为 $d=1$ 的方程，其解是度规 $g_{ij}$ 的测地线。

\SourceParagraph{ch03-3.2.6-intro-005}{155}
更一般地，在 $\mathrm{d}B=0$ 的条件下，这些方程就是定义调和映射的方程：
\[
\tau_\gamma(\varphi)
=\operatorname{tr}_\gamma\nabla\mathrm{d}\varphi=0,
\]
其中，$\gamma$ 是世界面的度规，$\nabla$ 使用 $\gamma$ 与靶度规 $g$ 的 Levi--Civita 联络。因此，另一个解，是取 $\varphi(\Sigma)$ 为一个光滑的体积极小的二维循环，并对它采用与诱导度规相容的共形参数化。

\subsubsection{微扰理论}\label{sec:3.2.6.1}

\SourceParagraph{ch03-3.2.6.1-001}{156}
这些经典解，在量子场论中也很重要，因为赋予泛函积分\BookEqRef{eq:3.14}{3.14}意义的最简单方法，就是半经典方法。假定 $\hbar$ 很小，将 $Z$ 写成对积分的“鞍点” $\varphi_i$ 的形式求和：
\begin{equation}\label{eq:3.77}
Z\sim\sum_{\delta S[\varphi_i]=0}
\frac{\mathrm{e}^{-S[\varphi_i]/\hbar}}
{\bigl(\det{}'\Delta_i\bigr)^{1/2}}.
\end{equation}
这里仅显示非零模的 Gaussian 因子，略去了测度的整体归一化；若鞍点有零模，还须对相应的集体坐标积分。

\SourceParagraph{ch03-3.2.6.1-002}{156}
前面的因子，通过将 $S[\varphi_i+\delta\varphi_i]$ 对小扰动 $\delta\varphi$ 展开到二阶，再像\BookSectionRef{sec:3.2.3}{3.2.3}中那样作 Gaussian 积分而得到。采用适当的坐标系，有助于完成这一计算。给定一般的 Riemann 度规 $g_{ij}$ 与一个指定点 $\varphi$，自然的选择，是以 $\varphi$ 为中心的 Riemann 法坐标。在这种坐标中，
\[
g_{ij}(\varphi+\delta\varphi)
=g_{ij}(\varphi)-\frac13R_{ikjl}(\varphi)
\delta\varphi^k\delta\varphi^l+\cdots,
\]
其中，$R_{ikjl}$ 是 Riemann 曲率张量，采用正曲率球面具有正截面曲率的约定。

\SourceParagraph{ch03-3.2.6.1-003}{156}
使用这一展开，并为简单起见取 $B=0$，可以将作用量\BookEqRef{eq:3.75}{3.75}中的二次项改写如下。对于一般背景映射，用指数映射表示扰动 $\xi\in\Gamma(\varphi^*TX)$，并令 $D$ 为拉回的 Levi--Civita 协变导数。将 $\delta^2S$ 记为 Taylor 展开的二次项，则
\begin{equation}\label{eq:3.78}
\delta^2S[\varphi]=\frac1{2\pi}\int_\Sigma\mathrm{d}^2z\,
\left(g_{ij}(\varphi)D_z\xi^iD_{\bar z}\xi^j
-R_{ikjl}(\varphi)\partial\varphi^i\bar\partial\varphi^j
\xi^k\xi^l\right).
\end{equation}
由此，定义了\BookEqRef{eq:3.77}{3.77}中相应鞍点处的二次算符 $\Delta_i$。

\SourceParagraph{ch03-3.2.6.1-004}{156}
同样，$\det{}'\Delta$ 的发散，必须通过重整化处理。在早期的讨论中，这被定义为一个形式过程：将发散的“反项”加入作用量\BookEqRef{eq:3.78}{3.78}，以抵消泛函积分产生的发散。从这一观点看，重整化理论的主要目标，是证明这些发散，可以写成有限种场的局域泛函之和，从而由具有有限种项的作用量抵消。对于一般非线性 sigma 模型，这些项的系数本身可以是靶坐标的函数。
\begin{definition*}\label{def:3.renormalizable-action}
如果所有必需的项，都已包含在作用量中，就称这个作用量可重整化。
\end{definition*}

\SourceParagraph{ch03-3.2.6.1-005}{156}
在较现代的 RG 语言中，表述有所不同。定义一族依赖额外参数 $\Lambda$ 的“截断”或正规化理论；通常取 $\Lambda$ 的量纲为长度的倒数，并从泛函测度中去掉所有长度尺度小于 $1/\Lambda$ 的涨落。随后，计算带截断的泛函积分对 $\Lambda$ 的变化。最后，令作用量 $S$ 产生一个补偿前述变化的变化，以落实最终结果应与 $\Lambda$ 无关的原则。

\SourceParagraph{ch03-3.2.6.1-006}{156}
这些计算的结果，可以归结为 beta 函数\BookEqRef{eq:3.74}{3.74}。粗略地说，它由配分函数的尺度依赖确定。沿用\BookEqRef{eq:3.71}{3.71}的生成函数约定，将显式的截断依赖与耦合的运行分开写，有
\[
\left(\Lambda\frac{\partial}{\partial\Lambda}
+\sum_i\beta^i\frac{\partial}{\partial g^i}\right)\log Z=0,
\]
因而
\[
\left.\Lambda\frac{\partial\log Z}{\partial\Lambda}\right|_g
=-\sum_i\beta^i
\left\langle\int_\Sigma\mathrm{d}^dx\,O_i(x)\right\rangle_g.
\]
这里的 $O_i$，是\BookEqRef{eq:3.72}{3.72}中定义的局部算符。精确定义，特别是将求和限制为抵消发散所需的有限种算符，还是允许更大的求和，取决于采用的形式体系。

\SourceParagraph{ch03-3.2.6.1-007}{157}
RG 的一般行为，强烈依赖于时空维数 $d$。通常，一个给定算符的标度维数，包括经典的维数，也称为“工程维数”，以及量子修正；在微扰理论中，后者由 $\hbar$ 的形式幂级数给出。临界维数，是使经典标度维数等于 $d$ 的时空维数选择；此时，RG 由量子修正控制。

\SourceParagraph{ch03-3.2.6.1-008}{157}
当一般作用量\BookEqRef{eq:3.75}{3.75}处于其临界维数时，非线性 sigma 模型最有意思。为此，所需的标度维数 $d$，应完全由导数的标度贡献，这样就可以令场 $\varphi$ 的经典标度维数为零。由于\BookEqRef{eq:3.75}{3.75}含有两个导数，这迫使 $d=2$。对于 $d>2$，在 Gaussian 固定点附近作幂次计数，可以看出，$\widehat G$ 的非线性相互作用是无关的；在这一微扰长距离极限中，理论局部趋向带 Euclidean 度规的自由场。

\SourceParagraph{ch03-3.2.6.1-009}{157}
在 $d=2$ 时，可以把度规张量 $\widehat G_{ij}$，看作无穷多个耦合常数的形式生成函数；或者更恰当地，把 QFT 空间看作由一个度规空间参数化。因此，beta 函数，是度规空间上的向量场。可以用\BookSectionRef{sec:3.2.3.1}{3.2.3.1}的形式体系计算它。对于 $\mathrm{d}B=0$ 的特殊情形，在主导的单圈阶，得到
\begin{equation}\label{eq:3.79}
\Lambda\frac{\partial g_{ij}}{\partial\Lambda}
=\beta_{ij}=\hbar R_{ij}[g]+O(\hbar^2).
\end{equation}
这里，$\Lambda$ 是能量尺度。若改用流向红外的参数 $t=-\log\Lambda$，便得到 $\partial_tg_{ij}=-\hbar R_{ij}+O(\hbar^2)$，即 Ricci 流。单圈项在有效作用量中属于 $\hbar^0$ 阶，而度规的 beta 函数从 $\hbar$ 阶开始。

\SourceParagraph{ch03-3.2.6.1-010}{157}
可以继续将重整化作用量、配分函数与关联函数，展开到 $\hbar$ 的更高阶，得到标准的弱耦合微扰展开。在 sigma 模型中，这也相当于导数展开。例如，\BookEqRef{eq:3.79}{3.79}中的下一项，是 Riemann 曲率张量的二次式，与 Ricci 张量相比，多出两个导数。因此，存在一个使修正相对于主导项趋于零的极限，称为“大体积极限”“大结构极限”或“超引力极限”。

\SourceParagraph{ch03-3.2.6.1-011}{157}
RG 流\BookEqRef{eq:3.79}{3.79}由 Friedan 发现\BookCite{ref160}。他进一步证明，重整化过程对各阶 $\hbar$ 都关于 $X$ 上的坐标变换协变。因此，该 QFT 只依赖于度规与 B 场的微分同胚等价类。从这个意义上说，可以把 sigma 模型，看作从 Riemann 几何到 QFT 的一个函子。

\SourceParagraph{ch03-3.2.6.1-012}{157}
超对称 sigma 模型，见\BookSectionRef{sec:3.3.2}{3.3.2}，具有类似的重整化理论；在主导阶，同样得到\BookEqRef{eq:3.79}{3.79}。在后面将讨论的 $(2,2)$ 模型中，单圈项之后的首个非零修正，出现在四圈阶，具有 Riemann 曲率四阶量相应的导数阶数\BookCite{ref198}。

\SourceParagraph{ch03-3.2.6.1-013}{157}
遗憾的是，无论玻色还是超对称 sigma 模型，微扰展开通常都是 $\hbar$ 的渐近级数，系数呈阶乘增长；因此，无论从数学上还是物理上，都不能将它本身，看作 QFT 的令人满意的定义。这一问题，与重整化问题完全不同，既长期存在，又十分根本，以致在这类一般性讨论中，往往被搁置不谈。

\SourceParagraph{ch03-3.2.6.1-014}{158}
对于当前情形，以及四维 Yang--Mills 理论等类似情形，已有许多物理论证支持这样一种看法：这在本质上是技术问题，确实应存在 QFT 与 CFT 的配分函数，而微扰理论为它们提供良好的渐近展开。此外，人们也相当清楚，物理论证中的哪些环节，本质上依赖微扰理论，哪些环节则可以不依赖微扰展开而证明，或者至少预期成立。当称一个物理论证是“启发式的”、却仍认为它有说服力时，隐含依靠的，正是这种理解。

\SourceParagraph{ch03-3.2.6.1-015}{158}
在更简单的情形中，例如\BookSectionRef{sec:3.6.8}{3.6.8}将讨论的 Landau--Ginzburg 理论，甚至存在通过收敛级数展开，对 QFT 作出的显式“构造”，尽管这些构造极其复杂\BookCite{ref182}。也许有一天，类似或更简单的论证，会为本书讨论的全部物理内容，提供坚实的数学基础。

\subsubsection{CFT 与重整化群}\label{sec:3.2.6.2}

\SourceParagraph{ch03-3.2.6.2-001}{158}
要使 sigma 模型成为 CFT，beta 函数必须为零。由\BookEqRef{eq:3.79}{3.79}可知，在主导阶，并且 $\mathrm{d}B=0$ 时，度规必须 Ricci 平坦。因此，在大体积极限下，任何 Ricci 平坦的靶空间 $X$，都可以用来定义 CFT。

\SourceParagraph{ch03-3.2.6.2-002}{158}
在有限但很大的体积下，必须研究\BookEqRef{eq:3.79}{3.79}的修正，并判断，是否有阻碍使解无法延伸到更高阶。对于这类一般模型，beta 函数尚没有精确的闭式结果；不过，对于下面将讨论的 $N=2$ 超对称 sigma 模型，使用超空间形式体系，不难看出，在复坐标中，beta 函数具有以下形式：
\[
\frac{\beta_{i\bar j}}{\hbar}
=R_{i\bar j}+\sum_{k\ge1}\hbar^k
\partial_i\partial_{\bar j}F_k,
\]
其中，$F_k$ 是 $X$ 上的全局定义的函数。文献\BookCite{ref374}的论证表明，对于所讨论的紧 Calabi--Yau 靶空间，可以逐阶修正度规，使所得 $\beta=0$ 条件，作为 $\hbar$ 的形式级数，在所有阶成立。

\SourceParagraph{ch03-3.2.6.2-003}{158}
另一方面，对于曲率非零、$\mathrm{d}B=0$ 且伸缩子为常值的非平凡玻色 sigma 模型，没有上述已知的共形构造。确实存在 $\mathrm{d}B\ne0$ 的非平凡模型，即 Wess--Zumino--Witten 模型。不过，由于主要关注超对称 sigma 模型，这里不再继续讨论这一话题。

\SourceParagraph{ch03-3.2.6.2-004}{158}
$\beta\ne0$ 的情形，最适合用 RG 的语言理解。自然地，可以将\BookEqRef{eq:3.74}{3.74}，看作在耦合常数空间上定义一个向量场或流，而这个空间，被看作 QFT 的空间。于是，$\beta=0$ 的零点，就是流的固定点；不过，流本身也有意义，因为带截断的 QFT 的基本定义，原则上给出了所有可观测量，即关联函数，作为 $\Lambda$ 的函数的定义。

\SourceParagraph{ch03-3.2.6.2-005}{159}
考虑一条流，它在 $\Lambda\to\infty$ 时渐近于固定点 $g_{\mathrm{UV}}$，在 $\Lambda\to0$ 时渐近于固定点 $g_{\mathrm{IR}}$。在物理上，可以通过取关联函数的 $\Lambda\to0$ 极限，达到后一个固定点；等价地，就是考虑极长距离尺度的极限，通常称为“红外”或 IR 极限。同样，$\Lambda\to\infty$ 的极限，定义“紫外”或 UV 极限。

\SourceParagraph{ch03-3.2.6.2-006}{159}
从物理上说，给定的泛函积分，是否存在满足 QFT 公理的真正定义，被认为等价于以下问题：能否定义一条 RG 流，使它从一个定义良好的 UV 固定点渐近出发。这带来红外自由与紫外自由理论之间的重要区别。对于一个小的正边缘相互作用耦合 $u$，若 $\beta(u)>0$，它向红外减小，理论是 IR 自由的；若 $\beta(u)<0$，它向紫外减小，理论是 UV 自由的。后一类提供独立 QFT 的一个通常的 UV 完备构造。前一类若向紫外不断增强、又没有其他 UV 固定点，则需要作为具有更多自由度的 CFT 或 UV 自由 QFT 的子理论来实现，例如嵌入靶空间 $X$ 维数更高的 sigma 模型。按照\BookSectionRef{sec:3.2.5}{3.2.5}的术语，前者是边缘相关形变，后者是边缘无关形变。

\SourceParagraph{ch03-3.2.6.2-007}{159}
\BookEqRef{eq:3.79}{3.79}的符号意味着，正 Ricci 曲率的流形，例如球面，提供 UV 自由的 sigma 模型；正第一 Chern 类 $c_1(X)>0$ 的复流形，也给出这类典型模型。这里的相互作用耦合，与度规的整体尺度成反比，因此，它的 beta 函数具有相反的符号。这些模型通常也称为“有质量模型”，因为许多这样的渐近自由模型，会形成“质量隙”，具有与下面讨论的模型相当不同的物理解释。还应说明，许多这类模型可以作拓扑扭曲，得到一种理论；虽然它可能缺少后面至关重要的一个要素，即边界 $U(1)$ 荷，但在其他方面，与将要讨论的理论类似。把它们纳入当前框架的一个简单方法，是改为考虑相关的 sigma 模型，以典范丛 $K_X$ 或余切丛 $T^*X$ 的全空间为靶空间。有关这些模型的直接讨论，见 \BookBibText{ref242}{MS1}。

\subsubsection{线性 sigma 模型}\label{sec:3.2.6.3}

\SourceParagraph{ch03-3.2.6.3-001}{159}
在 RG 方法中，可以用比 $\beta=0$ 更弱的条件来得到一个 CFT：不必坚持当前 QFT 已经处于固定点，而可以考虑一条渐近于固定点的流。虽然流途中的理论，并不就是 CFT，但取 IR 极限后，可以得到一个 CFT，即 IR 固定点。

\SourceParagraph{ch03-3.2.6.3-002}{159}
物理文献中大量使用的一个主要例子，是线性 sigma 模型。它是一条 UV 极限为自由玻色子理论的流，通过在作用量中加入一个势 $V(\varphi)$ 得到。在前面的归一化下，可以写为
\[
S_{\mathrm{LSM}}=\frac1{2\pi}\int_\Sigma\mathrm{d}^2z\,
\left(|\partial\varphi|^2+V(\varphi)\right).
\]
从经典上看，预期由这一作用量描述的粒子，会试图使能量最小，也就是停留在函数 $V(\varphi)$ 的极小值处。

\SourceParagraph{ch03-3.2.6.3-003}{160}
\begin{example*}
取
\[
V(\varphi)=f(\varphi)^2,
\]
则在零集为正则超曲面的情形下，这一极小值条件，在 $\mathbb{R}^n$ 中定义一个非平凡的靶空间 $X$，即实超曲面 $f=0$，也称为“约束曲面”。
\end{example*}

\SourceParagraph{ch03-3.2.6.3-004}{160}
当然，在 QFT 中讨论这一问题，还需要讨论重整化。出发点是，从原来自由玻色子理论的经典幂次计数看，不含导数的势项是相关形变；具体算符的量子标度维数，还须在 RG 中考虑。对于这里产生约束曲面的势，RG 会趋于消去参数化约束曲面法向涨落的场，得到一个以 $X$ 为靶空间的非线性 sigma 模型。在超对称情形中，这一描述尤其有价值，因为可以论证，多种可观测量，特别是经过拓扑扭曲的 B 模型中的可观测量，沿流保持不变。有关这些模型更详细的讨论，见 \BookBibText{ref242}{MS1}。

\section{超对称场论与拓扑场论}\label{sec:3.3}

\SourceParagraph{ch03-3.3-intro-001}{160}
正如上一节所见，虽然可以把 CFT 的框架，看作谱几何这一主题的一种自然变体，但从数学上看，它仍处于发展的早期阶段。它与数学的大部分联系，是通过与之相关的二维拓扑场论框架建立的。\BookChapterRef{2}中已经看到，这一框架可以精确定义。很快还会说明，它与具有扩展超对称的 CFT 直接相关。

\SourceParagraph{ch03-3.3-intro-002}{160}
随后，将详细讨论主要的例子，即超对称非线性 sigma 模型，以及使它们成为拓扑模型的 A 扭曲与 B 扭曲。

\subsection{扭曲与拓扑场论}\label{sec:3.3.1}

\SourceParagraph{ch03-3.3.1-001}{160}
与讨论量子力学时一样，现在希望遵循上同调型 TFT 的一般方法\BookCite{ref467,ref468,ref109}，在 QFT 中分离出一个较简单的“拓扑”扇区。

\SourceParagraph{ch03-3.3.1-002}{160}
于是，从一个具有度规依赖的通常 QFT 出发，寻找一个算符 $Q$，使它满足以下条件：
\begin{definition*}\label{def:3.cohomological-tft-sector}
\begin{itemize}
\item $Q^2=0$，因而可以定义其上同调。
\item 能动张量 $T$ 在 $Q$ 上同调中平凡，也就是说，对某个局部算符 $b$，有
\begin{equation}\label{eq:3.80}
T=\{Q,b\}.
\end{equation}
\end{itemize}
\end{definition*}
根据\BookEqRef{eq:3.29}{3.29}，所有度规依赖，都由能动张量决定。因此，取 $Q$ 上同调，就消去了对度规的依赖，从而得到一个满足\BookChapterRef{2}公理的 TFT。

\SourceParagraph{ch03-3.3.1-003}{160}
在物理上，当 $Q$ 是费米算符时，方程 $Q^2=\frac12\{Q,Q\}=0$ 自然地表现为一个反对易关系。结合\BookEqRef{eq:3.80}{3.80}，这提示，可以从超对称量子场论的超荷中得到 $Q$。与 SQM 一样，超对称 QFT 的基本定义性特征，是存在超荷，即一组线性无关的 Hermitian 算符 $Q_I$。在约化到量子力学的关系中，取适当的归一化，可以写为
\begin{equation}\label{eq:3.81}
\{Q_I,Q_J\}=\delta_{IJ}H,
\end{equation}
其中，$H$ 是 Hamilton 算符；完整的二维关系还包含空间动量。可以证明，在转动作用下，超荷按照若干自旋表示的直和变换。回忆，在二维中，有两个不等价的一维自旋表示，定向反转交换它们。在宇称不变的理论中，可以使用取值于这两个表示之和的 Dirac 旋量。如果有 $p$ 个这样的旋量超荷，就称理论具有 $p$ 重扩展超对称。等价地，可以称它具有 $N=(p,p)$ 超对称，表示有 $p$ 个左手超荷与 $p$ 个右手超荷。这里主要讨论 $N=(2,2)$ QFT，特别是具有 $N=(2,2)$ 超对称的 sigma 模型\footnote{在二维中，也可以有左右手超荷数不等的幺正理论，即 $p\ne q$ 的 $N=(p,q)$ 超对称。在这些理论中，\BookEqRef{eq:3.81}{3.81}须写成完整的手征关系
\[
\{Q_I^\pm,Q_J^\pm\}
=\tfrac12\delta_{IJ}(H\pm P),
\qquad\{Q_I^\pm,Q_J^\mp\}=0,
\]
其中，$Q^+$ 与 $Q^-$ 分别为右手与左手超荷，$P$ 是空间平移的生成元，即空间动量算符。}。

\SourceParagraph{ch03-3.3.1-004}{161}
与往常一样，假定超对称是局域实现的。这意味着，存在费米局部算符 $G_I^\mu(x)$，称为超流，在平坦时空中满足 $\partial_\mu G_I^\mu=0$，从而
\[
Q_I=\int_Y\mathrm{d}\operatorname{Vol}_Y\,G_I^\mu v_\mu.
\]
这里，$v^\mu$ 与前面一样，是单位类时向量。

\SourceParagraph{ch03-3.3.1-005}{161}
在\BookEqRef{eq:3.81}{3.81}的代数中，如果 $N\ge2$，可以尝试取 $Q=Q_1+\mathrm{i}Q_2$，得到一个 TFT。超对称代数随即给出 $Q^2=0$，以及 $H$ 是 $Q$ 恰当的。在完整二维理论中，必须选择相容的手征组合；例如，$(2,2)$ 超对称允许这样的幂零组合。不过，能动张量 $T$ 通常既不是 $Q$ 恰当的，甚至也不是 $Q$ 闭的。一般而言，要得到 TFT，需要比上述两个超荷更丰富的结构。

\SourceParagraph{ch03-3.3.1-006}{161}
所需的结构，是一个守恒的玻色流 $J_\mu$，满足 $\partial_\mu J^\mu=0$，并且对某个费米张量 $b_{\mu\nu}$，有\footnote{这里，$\epsilon_{\nu\alpha}$ 是二维反对称张量，取 $\epsilon_{12}=1$。}
\begin{equation}\label{eq:3.82}
\{Q,b_{\mu\nu}\}
=T_{\mu\nu}-\frac14\epsilon_{\mu\alpha}\partial^\alpha J_\nu
-\frac14\epsilon_{\nu\alpha}\partial^\alpha J_\mu.
\end{equation}

\SourceParagraph{ch03-3.3.1-007}{161}
于是，可以定义一个新的能动张量，它是守恒的、对称的，并且是 $Q$ 恰当的：
\begin{equation}\label{eq:3.83}
T'_{\mu\nu}=T_{\mu\nu}
-\frac14\epsilon_{\mu\alpha}\partial^\alpha J_\nu
-\frac14\epsilon_{\nu\alpha}\partial^\alpha J_\mu.
\end{equation}
在空间方向没有边界贡献的情形下，$T'$ 对应的 Hamilton 算符，与 $T$ 的相同。因此，\BookEqRef{eq:3.82}{3.82}的积分形式为
\[
\{Q,G\}=H,
\qquad
G=\int_Y\mathrm{d}\operatorname{Vol}_Y\,b_{\mu\nu}v^\mu v^\nu.
\]
这提示，$b$ 鬼场 $b_{\mu\nu}$，应是超流 $G_I^\mu$ 的某个线性组合。

\SourceParagraph{ch03-3.3.1-008}{162}
另一个约束，来自要求 $Q$ 在 Lorentz 变换下是标量\footnote{如果 $Q$ 不是标量，在一般弯曲流形上考虑理论时，它就不能作为上述方式的全局守恒标量超荷。}。在超对称场论中，$Q$ 是旋量的一个分量；不过，由于能动张量已经改变，Lorentz 变换的生成元也改变：
\[
\begin{aligned}
M_{\mu\nu}
&=\int_Y\mathrm{d}\operatorname{Vol}_Y\,
(x_\mu T_{\rho\nu}-x_\nu T_{\rho\mu})v^\rho\\
&\longmapsto
M'_{\mu\nu}=\int_Y\mathrm{d}\operatorname{Vol}_Y\,
(x_\mu T'_{\rho\nu}-x_\nu T'_{\rho\mu})v^\rho.
\end{aligned}
\]
如果
\[
[R,Q]=Q,\qquad[R,G]=-G,
\]
其中
\[
R=\int_Y\mathrm{d}\operatorname{Vol}_Y\,J_\mu v^\mu,
\]
那么，$Q$ 确实与 $M'_{\mu\nu}$ 对易。

\SourceParagraph{ch03-3.3.1-009}{162}
换句话说，$J_\mu$ 对应于 $U(1)$ R 对称，而 $Q$ 的 R 荷为1。这意味着，$Q$ 上同调按 R 荷分次。在上面采用的 R 荷归一化下，这一分次常常是有理数分次，而非整数分次。

\SourceParagraph{ch03-3.3.1-010}{162}
因此，将寻找具有 $U(1)$ R 对称的 $N=2$ 超对称量子场论，作为构造拓扑场论的候选。事实上，在 $d=2$ 中，将这一要求与共形对称结合，便定义了 $N=2$ 超共形代数，简称 SCA；这施加了非常强的约束。不过，先介绍基本例子，然后再展开它的一般结果。

\subsection{超对称 sigma 模型}\label{sec:3.3.2}

\SourceParagraph{ch03-3.3.2-001}{162}
这里最重要的一类超对称 QFT，是超对称 sigma 模型。可以按照与前面相同的程序，将 sigma 模型超对称化。在泛函积分的语言中，为每个玻色场 $X^i$，即 $M$ 的独立局部坐标，引入一个“伙伴费米子” $\psi^i$。随后，设定一个在以下一般形式的超对称作用下不变的作用量：
\begin{equation}\label{eq:3.84}
\delta X^i=\epsilon\psi^i,\qquad
\delta\psi^i=\epsilon\partial X^i.
\end{equation}
这里，$\epsilon$ 是费米参数。当然，$\Sigma$ 上的偏导数 $\partial$ 有两个分量，因此，还需要作更多选择。最对称的理论，允许使用任意一个分量，于是具有每种手征各一个超对称。因此，对任意带度规的流形 $M$，引入费米子 $\psi^i$，就可以得到具有 $(1,1)$ 超对称的 sigma 模型。其一般理论，与玻色 sigma 模型大体类似。特别地，如果 beta 函数为零，就会得到一个 SCFT；在主导阶，这一条件仍然要求 $M$ Ricci 平坦。不过，更高阶的修正不同。

\SourceParagraph{ch03-3.3.2-002}{163}
虽然这些模型具有两个超对称，但没有类似于流 $J$ 的结构；在\BookSectionRef{sec:3.3.1}{3.3.1}中已经看到，要得到拓扑理论，需要这样的流。因此，需要扩展超对称。与超对称量子力学一样，要在超对称 sigma 模型中实现代数\BookEqRef{eq:3.8}{3.8}，靶空间 $X$ 必须具有相容的复结构。在下面的 K\"ahler 情形中，标准的超对称 sigma 模型，具有 $(2,2)$ 超对称。

\SourceParagraph{ch03-3.3.2-003}{163}
更一般地，如果 $X$ 具有 $k$ 个线性无关、协变常值，并满足相应超对称代数相容关系的复结构，sigma 模型就具有 $(k+1,k+1)$ 世界面超对称。本质上，唯一的非平凡情形，是超 K\"ahler 几何与 $(4,4)$ 超对称。也可以通过加入不保持宇称的耦合，包括 B 场的某些特殊情形，得到 $p\ne q$ 的 $(p,q)$ 超对称。文献\BookCite{ref285,ref474}表明，$(0,2)$ 超对称，与文献\BookCite{ref345,ref187,ref189,ref188}在复流形 $X$ 上构造的手征代数层相关。

\SourceParagraph{ch03-3.3.2-004}{163}
下面只讨论 $(2,2)$ 超对称。存在一个相容且协变常值的复结构，等价于假定 $X$ 是 K\"ahler 流形。现在改用复坐标，记为 $\phi^i$，其复共轭记为 $\phi^{\bar\imath}$。

\SourceParagraph{ch03-3.3.2-005}{163}
为了定义费米子，引入\BookChapterRef{2}中讨论的两个 $\operatorname{spin}^c$ 结构。回忆，这里通过一对全纯线丛 $L_1$ 与 $L_2$，以及同构
\[
L_1\otimes L_2\cong K_\Sigma\equiv T^{*1,0}\Sigma,
\]
来描述所需的结构。基本例子是 $L_1\cong L_2\cong K_\Sigma^{1/2}$，即 $K_\Sigma$ 的平方根；不过，很快就会需要更一般的情形。第二个结构满足 $L_3\otimes L_4\cong K_\Sigma$。

\SourceParagraph{ch03-3.3.2-006}{163}
\begin{definition*}\label{def:3.sigma-fermion-bundles}
将费米子定义为 $\Sigma$ 上下列丛的截面：
\begin{equation}\label{eq:3.85}
\begin{aligned}
\psi_+^i&\in\Gamma(L_1\otimes\phi^*T_X),\\
\psi_+^{\bar\jmath}&\in\Gamma(L_2\otimes\phi^*\overline{T_X}),\\
\psi_-^i&\in\Gamma(\overline{L_3}\otimes\phi^*T_X),\\
\psi_-^{\bar\jmath}&\in\Gamma(\overline{L_4}\otimes\phi^*\overline{T_X}).
\end{aligned}
\end{equation}
\end{definition*}
这里，$T_X$ 是 $X$ 的全纯切丛，横线表示相应的反全纯丛。$D$ 表示协变导数，
\[
D\psi_-^i=\partial\psi_-^i+
\partial\phi^j\Gamma^i_{jk}\psi_-^k,
\]
其中，$\partial$ 与往常一样，是 de Rham 微分的全纯部分\footnote{文献中还有许多书写 $N=(2,2)$ 理论的约定。例如，有时用 $\pm$ 记号表示 $U(1)$ 荷的正负。}。

\SourceParagraph{ch03-3.3.2-007}{164}
作用量为
\begin{equation}\label{eq:3.86}
\begin{aligned}
S=\frac1{4\pi}\int_\Sigma\mathrm{d}^2z\,
\Bigl\{&g_{i\bar\jmath}
\bigl(\partial\phi^i\bar\partial\phi^{\bar\jmath}
+\bar\partial\phi^i\partial\phi^{\bar\jmath}\bigr)\\
&+B_{i\bar\jmath}
\bigl(\partial\phi^i\bar\partial\phi^{\bar\jmath}
-\bar\partial\phi^i\partial\phi^{\bar\jmath}\bigr)\\
&+\mathrm{i}g_{i\bar\jmath}\psi_-^{\bar\jmath}D\psi_-^i
+\mathrm{i}g_{i\bar\jmath}\psi_+^{\bar\jmath}\bar D\psi_+^i\\
&+R_{i\bar k j\bar l}\psi_+^i\psi_+^{\bar k}
\psi_-^j\psi_-^{\bar l}\Bigr\},
\end{aligned}
\end{equation}
其中，$g_{i\bar\jmath}$ 是 K\"ahler 度规。

\SourceParagraph{ch03-3.3.2-008}{164}
张量 $B_{i\bar\jmath}$ 称为 B 场。一般地，作用量中依赖于 B 的部分，就是将2形式 $B\in\Omega^2(X)$ 拉回到世界面 $\Sigma$ 后的积分。对于这里的目的，只需假定 $\mathrm{d}B=0$。\BookSectionRef{sec:3.4.1}{3.4.1}将说明，此时 sigma 模型，只依赖于 B 在 $H^2(X,\mathbb{R})/H^2(X,\mathbb{Z})$ 中的类。此外，在大部分例子中，$H^{2,0}(X)$ 平凡，因此，B 可以取为实的 $(1,1)$ 形式；在这一情形中，写为
\[
B=\tfrac12B_{i\bar\jmath}\,
\mathrm{d}\phi^i\wedge\mathrm{d}\phi^{\bar\jmath}.
\]

\SourceParagraph{ch03-3.3.2-009}{164}
超对称由以下变换给出：
\begin{equation}\label{eq:3.87}
\begin{aligned}
\delta\phi^i&=\mathrm{i}\alpha_-\psi_+^i
+\mathrm{i}\alpha_+\psi_-^i,\\
\delta\phi^{\bar\imath}&=\mathrm{i}\widetilde\alpha_-\psi_+^{\bar\imath}
+\mathrm{i}\widetilde\alpha_+\psi_-^{\bar\imath},\\
\delta\psi_+^i&=-\widetilde\alpha_-\partial\phi^i
-\mathrm{i}\alpha_+\psi_-^j\Gamma^i_{jk}\psi_+^k,\\
\delta\psi_+^{\bar\imath}&=-\alpha_-\partial\phi^{\bar\imath}
-\mathrm{i}\widetilde\alpha_+\psi_-^{\bar\jmath}
\Gamma^{\bar\imath}_{\bar\jmath\bar k}\psi_+^{\bar k},\\
\delta\psi_-^i&=-\widetilde\alpha_+\bar\partial\phi^i
-\mathrm{i}\alpha_-\psi_+^j\Gamma^i_{jk}\psi_-^k,\\
\delta\psi_-^{\bar\imath}&=-\alpha_+\bar\partial\phi^{\bar\imath}
-\mathrm{i}\widetilde\alpha_-\psi_+^{\bar\jmath}
\Gamma^{\bar\imath}_{\bar\jmath\bar k}\psi_-^{\bar k}.
\end{aligned}
\end{equation}
其中，费米参数 $\alpha_-$、$\alpha_+$、$\widetilde\alpha_-$ 与 $\widetilde\alpha_+$，分别是 $L_1^{-1}$、$\overline{L_3}^{\,-1}$、$L_2^{-1}$ 与 $\overline{L_4}^{\,-1}$ 的截面。

\SourceParagraph{ch03-3.3.2-010}{164}
在 $L_1=K_\Sigma^{1/2}$ 的特殊情形中，可以将费米场组合成一个取值于 $X$ 的实切丛的 Dirac 旋量。费米参数 $\alpha_+,\widetilde\alpha_+$ 与 $\alpha_-,\widetilde\alpha_-$，分别为右手与左手旋量。这对应于通常的、未经扭曲的 $N=(2,2)$ 超对称 sigma 模型。

\SourceParagraph{ch03-3.3.2-011}{164}
可以定义守恒的超流 $G$，使它们在以下意义上，生成对称变换\BookEqRef{eq:3.87}{3.87}：
\begin{equation}\label{eq:3.88}
\delta W=-\mathrm{i}[Q(\alpha),W]_{\mathrm{s}},
\qquad Q(\alpha)=\int\mathrm{d}\sigma_2\,G\alpha.
\end{equation}
这里，$W$ 可以是任意算符，括号按照算符的奇偶性取分次对易子。由于有四个超对称，就有四个独立超流，记为 $G_+$、$\widetilde G_+$、$G_-$ 与 $\widetilde G_-$。此外，还有一个 $U(1)$ 流 $J_\mu$，其全纯与反全纯分量，分别记为 $J$ 与 $\bar J$。在经典层面，能动张量、超流与 R 流为
\begin{equation}\label{eq:3.89}
\begin{aligned}
T(z)&=-g_{i\bar\jmath}\partial\phi^i\partial\phi^{\bar\jmath}
+\tfrac12g_{i\bar\jmath}\psi_+^i\partial\psi_+^{\bar\jmath}
+\tfrac12g_{i\bar\jmath}\psi_+^{\bar\jmath}\partial\psi_+^i,\\
G_+(z)&=\tfrac12g_{i\bar\jmath}\psi_+^i\partial\phi^{\bar\jmath},\\
\widetilde G_+(z)&=\tfrac12g_{i\bar\jmath}\psi_+^{\bar\jmath}\partial\phi^i,\\
J(z)&=\tfrac14g_{i\bar\jmath}\psi_+^i\psi_+^{\bar\jmath}.
\end{aligned}
\end{equation}
反全纯的 $\bar T(\bar z)$、$G_-(\bar z)$、$\widetilde G_-(\bar z)$ 与 $\bar J(\bar z)$，具有类似的表达式。注意，左手超流是全纯的，右手超流是反全纯的。这是 sigma 模型在经典层面具有超共形对称的结果。如前所述，当度规与 B 场的 beta 函数为零时，量子化后的理论仍具有这一对称。

\SourceParagraph{ch03-3.3.2-012}{165}
量子化时，必须使这些算符能够经过重整化，同时保持 $(2,2)$ 代数。可以在 $\alpha'$ 展开的所有阶证明这一点；超场形式体系的存在，使这项工作相对容易\BookCite{ref010}。

\subsection{\texorpdfstring{$(2,2)$}{(2,2)} 超共形代数}\label{sec:3.3.3}

\SourceParagraph{ch03-3.3.3-001}{165}
弦态属于超共形代数的表示。在经典层面，这是由变换\BookEqRef{eq:3.87}{3.87}生成的对称代数，其中，参数 $\alpha_-$ 与 $\widetilde\alpha_-$ 取为全纯截面，而 $\alpha_+$ 与 $\widetilde\alpha_+$ 取为反全纯截面。这些超对称变换的反对易，生成共形代数以及 $U(1)\times U(1)$ R 对称变换。

\SourceParagraph{ch03-3.3.3-002}{165}
为确定起见，先关注全纯模式；在另行说明以前，省去全纯超流 $G_+$ 与 $\widetilde G_+$ 的手征下标。写出 $N=2$ SCA 的对易关系以前，回忆，在 $N=2$ 超共形场论中，圆柱上的 Hilbert 空间，依赖于一个 $\lambda\in\mathbb{C}^*$ 的选择，它标记线丛 $L_1$ 的同构类。也写为
\begin{equation}\label{eq:3.90}
\lambda=\mathrm{e}^{2\pi\mathrm{i}a}.
\end{equation}
在这样的扇区中，流的模式展开为
\[
\begin{aligned}
T(z)&=\sum_{n\in\mathbb{Z}}L_nz^{-n-2},\\
G(z)&=\sum_{n\in\mathbb{Z}}G_{n-a}z^{-n+a-3/2},\\
\widetilde G(z)&=\sum_{n\in\mathbb{Z}}\widetilde G_{n+a}z^{-n-a-3/2},\\
J(z)&=\sum_{n\in\mathbb{Z}}J_nz^{-n-1}.
\end{aligned}
\]
有两个特殊情形尤其重要：Ramond 扇区，$\lambda=1$ 或 $a=0$；Neveu--Schwarz 扇区，$\lambda=-1$ 或 $a=1/2$。它们通常分别简称 R 扇区与 NS 扇区。

\SourceParagraph{ch03-3.3.3-003}{166}
\begin{definition*}\label{def:3.n2-superconformal-algebra}
用这些模式表示，$N=2$ SCA 的对易关系为
\begin{equation}\label{eq:3.91}
\begin{aligned}
[L_m,L_n]&=(m-n)L_{m+n}
+\frac{\widehat c}{4}m(m^2-1)\delta_{m+n,0},\\
[J_m,J_n]&=\widehat c\,m\delta_{m+n,0},\\
[L_n,J_m]&=-mJ_{m+n},\\
[L_n,G_{m-a}]&=(n/2-m+a)G_{m+n-a},\\
[L_n,\widetilde G_{m+a}]&=(n/2-m-a)\widetilde G_{m+n+a},\\
[J_n,G_{m-a}]&=G_{m+n-a},\\
[J_n,\widetilde G_{m+a}]&=-\widetilde G_{m+n+a},\\
\{G_{n-a},\widetilde G_{m+a}\}
&=2L_{m+n}+(n-m-2a)J_{m+n}\\
&\quad+\widehat c\left((n-a)^2-\tfrac14\right)\delta_{m+n,0}.
\end{aligned}
\end{equation}
同类超流模式之间的反对易子为零。
\end{definition*}

\SourceParagraph{ch03-3.3.3-004}{166}
与 Virasoro 代数一样，可以用这一表述定义最高权表示。不可约幺正最高权表示，可以完整分类\BookCite{ref054}。这里只列出后面需要的几项结果。

\SourceParagraph{ch03-3.3.3-005}{166}
首先，SCA 的 Cartan 子代数，由三个元素 $L_0$、$\widehat c$ 与 $J_0$ 生成。其中两个，与讨论 Virasoro 代数时相同；特别地，$L_0$ 的本征值 $h$ 相同。对于中心荷，通常采用另一个归一化，$\widehat c=c/3$，使 sigma 模型的 $\widehat c$，等于 $X$ 的复维数。新的 Cartan 元素，或者说守恒荷，是 $J_0$。其本征值 $q$ 是 R 荷，有时也称为“$U(1)$ 荷”。

\SourceParagraph{ch03-3.3.3-006}{166}
闭弦同时具有左行与右行的 $N=2$ 代数。因此，一个闭弦态，具有左行权重与荷，记为 $h_L$ 与 $q_L$，以及右行权重与荷，记为 $h_R$ 与 $q_R$。最后，可以使 $a_L$ 与 $a_R$ 彼此独立。在超弦理论中，正是这样区分 NS--NS、NS--R、R--NS 与 R--R 扇区；不过，这里的讨论并不主要依靠这一点。

\SourceParagraph{ch03-3.3.3-007}{166}
正如\BookSectionRef{sec:3.2.3.4}{3.2.3.4}中所述，荷 $(q_L,q_R)$ 在 CFT 上定义双分次。分解成 $a$ 扇区，也给出一个 $\mathbb{R}/\mathbb{Z}$ 分次，或者双分次；不过，此时按加法守恒的是 $a+1/2$。换句话说，NS 算符的作用，保持态的 $a$ 值，而 R 算符及其他扇区的算符，则改变它。这是因为，$L_1$ 与 $L_2$ 用来推广自旋结构，因此，自然的选择，是 $L_1\cong L_2\cong K_\Sigma^{1/2}$。更具体地，需要确定，在什么条件下，\BookSectionRef{sec:2.1.6}{2.1.6}中定义的从 $S^1\sqcup S^1$ 到 $S^1$ 的配边存在，同时保持 $L_1\otimes L_2\cong K_\Sigma$。最简单的相容情形，就是 $L_1\cong L_2\cong K_\Sigma^{1/2}$ 所给出的通常自旋结构。

\subsection{拓扑扭曲与手征代数}\label{sec:3.3.4}

\SourceParagraph{ch03-3.3.4-intro-001}{167}
现在说明，\BookSectionRef{sec:3.3.1}{3.3.1}的讨论，如何在 $N=2$ SCA 中实现。将能动张量的改进\BookEqRef{eq:3.83}{3.83}写成 $\Sigma$ 上的复坐标形式，并选取使左行 $G$ 与右行 $\widetilde G$ 成为1形式的符号约定，得到
\[
T'=T-\tfrac12\partial J,\qquad
\bar T'=\bar T+\tfrac12\bar\partial\bar J.
\]
这提示，可以定义扭曲后的 Virasoro 生成元。在 Ramond 圆柱的整数模式表述中，它们为
\[
L'_m=L_m-\frac m2J_m-\frac c{24}\delta_{m,0},
\]
满足
\begin{equation}\label{eq:3.92}
[L'_m,G_n]=-nG_{m+n}.
\end{equation}
因此，零模 $G_0$ 在扭曲后的能动张量生成的共形变换下守恒，这提示选取
\begin{equation}\label{eq:3.93}
Q_L=G_{L,0}.
\end{equation}

\SourceParagraph{ch03-3.3.4-intro-002}{167}
对 $Q_R$ 可以作类似论证，从而选择 $G_{R,0}$ 或 $\widetilde G_{R,0}$。左右行部分的组合，留到\BookSectionRef{sec:3.3.4.4}{3.3.4.4}讨论；目前，写成 $Q=G_0$。

\SourceParagraph{ch03-3.3.4-intro-003}{167}
注意，$G_0$ 属于 Ramond 扇区。还可以通过将定义改写为
\[
Q_L=\frac1{2\pi\mathrm{i}}\oint\mathrm{d}z\,s(z)G(z)
\]
来理解这一选择为何自然。这里，$s(z)$ 是 $L_1^{-1}$ 的截面；需要它，是因为 $G(z)\,\mathrm{d}z$ 是取值于 $L_1$ 的全纯1形式。一般地，只有在 $L_1$ 平凡，即相应的 Ramond 表述中，才能取一个全局定义、处处非零的这样的截面。

\subsubsection{Ramond 基态}\label{sec:3.3.4.1}

\SourceParagraph{ch03-3.3.4.1-001}{167}
现在，将拓扑理论的 Hilbert 空间，定义为 $Q$ 上同调，并尝试沿用\BookSectionRef{sec:3.1.1.1}{3.1.1.1}中的其他定义。例如，考虑\BookEqRef{eq:3.91}{3.91}中的以下关系，就可以看出它与 Hodge 理论的联系：
\begin{equation}\label{eq:3.94}
\{G_0,\widetilde G_0\}=2L_0-\frac c{12}.
\end{equation}
用那里相同的论证，可以为 $Q$ 上同调类 $s$，找到一个典范代表，即满足 $\widetilde G_0|s\rangle=0$ 的态 $|s\rangle$。这个态具有 $h=c/24$。此外，由于 $\widetilde G_0=G_0^\dagger$，\BookSectionRef{sec:3.2.2.1}{3.2.2.1}的论证可以用来说明，在 Ramond 扇区中，
\[
h\ge\frac c{24}.
\]
因此，拓扑理论的 Hilbert 空间，由能量最小的态组成，即 Ramond 基态。

\subsubsection{手征算符}\label{sec:3.3.4.2}

\SourceParagraph{ch03-3.3.4.2-001}{168}
接下来，希望定义上同调代数的类似物。不过，正如已经讨论的那样，Ramond 算符并不构成一个代数；事实上，两个 Ramond 算符的乘积，是 NS 算符。

\SourceParagraph{ch03-3.3.4.2-002}{168}
由于 NS 算符构成一个封闭的子代数，它们在这方面更有希望。不过，仍可能担心，算符乘积中通常出现的奇异性。

\SourceParagraph{ch03-3.3.4.2-003}{168}
此外，限制到上同调，也不是那么简单。$G(z)$ 的每个模式都是幂零的，因此，可以尝试定义
\[
Q_{\mathrm{NS}}=G_{-1/2},
\]
并取其上同调。不过，模式数 $-1/2$ 的选择，乍看似乎是任意的；若把它直接代入 Ramond 模式的关系\BookEqRef{eq:3.92}{3.92}，它也并非不变。因此，NS 表述需要同时采用相应的模式重标记。在平面的 NS 扭曲中，$L'_m=L_m-(m+1)J_m/2$，而 $G_{-1/2}$ 正好是扭曲后权重为1的超流的零模。

\SourceParagraph{ch03-3.3.4.2-004}{168}
尽管如此，进行这一构造仍然很有用。
\begin{definition*}\label{def:3.chiral-primary-state}
手征初级态，是满足
\[
G_{-1/2}|h,q\rangle=0
\]
的超共形初级态。
\end{definition*}
由初级条件（见\BookSectionRef{sec:3.2.2.2}{3.2.2.2}）可推出 $\widetilde G_{1/2}|h,q\rangle=0$。因此，它们确实是 Hodge 理论意义下的典范代表。

\begin{proof}[权重与 R 荷的界]
\SourceParagraph{ch03-3.3.4.2-005}{168}
\BookEqRef{eq:3.94}{3.94}的类似关系为
\[
\{\widetilde G_{1/2},G_{-1/2}\}=2L_0-J_0.
\]
在最高权态 $|h,q\rangle$ 上取期望值，得到
\[
\langle h,q|\widetilde G_{1/2}G_{-1/2}|h,q\rangle=2h-q.
\]
由于 $\widetilde G_m^\dagger=G_{-m}$，左端必须非负，因此
\begin{equation}\label{eq:3.95}
h\ge\frac q2.
\end{equation}
等号
\begin{equation}\label{eq:3.96}
h=\frac q2
\end{equation}
恰在手征态上成立。由于 $h\ge0$，还得到，手征态满足
\begin{equation}\label{eq:3.97}
q\ge0.
\end{equation}

\SourceParagraph{ch03-3.3.4.2-006}{168}
如果改为要求
\[
\widetilde G_{-1/2}|h,q\rangle=0,
\]
就定义了反手征初级态。此时，考虑
\[
\{G_{1/2},\widetilde G_{-1/2}\}=2L_0+J_0,
\]
还可以导出
\[
h\ge-\frac q2,
\]
等号在反手征初级态上成立。

\SourceParagraph{ch03-3.3.4.2-007}{169}
最后，考虑
\[
\{\widetilde G_{3/2},G_{-3/2}\}
=2L_0-3J_0+2\widehat c,
\]
就可以得到手征算符的 R 荷上界：
\begin{equation}\label{eq:3.98}
q\le\widehat c.
\end{equation}
\end{proof}

\subsubsection{手征代数}\label{sec:3.3.4.3}

\SourceParagraph{ch03-3.3.4.3-001}{169}
手征算符的关键特征，是它们在算符乘积下具有良好的封闭性，构成手征代数；较不精确地，也称“手征环”。由\BookEqRef{eq:3.95}{3.95}可知，两个手征算符的乘积具有以下形式：
\begin{equation}\label{eq:3.99}
O_1(z_1)O_2(z_2)
\longrightarrow O_3(z_1)+O\bigl((z_1-z_2)^\epsilon\bigr),
\end{equation}
其中，$\epsilon>0$，$h_3=h_1+h_2$；$O_3$ 可以代表相应手征算符的线性组合。

\SourceParagraph{ch03-3.3.4.3-002}{169}
注意，这一展开主导项的系数，与算符的位置无关。此外，次主导项不满足\BookEqRef{eq:3.96}{3.96}，不能是手征初级算符；作为 $Q$ 闭乘积的非手征部分，它们在这里是 $Q$ 恰当的。因此，限制到 $Q$ 上同调以后，算符乘积代数就与算符的位置无关，正如 TFT 所要求的那样。

\SourceParagraph{ch03-3.3.4.3-003}{169}
当然，也可以对反手征算符作类似构造；这时取 $Q=\widetilde G$，而不是 $Q=G$。与范数正性有关的物理相容条件，即 CPT 不变性，意味着，所得“反手征环”是手征环的复共轭，因此不提供额外信息。

\SourceParagraph{ch03-3.3.4.3-004}{169}
这样，$N=2$ SCFT 的 $Q$ 上同调，具有一个超交换代数结构。这是下面定义的特殊情形。
\BookStatementNumber{2}
\begin{definition}\label{def:3.2}
设一个 QFT 具有满足 $Q^2=0$ 的超对称 $Q$。与它相关的 $Q$ 手征代数，是 $Q$ 的上同调，配以从算符乘积展开诱导的乘法。
\end{definition}

\SourceParagraph{ch03-3.3.4.3-005}{169}
由于历史原因，$Q$ 常称为“BRST 荷”。

\SourceParagraph{ch03-3.3.4.3-006}{169}
对于这里的目的，使用 $Q=Q_{\mathrm{NS}}$ 的这一定义就足够了。不过，这也引出了一个问题：如何与 Ramond 扇区，以及定义\BookEqRef{eq:3.93}{3.93}相联系？\BookSectionRef{sec:3.3.5.1}{3.3.5.1}还将进一步讨论。

\subsubsection{A 模型与 B 模型}\label{sec:3.3.4.4}

\SourceParagraph{ch03-3.3.4.4-001}{169}
现在，考虑一个同时具有左右行扇区的 $(2,2)$ SCFT。算符乘积具有类似\BookEqRef{eq:3.99}{3.99}的展开，不过，此时的次主导项，对坐标既有全纯依赖，也有反全纯依赖。要约化为 TFT，需要消去这种依赖。

\SourceParagraph{ch03-3.3.4.4-002}{169}
仍然可以通过取上同调来做到这一点，不过，现在涉及一个选择。假定通过取左行超荷 $Q_L=G_L$ 的上同调，消去全纯依赖。接着，可以取 $Q_R$ 的上同调，消去反全纯依赖；然而，$Q_R$ 可以选为 $G_R$ 或 $\widetilde G_R$。由于选定 $G_L$，已经不再保持 CPT 共轭下的对称处理，这两种选择将给出不同的、一般并不同构的代数。

\SourceParagraph{ch03-3.3.4.4-003}{170}
由于左右行 $N=2$ SCA 彼此分次对易，这一双重上同调，可以用总超荷 $Q_L+Q_R$ 的上同调表述。因此，可以选用下列任何一个算符，作为 BRST 荷：
\begin{equation}\label{eq:3.100}
\begin{aligned}
Q_A&=G_{+,0}+\widetilde G_{-,0},\\
\widetilde Q_A&=\widetilde G_{+,0}+G_{-,0},\\
Q_B&=\widetilde G_{+,0}+\widetilde G_{-,0},\\
\widetilde Q_B&=G_{+,0}+G_{-,0}.
\end{aligned}
\end{equation}
这里使用 Ramond 表述；NS 表述采用相应的半整数模式。由于 sigma 模型具有左右对称性，$Q$ 与 $\widetilde Q$ 的手征代数实际上同构；不过，$Q_A$ 与 $Q_B$ 的手征代数，一般并不同构。它们分别定义 A 拓扑模型与 B 拓扑模型。

\SourceParagraph{ch03-3.3.4.4-004}{170}
可以将关系\BookEqRef{eq:3.96}{3.96}适配到这一情形，将 $q$ 定义为下列算符的荷，即本征值：
\begin{equation}\label{eq:3.101}
\begin{aligned}
J_A&=J_{L,0}-J_{R,0},&&\text{A 模型},\\
J_B&=-J_{L,0}-J_{R,0},&&\text{B 模型}.
\end{aligned}
\end{equation}
按照这一定义，两种情形都满足
\[
[J_X,Q_X]=Q_X,\qquad X\in\{A,B\}.
\]

\SourceParagraph{ch03-3.3.4.4-005}{170}
由此，A 模型的 Hilbert 空间，可以经谱流，与满足 $h=q/2$、$\bar h=-\bar q/2$ 的超共形初级态空间相识别，也就是左手手征、右手反手征的态空间。同样，B 模型的 Hilbert 空间，与左右手均反手征的初级态空间相识别。

\SourceParagraph{ch03-3.3.4.4-006}{170}
因此，扭曲后的理论具有 $U(1)$ 对称，或者分次。相应的荷，通常称为鬼数；其定义使 BRST 荷的鬼数为1\footnote{这在本质上，就是量子力学讨论中出现的同一个费米子数与 $\mathbb{Z}$ 分次。读者也许会疑惑，为什么物理术语中会突然出现这种带神秘色彩的名称。\BookSectionRef{sec:3.3.5.4}{3.3.5.4}将揭晓其来历。}。

\subsubsection{与扭曲的一般讨论比较}\label{sec:3.3.4.5}

\SourceParagraph{ch03-3.3.4.5-001}{170}
本小节借助 SCFT 的特殊性质，简化了讨论。现在，将它与\BookSectionRef{sec:3.3.1}{3.3.1}中的一般讨论联系起来；这样做，部分也是为了说明，其中哪些内容适用于更一般的二维超对称理论。

\SourceParagraph{ch03-3.3.4.5-002}{170}
事实上，刚才讨论的构造，与先前的构造等价。不过，在\BookEqRef{eq:3.83}{3.83}中，用来修改能动张量的 $U(1)$ 流，对 A 模型与 B 模型有所不同。暂时将它记为 $\mathcal{J}_\mu$，以区别于原先的 R 流 $J_\mu$；二者先验并不相同。采用前面的手征荷约定，在 $\Sigma$ 的复坐标中，改进后的能动张量写为
\[
T'=T-\tfrac12\partial\mathcal{J},\qquad
\bar T'=\bar T+\tfrac12\bar\partial\overline{\mathcal{J}}.
\]
重复\BookSectionRef{sec:3.3.4}{3.3.4}以下的讨论，可以看出，当 $\mathcal{J}=J$、$\overline{\mathcal{J}}=\bar J$ 时，A 模型的 BRST 算符 $Q_A=G_{+,0}+\widetilde G_{-,0}$，在任意共形变换下不变。同样容易看出，在这一选择下，其他三个 BRST 算符并非都具有这种共形不变性。要使 $\widetilde Q_A$ 共形不变，必须反转 $\mathcal{J}_\mu$ 的符号；要使 $Q_B$ 共形不变，则取 $\mathcal{J}=-J$、$\overline{\mathcal{J}}=\bar J$。在实正交坐标及相应约定下，这等价于 $\mathcal{J}_\mu=-\epsilon_{\mu\nu}J^\nu$。这一流通常称为轴向 R 流，而 sigma 模型中通常的 R 流，称为向量 R 流。$N=2$ 超共形不变性，要求向量与轴向 R 流都守恒，因此，共形不变的 $N=2$ sigma 模型，可以扭曲为 A 模型或 B 模型。不过，如果考虑 beta 函数非零的更一般的 $N=2$ sigma 模型，例如，靶空间 $X$ 为 $c_1(X)\ne0$ 的 K\"ahler 流形，那么，通常只有向量 R 流在量子层面守恒，也就只能定义 A 模型。

\SourceParagraph{ch03-3.3.4.5-003}{171}
注意，失去共形不变性时，轴向 R 荷一般不守恒，因此，在非 Calabi--Yau 靶空间上，A 模型通常不具有完整的 $\mathbb{Z}$ 分次。不过，它仍然具有 $\mathbb{Z}_2$ 分次，因为费米子与玻色子的区别仍然存在。

\subsubsection{\texorpdfstring{$U(1)$}{U(1)} 子代数与谱流}\label{sec:3.3.4.6}

\SourceParagraph{ch03-3.3.4.6-001}{171}
注意，除去一个整体常数，关系
\[
[J_m,J_n]=\widehat c\,m\delta_{m+n,0}
\]
就是 $U(1)$ 流代数\BookEqRef{eq:3.46}{3.46}。因此，可以使用\BookSectionRef{sec:3.2.4}{3.2.4}中的因子分解，用一个自由玻色子，描述任意 $(2,2)$ SCFT 的“$U(1)$ 扇区”。

\SourceParagraph{ch03-3.3.4.6-002}{171}
为处理这个常数，定义 $J_n=\sqrt{\widehat c}\,\alpha_n$；等价地，采用\BookSectionRef{sec:3.2.3.2}{3.2.3.2}的实自由玻色子归一化，写成
\begin{equation}\label{eq:3.102}
J=\mathrm{i}\sqrt{2\widehat c}\,\partial\varphi.
\end{equation}
这给出
\begin{equation}\label{eq:3.103}
q=\sqrt{\widehat c}\,p,
\end{equation}
将先前讨论中定义的“$U(1)$ 荷” $p$，与当前使 $G$、$\widetilde G$ 的荷分别为 $+1$、$-1$ 的定义联系起来。

\SourceParagraph{ch03-3.3.4.6-003}{171}
因此，可以把 $(2,2)$ SCFT $X$ 的 Hilbert 空间，写成\BookEqRef{eq:3.70}{3.70}那样的分次直和：
\[
\mathcal{H}^X=\widehat{\bigoplus}_{q_L,q_R}
\mathcal{H}^{X'}_{q_L,q_R}\widehat\otimes V_{q_L,q_R}.
\]
这里以 R 荷重新标记自由玻色子表示；连续谱时，仍需使用直积分。$X$ 中的一般手征态，在两个因子中，都可能分解出非平凡的态。当然，真空在两个因子中都平凡。

\SourceParagraph{ch03-3.3.4.6-004}{172}
是否存在这样的手征算符：在 $X'$ 中为恒等算符，在 $U(1)$ 扇区中却非平凡？利用 $h_X=h_{X'}+h_{U(1)}$ 与\BookSectionRef{sec:3.2}{3.2}的\BookRef{axiom:3.3}{公理3}，只需判断，维数 $h_X$ 是否等于 $h_{U(1)}$。这样的算符，必须同时满足 $h=q/2$ 与\BookEqRef{eq:3.52}{3.52}。考虑\BookEqRef{eq:3.103}{3.103}，后一个关系成为
\[
h=\frac{q^2}{2\widehat c}+N,
\qquad N\in\mathbb{Z}_{\ge0}.
\]
对于相应的流代数初级算符，$N=0$；除恒等算符 $q=0$ 外，非平凡的解为 $q=\widehat c$。

\SourceParagraph{ch03-3.3.4.6-005}{172}
将这样一个具有 $q_L=\widehat c$、$q_R=0$ 的局部算符，记为 $\Omega_L$。在具有全纯体积形式的 sigma 模型中，可以将这个算符，与全纯 $(d,0)$ 形式相识别。具体地，
\begin{equation}\label{eq:3.104}
\Omega_L(z)=\Omega_{i_1\cdots i_d}
\psi_+^{i_1}\cdots\psi_+^{i_d}.
\end{equation}
类似地，可以定义 $(0,d)$ 形式的对应物 $\Omega_R$。在大体积的半经典极限中，这一识别很清楚，因为没有其他具有这些荷与维数的候选算符。在相应的量子理论中，它继续成立，是一种“非重整化定理”，与\BookSectionRef{sec:3.2.4}{3.2.4}的因子分解的非平凡结果相关。

\SourceParagraph{ch03-3.3.4.6-006}{172}
在用于弦理论紧化的 SCFT 中，甚至可以找到以下关系的算符类似物：全纯 $(d,0)$ 形式，与协变常值旋量之间满足
\[
\Omega_{i_1\cdots i_d}
=\epsilon^\dagger\Gamma_{i_1\cdots i_d}\epsilon.
\]
相应地，存在一个满足
\begin{equation}\label{eq:3.105}
\Upsilon_L^2=\Omega_L
\end{equation}
的算符 $\Upsilon_L$，其 $U(1)$ 荷为 $q_L=\widehat c/2$。这里的平方，表示相应的 OPE 复合。这一局部算符，称为谱流算符。它在 Ramond 与 Neveu--Schwarz 扇区之间起交织作用，并负责超弦理论中的时空超对称。

\SourceParagraph{ch03-3.3.4.6-007}{172}
这一构造的另一个应用，是通过给出 R 扇区与 NS 扇区之间的显式同构，解决前面提到的态—算符对应问题。最简单的表述，不使用 $\Upsilon$，而是在分解\BookEqRef{eq:3.70}{3.70}中，将谱流定义为交织映射
\begin{equation}\label{eq:3.106}
S_{\delta q_L}\colon V_{q_L,q_R}
\longrightarrow V_{q_L+\delta q_L,q_R}.
\end{equation}
它在基\BookEqRef{eq:3.47}{3.47}中，只平移最高权的荷 $q$，保持所有占据数 $N_n$ 不变。于是，NS--R 对应，就是 $\delta q_L=\widehat c/2$ 的谱流。完整 SCA 的权重随之按 $h\mapsto h+\eta q+\widehat c\eta^2/2$ 改变，其中 $\delta q=\widehat c\eta$。

\SourceParagraph{ch03-3.3.4.6-008}{172}
重复这一操作，得到 $\delta q=\widehat c$ 的谱流，将 NS 扇区映到自身，或者将 R 扇区映到自身；这对应于\BookEqRef{eq:3.105}{3.105}中 $\Omega$ 的作用。稍后将回到这一点。

\subsection{Frobenius 结构与扭曲}\label{sec:3.3.5}

\SourceParagraph{ch03-3.3.5-intro-001}{173}
回忆\BookChapterRef{2}，代数上的 Frobenius 结构，是一个线性泛函，即迹；将它与代数乘法复合，便诱导一个非退化的双线性配对。

\SourceParagraph{ch03-3.3.5-intro-002}{173}
在 CFT 中，可以使用态—算符对应，得到一个自然的迹：
\[
\operatorname{tr}O\equiv\langle0|O\rangle,
\]
其中，$|0\rangle$ 是真空。这也就是球面上 $O$ 的一点函数。借助幺正 CFT 的公理，将这一迹与 OPE 复合，并对其中一个算符作相应共轭，可以恢复 Hilbert 空间的内积。

\SourceParagraph{ch03-3.3.5-intro-003}{173}
虽然可以将这一迹，限制到 SCFT 的手征初级算符代数，但所得配对高度退化，不能定义 Frobenius 结构。在 TFT 中，改为使用\BookSectionRef{sec:3.3.4.6}{3.3.4.6}定义的典范态 $\Omega=\Omega_L\otimes\Omega_R$，替代真空 bra。给定 $\Omega$，其中涉及归一化的选择，定义
\begin{equation}\label{eq:3.107}
\langle O\rangle_{\mathrm{TFT}}
=\langle\Omega|O|0\rangle_{\mathrm{SCFT}}.
\end{equation}
这就定义了配对
\begin{equation}\label{eq:3.108}
\langle O_iO_j\rangle_{\mathrm{TFT}}.
\end{equation}
利用 $U(1)$ 子代数的因子分解，它与 $X'$ 理论中的内积成比例（记号同\BookSectionRef{sec:3.2.4}{3.2.4}），因而非退化。

\SourceParagraph{ch03-3.3.5-intro-004}{173}
配对\BookEqRef{eq:3.108}{3.108}，对于左手与右手 R 荷，分别具有次数 $\widehat c$，给出 Kodaira--Serre 对偶关系的 SCFT 推广：
\begin{equation}\label{eq:3.109}
H^{0,p}_{\bar\partial}(\Omega^q)
\cong H^{0,d-p}_{\bar\partial}(\Omega^{d-q})^*.
\end{equation}

\SourceParagraph{ch03-3.3.5-intro-005}{173}
显然，这与上一小节末尾提到的 $\delta q=\widehat c$ 的谱流相关。注意，由于\BookEqRef{eq:3.109}{3.109}中 $U(1)$ 荷的符号反转，这本身并不是 SCFT 的对称性。将它和 $J\to-J$ 复合，或者再迭代一次，都可以得到对称性；由此可知，以 $\delta q=2\widehat c$ 作的总体谱流，是 SCFT 的对称性。\BookSectionRef{sec:5.3.4}{5.3.4}将使用这一点。

\subsubsection{泛函积分中的扭曲}\label{sec:3.3.5.1}

\SourceParagraph{ch03-3.3.5.1-001}{173}
虽然\BookEqRef{eq:3.107}{3.107}看起来有些像临时规定，但可以从反常得到更具物理意义的推导；它还说明，谱流如何自然地出现在扭曲后的理论中。关键是，仔细考虑扭曲\BookEqRef{eq:3.83}{3.83}对泛函积分的作用。回忆，在算符语言中，这通过依赖 R 流 $J_\mu$ 的项，修改能动张量。由于能动张量决定理论与世界面曲率的耦合，球面上的扭曲理论，与未扭曲理论不同。在泛函积分方法中，这表现为，费米子取值于经过适当 $\operatorname{spin}^c$ 结构扭曲的丛，正如\BookSectionRef{sec:3.3.2}{3.3.2}中所述。

\SourceParagraph{ch03-3.3.5.1-002}{174}
在扭曲理论中，R 流具有“引力反常”，也就是说，R 荷在弯曲世界面上并不守恒。例如，可以证明，球面上的一点函数非零，当且仅当算符的总 R 荷为 $2\widehat c$，此时它与\BookEqRef{eq:3.107}{3.107}成比例。更一般地，一个区域 $R\subset\Sigma$，包含“背景 R 荷” $\delta q(R)$，它正比于度规的曲率2形式 $\mathcal{R}$ 的积分：
\begin{equation}\label{eq:3.110}
\delta q(R)=-\frac{\widehat c}{2\pi}\int_R\mathcal{R}.
\end{equation}
例如，在 Euler 示性数为 $\chi$ 的闭曲面上，关联函数非零的必要条件，是算符的总 $U(1)$ 荷为 $\chi\widehat c$。此外，在将 $\Sigma$ 分解为几何态射时，计入这一背景 R 荷的影响，就会发现，中间的 Hilbert 空间上，会诱导各种谱流；它们将上面给出的 NS 与 R 扇区表述，交织为同一个 TFT。

\subsubsection{\texorpdfstring{$N=2$}{N=2} 超对称形变}\label{sec:3.3.5.2}

\SourceParagraph{ch03-3.3.5.2-001}{174}
这样的形变，对应于荷为零的局部2形式算符，其在 $\Sigma$ 上的积分是 $Q$ 闭的。事实上，可以按如下方式，将这些2形式，与前面讨论的标量算符相联系。

\SourceParagraph{ch03-3.3.5.2-002}{174}
设 $W_a$ 是荷为 $p$ 的局部算符。算符 $\mathrm{d}W_a$，其中 $\mathrm{d}$ 是世界面的 de Rham 算符，与其他算符的关联函数平凡，因为在 TFT 中，算符插入的位置并不重要。由此可知，它必须是 $Q$ 恰当的，即存在一个鬼数为 $p-1$ 的算符1形式 $W_a^{(1)}$，使得
\begin{equation}\label{eq:3.111}
\mathrm{d}W_a=\{Q,W_a^{(1)}\}.
\end{equation}
可以再次重复这一过程，令
\begin{equation}\label{eq:3.112}
\mathrm{d}W_a^{(1)}=\{Q,W_a^{(2)}\},
\end{equation}
其中，$W_a^{(2)}$ 是鬼数为 $p-2$ 的算符2形式。由于它是2形式，自然可以在 $\Sigma$ 上积分。因此，得到以下理论形变：
\begin{equation}\label{eq:3.113}
S\longmapsto S+t\int_\Sigma W_a^{(2)}\,\mathrm{d}^2z.
\end{equation}
为了保持分次，必须取 $p=2$。

\subsubsection{关联函数与预势}\label{sec:3.3.5.3}

\SourceParagraph{ch03-3.3.5.3-001}{174}
与\BookSectionRef{sec:3.2.5}{3.2.5}一样，现在希望为 TFT 中的关联函数定义生成函数。对于闭弦关联函数，它通常称为预势。标准定义，将前面的要素组合起来。最初，可能希望写成
\begin{equation}\label{eq:3.114}
F[t]=\left\langle
\exp\left(\sum_i t_i\int_\Sigma W_i^{(2)}\right)
\right\rangle.
\end{equation}
这里的积分算符，就是\BookEqRef{eq:3.113}{3.113}中定义的 $W_i^{(2)}$。不过，由于 $W_i^{(2)}$ 的 R 荷为零，直接按这一表达式取球面期望值，将因荷选择规则而平凡。

\SourceParagraph{ch03-3.3.5.3-002}{175}
要得到非平凡的期望值，必须补足总 R 荷 $q_L+q_R=2\widehat c$。一种看似临时规定的方法，是将\BookEqRef{eq:3.114}{3.114}指数展开中的 $n=\widehat c$ 个算符，改取为原来的标量 $W_i$；这就提供了总鬼数 $2\widehat c$。

\SourceParagraph{ch03-3.3.5.3-003}{175}
当然，随后必须追问，结果如何依赖于这 $n=\widehat c$ 个特殊算符的选择。一个颇不平凡的结论是，在所讨论的下降构造中，结果与这一选择无关；例如，文献\BookCite{ref119}第8节给出了证明。这可以通过 $N=2$ SCA 的性质理解，也可以更几何地，通过发展\BookSectionRef{sec:3.3.5.1}{3.3.5.1}中与扭曲的联系理解。因此，可以将\BookEqRef{eq:3.114}{3.114}看作预势的形式定义，隐含将鬼数反常 $2\widehat c$ 分配给各个插入算符、以得到非零结果的步骤。

\subsubsection{一段历史插叙}\label{sec:3.3.5.4}

\SourceParagraph{ch03-3.3.5.4-001}{175}
为什么 $U(1)_R$ 荷经常称为“鬼数”？这一术语的历史\BookCite{ref438}，始于 Feynman 在1963年对 Yang--Mills 理论量子化的研究。他注意到，原先用图示方法计算 QFT 的方法，不能直接用于这些理论，于是引入额外的虚构粒子来解决问题。这些粒子在物理上不可观测，最终得到“鬼粒子”的名称。同时，Faddeev、Popov、DeWitt 等人的工作表明，在具有非 Abelian 规范对称的理论中，鬼粒子是泛函积分图展开的不可避免的要素；它们表示从规范固定条件换元到规范参数时出现的 Jacobian。

\SourceParagraph{ch03-3.3.5.4-002}{175}
另一个重要例子，是弦理论的世界面理论，其规范对称是二维一般坐标不变性。从1980年代初 Friedan 与 Alvarez 的工作开始，Faddeev--Popov 方法就被用于弦的量子化\BookCite{ref159,ref009}。所得“世界面重参数化鬼场”，具有像\BookEqRef{eq:3.110}{3.110}那样反常的 $U(1)$ 对称，不过，其中系数 $\widehat c$ 被3替代。这一常数，与 $3=\dim_{\mathbb{C}}\operatorname{SL}(2,\mathbb{C})$ 相联系；这一群的射影形式，是 $\mathbb{CP}^1$ 的共形对称群。它还与亏格为 $g\ge2$ 的 Riemann 曲面模空间维数公式 $\dim_{\mathbb{C}}\mathcal{M}_g=3g-3$ 中的系数相联系。

\SourceParagraph{ch03-3.3.5.4-003}{175}
同样的结构，也出现在超弦理论中；前面讨论的大部分结构，正是在这一情形中，从物理上发现的。例如，无论从数学上还是物理上，用带三个标记点的球面来定义关联函数，比直接在没有标记点的 $\mathbb{CP}^1$ 上定义更自然，这就说明了\BookSectionRef{sec:3.3.5.3}{3.3.5.3}中 $\widehat c=3$ 情形的规定。更一般地，从 $\widehat c=3$ 的 TFT 出发，系统地让 R 荷起鬼数的作用，就可以定义任意亏格的拓扑弦理论；细节见文献\BookCite{ref119,ref242}。完整超弦理论，也有类似但复杂得多的讨论\BookCite{ref395}。

\subsection{\texorpdfstring{$N=2$}{N=2} 最小模型与 Gepner 模型}\label{sec:3.3.6}

\SourceParagraph{ch03-3.3.6-intro-001}{176}
与前面讨论的 Virasoro 表示论直接类似，$N=2$ SCA 的中心荷，也有一个临界值 $\widehat c=1$。低于它时，表示论变得简单，使人们能够分类并精确求解这些幺正 SCFT，即“$N=2$ 最小模型”。这里简要讨论这些结果，\BookSectionRef{sec:3.6.8}{3.6.8}将给出更详细的说明。

\SourceParagraph{ch03-3.3.6-intro-002}{176}
首先，小于1的允许中心荷，组成一个离散集合：
\[
\widehat c=1-\frac2{k+2},\qquad k\in\mathbb{Z}_{>0}.
\]
对于每个 $\widehat c$，A 系列手征初级态的允许 R 荷，也组成一个离散集合：
\[
q=\frac n{k+2},\qquad n=0,1,\ldots,k.
\]
这些 $N=2$ SCA 表示，可以组合成完整的 SCFT，必须满足模不变性约束。这给出一个很简单的结果：$\widehat c<1$ 的 $N=2$ SCFT，与单边（ADE）Dynkin 图一一对应；其中 $k$ 是对偶 Coxeter 数，手征算符则与图的结点相对应。

\SourceParagraph{ch03-3.3.6-intro-003}{176}
起初，这一结果似乎颇为神秘。不过，很快人们就认识到，ADE 分类的原因在于，$N=2$ 最小模型，与同样按 ADE 分类的简单复奇点相对应。看出这一点的最简单方法，是考虑 \BookBibText{ref242}{MS1} 第13章与第16章讨论的“Landau--Ginzburg”构造。它从一个自由 sigma 模型出发，即\BookEqRef{eq:3.86}{3.86}中 $X\cong\mathbb{C}^n$、度规为 Euclidean 度规的理论，再通过超势修改作用量；超势是一个全纯函数 $W\colon X\to\mathbb{C}$。

\SourceParagraph{ch03-3.3.6-intro-004}{176}
这有两个主要的物理效果。首先，按照\BookSectionRef{sec:3.2.5}{3.2.5}的术语，所考虑的超势项是相关形变，因此，原则上可以通向一个具有 $\widehat c<n$ 的新 CFT。不过，单靠 RG 来刻画这一过程，并不容易。

\SourceParagraph{ch03-3.3.6-intro-005}{176}
然而，由于具有 $N=2$ 超对称，可以采取如下方法。可以证明，所得理论的手征环，是一个商环，有时称为 $W$ 的 Jacobian 环：
\[
\mathcal{R}\cong
\frac{\mathbb{C}[\phi_1,\ldots,\phi_n]}
{(\partial W/\partial\phi_1,\ldots,\partial W/\partial\phi_n)}.
\]
由于 CFT 中的手征环带有分次，要求 $W$ 为拟齐次函数。事实上，QFT 的定义，要求 $W$ 的单边 $U(1)$ 荷为1\footnote{更精确地，从超势项的超场形式 $\int\mathrm{d}^2\theta\,W$ 看，$W$ 的轴 R 荷为2。在 SCFT 中，它等于相等的左右 $U(1)$ 荷之和。}，从而确定这些分次。最后，荷为1的相应超对称下降形变是边缘的，因此，为了得到这里的最小模型，选择 $W$ 使这类模参数不存在，便导向简单奇点。将这些手征环，与从 SCFT 表示论导出的精确解比较，可以确认这一识别。

\SourceParagraph{ch03-3.3.6-intro-006}{177}
最简单的情形是 $W=\phi^{k+2}$，对应于最小模型的 $A_k$ 系列。很容易看出，这些结果与前面的陈述一致：手征环由 R 荷为 $1/(k+2)$ 的算符 $\phi$ 生成，而最大 R 荷 $\widehat c=k/(k+2)$，由 $\phi^k$ 实现。

\SourceParagraph{ch03-3.3.6-intro-007}{177}
由于前面已经讨论，$\widehat c$ 可以与 Calabi--Yau 的维数相识别，这提供了一个简单的物理定义，从某种意义上，得到具有某些非整数“维数”的 Calabi--Yau 的 SCFT 类似物。

\subsubsection{Gepner 模型与线性 sigma 模型}\label{sec:3.3.6.1}

\SourceParagraph{ch03-3.3.6.1-001}{177}
继续以前，应当说明，在这里引入 Gepner 模型，与逻辑顺序稍有不同，因为它们的现代解释，需要用到许多尚未讨论的话题。不过，仍在这里介绍，部分也是为了展示这段历史：Gepner 在1988年的工作\BookCite{ref175}，早于这里讨论的几乎所有其他工作。它也代表了单靠 SCFT 的代数方法能够达到的重要成果，因此，不依赖于其余讨论所借助的启发式物理论证。事实上，对 SCFT 的数学理解\BookCite{ref031,ref247,ref248}，正向足以将这些结果严格化的程度发展。

\SourceParagraph{ch03-3.3.6.1-002}{177}
考虑由 $W=\phi^5$ 定义的、层数 $k=3$ 的 A 系列模型，其中心荷为 $\widehat c=3/5$。为将它与通常的 Calabi--Yau 靶空间的 sigma 模型所得 SCFT 比较，可以尝试取五个这样的模型的张量积。由\BookEqRef{eq:3.91}{3.91}可知，将每个模型的生成元相加，例如
\[
T=\sum_{i=1}^5T_i,
\]
以及对应的 $G$、$\widetilde G$ 与 $J$，便得到中心荷为 $\widehat c=3$ 的对角 $N=2$ SCA。从某种意义上，层数为3的这个最小模型，是否可以看作“Fermat 五次超曲面的五次方根”？

\SourceParagraph{ch03-3.3.6.1-003}{177}
可以，不过，要精确表述，还存在一个细节。严格地说，$A_k$ 最小模型并不包含\BookEqRef{eq:3.105}{3.105}中的算符 $\phi^k=\Upsilon^2$，也不包含任何 $\phi^n$ 算符。完整理论包含的是满足模不变性约束的左右表示配对，可以记为 $\phi^n\bar\phi^m$，相应的左右 $U(1)$ 荷为 $m/(k+2)$ 与 $n/(k+2)$。在这里选用的对角 A 系列配对中，取 $m=n$，因而没有纯左荷标签 $(k,0)$。所以，五个模型的直接张量积，也没有对角 SCA 所需的那个纯左行 $\Upsilon^2$\footnote{这也意味着，最小模型或它们的乘积，不能被扭曲为二维 TFT。}。

\SourceParagraph{ch03-3.3.6.1-004}{177}
不过，Gepner 证明，这个最小模型乘积的一个 $\mathbb{Z}_5$ 轨形，包含所需的 $\Upsilon^2$ 算符，因此，可以将它与一个以五次 Calabi--Yau 超曲面为靶空间的 sigma 模型相识别。他还给出多种证据，表明，在 K\"ahler 模参数的某个“弦性”取值处，这个 SCFT，与靶空间为 Fermat 五次超曲面的 sigma 模型同构；在这一取值处，经典几何的考察可能失效。

\SourceParagraph{ch03-3.3.6.1-005}{178}
后来，人们通过线性 sigma 模型，更好地理解了这一构造，见 \BookBibText{ref242}{MS1} 第15章。线性 sigma 模型，引入一个复化的 K\"ahler 参数，可以用来在 Gepner 模型，或者 LG 轨形，与非线性 sigma 模型之间插值。因此，它展示了 TFT 中可见的全部结构，并为后续的物理发展，提供了更灵活的基础。

\SourceParagraph{ch03-3.3.6.1-006}{178}
尽管有这些优点，精确结果本身也十分有价值。后面将介绍 Gepner 模型边界态的结果\BookCite{ref402,ref069}；它们对于形成这里将讨论的一般图景，非常重要。另一个例子是，Gepner 模型的定向轨形，也已经有许多精确结果\BookCite{ref073}。这是一个相关话题，本书不作讨论；对于它，其他表述方法仍在发展中。

\section{闭弦的拓扑 sigma 模型}\label{sec:3.4}

\SourceParagraph{ch03-3.4-intro-001}{178}
现在，回到基于 $(2,2)$ 非线性 sigma 模型的讨论。不必先构造这个 SCFT，再约化到拓扑子扇区，而可以直接写出扭曲后的 sigma 模型作用量；这样，可以稍微简化讨论。

\subsection{A 模型}\label{sec:3.4.1}

\SourceParagraph{ch03-3.4.1-001}{178}
在 sigma 模型中，可以通过修改费米子取值的丛，进行扭曲。先从 A 模型开始。从现在起，一般将 A 模型的靶空间记为 $Y$，将 B 模型的靶空间记为 $X$\footnote{这一看似颠倒的约定，使后面某些代数几何小节的记号更为标准。}。

\SourceParagraph{ch03-3.4.1-002}{178}
修改费米子取值的丛，就“扭曲”了超共形场论。现在，不再采用\BookEqRef{eq:3.85}{3.85}的赋值，而取
\begin{equation}\label{eq:3.115}
\begin{aligned}
\chi^i=\psi_+^i&\in\Gamma(\phi^*T_Y),\\
\chi^{\bar\imath}=\psi_-^{\bar\imath}
&\in\Gamma(\phi^*\overline{T_Y}),\\
\psi_z^{\bar\imath}=\psi_+^{\bar\imath}
&\in\Gamma(K_\Sigma\otimes\phi^*\overline{T_Y}),\\
\psi_{\bar z}^i=\psi_-^i
&\in\Gamma(\overline{K_\Sigma}\otimes\phi^*T_Y).
\end{aligned}
\end{equation}
注意，在这一赋值下，作用量\BookEqRef{eq:3.86}{3.86}仍有意义，也就是说，它在世界面的转动下不变。

\SourceParagraph{ch03-3.4.1-003}{178}
现在，在\BookEqRef{eq:3.87}{3.87}中，取 $\alpha=\alpha_-=\widetilde\alpha_+$，并令 $\alpha_+=\widetilde\alpha_-=0$，所得对称满足 BRST 对称的条件；生成这一对称的算符 $Q$，满足 $Q^2=0$。

\SourceParagraph{ch03-3.4.1-004}{178}
将拓扑项的 B 场与 K\"ahler 参数，采用周期为1的上同调归一化后，作用量可以写为
\begin{equation}\label{eq:3.116}
S=\mathrm{i}\int_\Sigma\{Q,D\}
-2\pi\mathrm{i}\int_\Sigma\phi^*(B+\mathrm{i}\omega),
\end{equation}
其中
\begin{equation}\label{eq:3.117}
D=2\pi g_{i\bar\jmath}
\left(\psi_z^{\bar\jmath}\bar\partial\phi^i
+\partial\phi^{\bar\jmath}\psi_{\bar z}^i\right).
\end{equation}
$\omega=\mathrm{i}g_{i\bar\jmath}\mathrm{d}z^i\wedge\mathrm{d}\bar z^j$ 是 K\"ahler 形式，而 $B+\mathrm{i}\omega\in H^2(Y,\mathbb{C})$，是复化的 K\"ahler 类。

\SourceParagraph{ch03-3.4.1-005}{179}
由于复结构只出现在 $D$ 中，改变它，就使作用量发生一个 $Q$ 恰当的变化；在拓扑理论中，这种变化平凡。因此，这一拓扑 A 模型的关联函数，只依赖于复化的 K\"ahler 形式 $B+\mathrm{i}\omega$。事实上，由\BookEqRef{eq:3.116}{3.116}可知，它们只依赖于这一形式的上同调类。

\SourceParagraph{ch03-3.4.1-006}{179}
此外，由于泛函积分\BookEqRef{eq:3.14}{3.14}只依赖于 $\mathrm{e}^{-S}$，将 B 平移一个整上同调元素，并不会改变它。因此，A 模型的初步模空间\footnote{加入世界面瞬子修正以后，可以通过解析延拓，越过 K\"ahler 锥的边界。}，是 $Y$ 的复化 K\"ahler 模空间。其定义，是 $H^2(Y,\mathbb{C})/H^2(Y,\mathbb{Z})$ 中，$\omega$ 满足 K\"ahler 几何正性条件的区域\footnote{特别地，对于任意全纯曲线 $C\subset Y$，有 $\int_C\omega>0$。}。

\SourceParagraph{ch03-3.4.1-007}{179}
局部算符，是场 $\phi$ 与 $\psi$ 的一般函数。正如已经讨论的那样，$Q$ 上同调由标量算符生成；其中不能包含 $\psi_z^{\bar\imath}$ 或 $\psi_{\bar z}^i$。它们也不能包含场的导数 $\partial^n\phi$ 或 $\partial^n\chi$；可以用 SCA，将这些导数写成后裔项。

\SourceParagraph{ch03-3.4.1-008}{179}
于是，手征算符都可以写为
\begin{equation}\label{eq:3.118}
W[a]=a_{I_1I_2\cdots I_p}\chi^{I_1}\chi^{I_2}\cdots\chi^{I_p},
\end{equation}
其中，$a=a_{I_1I_2\cdots I_p}\mathrm{d}\phi^{I_1}\wedge\mathrm{d}\phi^{I_2}\wedge\cdots\wedge\mathrm{d}\phi^{I_p}$，是 $Y$ 上的 $p$ 形式。$I_n$ 是实坐标指标；在复化后的写法中，它们可以是全纯或反全纯指标。随后，可以计算得到
\begin{equation}\label{eq:3.119}
\{Q,W[a]\}=-W[\mathrm{d}a].
\end{equation}
也就是说，对于 A 模型，$Q$ 上同调就是 de Rham 上同调，算符空间为 $H^*(Y,\mathbb{C})$。

\subsubsection{关联函数}\label{sec:3.4.1.1}

\SourceParagraph{ch03-3.4.1.1-001}{179}
希望通过对映射 $\phi\colon\Sigma\to Y$ 积分，计算
\begin{equation}\label{eq:3.120}
\langle W_aW_bW_c\cdots\rangle
=\int[\mathrm{d}\phi\,\mathrm{d}\chi\,\mathrm{d}\psi]\,
\mathrm{e}^{-S}W_aW_bW_c\cdots.
\end{equation}
首先注意，\BookEqRef{eq:3.116}{3.116}的第二项，只依赖于类 $\phi_*[\Sigma]\in H_2(Y,\mathbb{Z})$，因此，可以将这一因子提出积分：
\[
Z=\sum_{\beta=\phi_*[\Sigma]}
\mathrm{e}^{2\pi\mathrm{i}\int_\beta(B+\mathrm{i}\omega)}
\int_{\phi_*[\Sigma]=\beta}[\mathrm{d}\phi\cdots]\,
\mathrm{e}^{-\mathrm{i}\int_\Sigma\{Q,D\}}.
\]
剩下的作用量项是 $Q$ 恰当的，因此，它的系数在拓扑关联函数中不重要。虽然第一反应，可能是将这个平凡的项置为零，但实际上采取相反的做法：用一个趋于无穷大的系数，放大它！在通常选取的积分轮廓上，局域化项的玻色部分半正定，这就把泛函积分限制到这一作用量部分为零的映射，即世界面瞬子。换句话说，拓扑场论中的瞬子鞍点局域化是精确的。在\BookEqRef{eq:3.117}{3.117}对应的局域化方程中，这些世界面瞬子，是满足 $\bar\partial\phi^i=0$ 的全纯映射。

\SourceParagraph{ch03-3.4.1.1-002}{180}
因此，作泛函积分时，所有映射 $\phi\colon\Sigma\to Y$ 的无穷维空间，被全纯映射的有限维模空间替代。超对称抵消了费米泛函积分的 Pfaffian，与 $\phi$ 积分中剩下的相应行列式，从而在全纯映射模空间上，得到自然的积分测度；存在阻碍或奇异性时，需用相应的虚拟积分类处理。细节见文献\BookCite{ref468}与 \BookBibText{ref242}{MS1}。

\SourceParagraph{ch03-3.4.1.1-003}{180}
三点函数的结果为
\begin{equation}\label{eq:3.121}
\langle W_aW_bW_c\rangle
=\int_Ya\wedge b\wedge c
+\sum_{\beta\in I}N_{abc}^\beta
\mathrm{e}^{2\pi\mathrm{i}\int_\beta(B+\mathrm{i}\omega)},
\end{equation}
其中，$I$ 是非零的有效瞬子类的集合，$N_{abc}^\beta$ 是由 $Y$ 中有理曲线的映射模空间的相交理论给出的数；这里包括多重覆盖的可能性\BookCite{ref085,ref024}。在 Calabi--Yau 的情形中，假定 $a$、$b$、$c$ 的实微分形式次数之和，为 $2d$，其中 $d=\dim_{\mathbb{C}}Y$；否则，三点函数为零。一点函数为
\begin{equation}\label{eq:3.122}
\langle W_a\rangle=\int_Ya.
\end{equation}
仅当 $a$ 是最高次形式时，它才可能非零。这两个关联函数，共同确定 $H^*(Y,\mathbb{C})$ 上的 Frobenius 代数结构，因而确定其他关联函数。

\SourceParagraph{ch03-3.4.1.1-004}{180}
在大体积极限中，可以忽略\BookEqRef{eq:3.121}{3.121}的瞬子修正，此时代数结构，就是微分形式的楔积。在有限体积下，形变后的环，称为 $Y$ 的“量子上同调环”。这里已经触及一个非常广阔的主题，没有篇幅继续展开。详细介绍，见文献\BookCite{ref101}及其所引文献。

\SourceParagraph{ch03-3.4.1.1-005}{180}
对于弦理论，最重要的情形，是 $Y$ 为 Calabi--Yau 三维流形。对于通常完整群恰为 $\operatorname{SU}(3)$ 的情形，量子上同调环的结构，完全由次数为2的形式的三点函数编码。从二维 TFT 的观点看，2形式具有特殊地位，因为它们给出保持鬼数的 A 模型形变。可以证明，这些形变，正是改变辛形式 $\omega$ 与 B 场 $B$；二者都是闭2形式。

\SourceParagraph{ch03-3.4.1.1-006}{180}
再次强调，算符代数的结构，只依赖于 $B+\mathrm{i}\omega$，而不依赖于 $Y$ 的复结构。事实上，文献\BookCite{ref468}还说明，不需要 $Y$ 具有可积的复结构，也不要求 Calabi--Yau 条件。$Y$ 可以是任意带有相容概复结构的辛流形\footnote{如果 $Y$ 不是 Calabi--Yau 流形，A 模型一般只保留 $\mathbb{Z}_2$ 分次，或者由反常允许的剩余有限分次，正如前面所解释的那样。}。此时，瞬子对应于伪全纯曲线。由于拓扑 A 模型，只知道未扭曲理论数据的一小部分，它能够适用于更广的靶空间类别，也就并不令人意外。

\subsection{B 模型}\label{sec:3.4.2}

\SourceParagraph{ch03-3.4.2-001}{181}
为了得到 B 模型，采用不同的重新赋值。令 $\psi_\pm^{\bar\jmath}$ 为 $\phi^*\overline{T_X}$ 的截面，而 $\psi_+^j$ 为 $K_\Sigma\otimes\phi^*T_X$ 的截面，$\psi_-^j$ 为 $\overline{K_\Sigma}\otimes\phi^*T_X$ 的截面。定义世界面标量
\begin{equation}\label{eq:3.123}
\begin{aligned}
\eta^{\bar\jmath}&=\psi_+^{\bar\jmath}+\psi_-^{\bar\jmath},\\
\theta_j&=g_{j\bar k}
\left(\psi_+^{\bar k}-\psi_-^{\bar k}\right),
\end{aligned}
\end{equation}
并定义1形式 $\rho^j$，其 $(1,0)$ 形式部分为 $\psi_+^j$，$(0,1)$ 形式部分为 $\psi_-^j$。

\SourceParagraph{ch03-3.4.2-002}{181}
现在，考虑在原来的超对称变换中，取 $\alpha_\pm=0$、$\widetilde\alpha_\pm=\alpha$ 所得到的变换。与 A 模型一样，这给出 BRST 荷 $Q$，在运动方程成立时满足 $Q^2=0$。

\SourceParagraph{ch03-3.4.2-003}{181}
此时，可以将作用量写为
\begin{equation}\label{eq:3.124}
S=\mathrm{i}\int_\Sigma\{Q,D\}+U,
\end{equation}
其中
\begin{equation}\label{eq:3.125}
\begin{aligned}
D&=g_{j\bar k}
\left(\rho_z^j\bar\partial\phi^{\bar k}
+\rho_{\bar z}^j\partial\phi^{\bar k}\right),\\
U&=\int_\Sigma\left(
-\theta_jD\rho^j
-2\mathrm{i}R_{j\bar l k\bar n}\rho^j\wedge\rho^k
\eta^{\bar l}\theta_mg^{m\bar n}\right).
\end{aligned}
\end{equation}

\SourceParagraph{ch03-3.4.2-004}{181}
B 模型还有一个额外的复杂之处：费米子的扭曲方式，比 A 模型更不对称。对于一般靶空间 $X$，费米泛函积分的 Pfaffian 的相位定义存在歧义，对应于一个手征反常。如果要求 $c_1(T_X)=0$，即这里的 Calabi--Yau 条件，这一反常消失\footnote{更精确地，当 $\Sigma$ 的边界为空时，消去所述反常，只需要求 $2c_1(T_X)=0$。不过，若像后面那样，考虑 $\Sigma$ 具有非空边界的情形，就需要更强的条件 $c_1(T_X)=0$。}。

\SourceParagraph{ch03-3.4.2-005}{181}
由\BookEqRef{eq:3.125}{3.125}可以明显看出，改变 $\Sigma$ 的度规，只通过 $Q$ 恰当项改变作用量。另一个不那么明显、但确实成立的事实，是改变 K\"ahler 形式 $\omega$，以及 B 场的 $(1,1)$ 分量，也只通过 BRST 恰当项改变作用量。因此，B 模型的关联函数，与这些参数无关。在 SCFT 语言中，这是因为，这些形变算符在 $Q_B$ 上同调中平凡。也可以在 sigma 模型中加入辅助场，得到一个等价表述，使这些变化是 $Q$ 恰当的\BookCite{ref318}。

\SourceParagraph{ch03-3.4.2-006}{182}
现在，将局部可观测量写为
\begin{equation}\label{eq:3.126}
W[A]=\eta^{\bar k_1}\cdots\eta^{\bar k_q}
A_{\bar k_1\cdots\bar k_q}^{j_1\cdots j_p}
\theta_{j_1}\cdots\theta_{j_p},
\end{equation}
其中
\begin{equation}\label{eq:3.127}
A=\mathrm{d}\bar z^{\bar k_1}\wedge\cdots\wedge\mathrm{d}\bar z^{\bar k_q}
A_{\bar k_1\cdots\bar k_q}^{j_1\cdots j_p}
\frac{\partial}{\partial z^{j_1}}\wedge\cdots\wedge
\frac{\partial}{\partial z^{j_p}},
\end{equation}
是 $X$ 上取值于 $\bigwedge^pT_X$ 的 $(0,q)$ 形式。也可以将 $A$ 称为“$(-p,q)$ 形式”。注意，通过与全纯 $d$ 形式 $\Omega$ 收缩，可以在 $(-p,q)$ 形式空间，与 $(d-p,q)$ 形式空间之间，得到一个同构。讨论镜像对称时，经常隐含或显式地使用这个同构，正如\BookSectionRef{sec:3.4.3}{3.4.3}中将看到的那样。

\SourceParagraph{ch03-3.4.2-007}{182}
现在，有
\begin{equation}\label{eq:3.128}
\{Q,W[A]\}=-W[\bar\partial A].
\end{equation}
因此，对于 B 模型，$Q$ 上同调，是取值于全纯切丛各阶外幂的形式的 Dolbeault 上同调。

\SourceParagraph{ch03-3.4.2-008}{182}
B 模型中的瞬子平凡。\BookEqRef{eq:3.125}{3.125}所给出的局域化条件，要求 $\bar\partial\phi^{\bar k}=\partial\phi^{\bar k}=0$，也就是说，$\phi$ 是常值映射，将 $\Sigma$ 映到 $X$ 中的一个点。事实上，这里的拓扑关联函数，不受到量子修正。这可以从超对称微扰理论的结构看出，也可以用以下论证理解：$Q$ 恰当项的系数，可以任意缩放而不改变关联函数。利用它将泛函积分局域化到常值映射，并结合超对称对非零模的抵消，结果就由经典的零模积分给出，相当于相应的 $\hbar\to0$ 极限。

\SourceParagraph{ch03-3.4.2-009}{182}
可以证明，球面上的 B 模型关联函数为
\[
\langle W_AW_B\cdots\rangle
=\int_X\Omega\wedge\iota_{AB\cdots}\Omega.
\]
这里假定，$A$、$B$ 等的微分形式次数之和，以及多重向量次数之和，都等于 $d$。于是，$AB\cdots$ 是一个取值于 $\bigwedge^dT_X$ 的 $(0,d)$ 形式；它与 $(d,0)$ 形式 $\Omega$ 收缩，得到普通的 $(0,d)$ 形式。

\SourceParagraph{ch03-3.4.2-010}{182}
在最有意思的情形中，$X$ 是 Calabi--Yau 三维流形，$A$、$B$、$C$ 是取值于 $T_X$ 的 $(0,1)$ 形式，此时，可以写得更显式。三点函数为
\begin{equation}\label{eq:3.129}
\langle W_AW_BW_C\rangle
=\int_X\Omega_{jkl}A^j\wedge B^k\wedge C^l\wedge\Omega,
\end{equation}
其中，$A=A^j\partial/\partial\phi^j\in H^1_{\bar\partial}(X,T_X)$，其余类似；$\Omega_{jkl}$ 是全纯3形式的分量。也可以借助 K\"ahler 度规升降指标，将相关张量与反全纯3形式 $\bar\Omega$ 联系起来；虽然这种表示表面上引入了 K\"ahler 度规依赖，关联函数中的这一依赖，在上同调中平凡。B 模型中这些 $(-1,1)$ 形式所对应的形变，就是 $X$ 的复结构形变；回忆，复结构的无穷小形变，由 $H^1(X,T_X)$ 中的元素分类。

\SourceParagraph{ch03-3.4.2-011}{183}
如果上同调群 $H^0(X,\bigwedge^2T_X)$ 或 $H^2(X,\mathcal{O})$ 非零，那么，还有一些不那么直观的 B 模型形变，分别与 $(-2,0)$ 形式及 $(0,2)$ 形式相关。在 完整群恰为 $\operatorname{SU}(3)$ 的流形上，没有这些形变；不过，当 $X$ 是复环面、K3 曲面，或者它们的乘积时，它们存在。可以使用 Hitchin 引入的广义复结构概念\BookCite{ref236}，解释其几何意义，见\BookSectionRef{sec:6.2.5}{6.2.5}。

\SourceParagraph{ch03-3.4.2-012}{183}
注意，B 模型确实要求 $X$ 具有复结构，并满足相应的 Calabi--Yau 条件。不过，它不要求指定 K\"ahler 形式 $\omega$。这意味着，B 模型只依赖于 $X$ 的复结构。与 A 模型类似，可以猜想，对不一定具有 K\"ahler 形式、但满足 $c_1(T_X)=0$ 的复流形，也能够定义相应的 B 模型。这确实成立\BookCite{ref293}。

\subsection{闭弦镜像对称}\label{sec:3.4.3}

\SourceParagraph{ch03-3.4.3-intro-001}{183}
这里只对后面需要的内容，作极简要的介绍；其余参见 \BookBibText{ref242}{MS1} 与文献\BookCite{ref101}。镜像对称有多种强弱不同的定义。这里仅需要一个较弱的定义。
\begin{definition*}\label{def:3.closed-string-mirror}
如果以 $Y$ 为靶空间的 A 模型的算符代数，与以 $X$ 为靶空间的 B 模型的算符代数同构，就称两个 Calabi--Yau 流形 $X$ 与 $Y$ 互为镜像。
\end{definition*}

\SourceParagraph{ch03-3.4.3-intro-002}{183}
原来的定义更强，是一个有关共形场论的陈述。从弦理论的观点看，最有意思的情形，是 $X$ 与 $Y$ 为 Calabi--Yau 三维流形。此时，最强的定义要求：将 IIA 型弦理论紧化在 $Y$ 上，与将 IIB 型弦理论紧化在 $X$ 上，在四维中给出等价的物理。

\SourceParagraph{ch03-3.4.3-intro-003}{183}
简单分析算符代数各个向量空间的维数，就得到
\[
h^{p,q}(Y)=h^{d-p,q}(X),
\]
因而
\[
\chi(Y)=(-1)^d\chi(X).
\]

\subsubsection{镜像映射}\label{sec:3.4.3.1}

\SourceParagraph{ch03-3.4.3.1-001}{183}
$Y$ 上 A 模型的算符代数，依赖于 $Y$ 上 $B+\mathrm{i}\omega$ 的选择，而 $X$ 上 B 模型的算符代数，依赖于 $X$ 的复结构选择。因此，镜像对称的精确陈述，必须将 $Y$ 的 $B+\mathrm{i}\omega$ 模空间，与 $X$ 的复结构模空间相联系。这一映射，称为“镜像映射”。\BookSectionRef{sec:5.6.1}{5.6.1}将对一个最简单的例子作较详细的计算：以线丛 $\mathcal{O}_{\mathbb{P}^2}(-3)$ 的全空间为非紧 Calabi--Yau 流形的情形。

\SourceParagraph{ch03-3.4.3.1-002}{183}
这里按照文献\BookCite{ref193,ref085}，介绍研究最多的一对镜像 Calabi--Yau 三维流形。取 $Y$ 为“五次三维流形”，即 $\mathbb{P}^4$ 中，由一个五次齐次多项式的零点定义的超曲面。由于 $h^{1,1}(Y)=1$，复化 K\"ahler 类的模空间是一维的。将 $H^2(Y,\mathbb{Z})$ 的正生成元记为 $e$\footnote{也就是说，$\int_Ye^3>0$。}。略微混用记号，将复化 K\"ahler 形式的上同调类，写为 $(B+\mathrm{i}\omega)e$；在五次超曲面这一语境中，$B$ 与 $\omega$ 是实数。因此，可以把 $Y$ 的大小，即 $\omega$，看作由周围 $\mathbb{P}^4$ 的大小决定。$Y$ 还有 $h^{2,1}=101$，因而具有101个复结构形变，不过，这里不关注它们。

\SourceParagraph{ch03-3.4.3.1-003}{184}
可以通过取 $Y$ 的 $(\mathbb{Z}_5)^3$ 轨形商，构造其镜像 $X$。所得空间具有轨形奇点；要得到光滑的 Calabi--Yau 流形，需要消解这些奇点，从而产生 $B+\mathrm{i}\omega$ 的100个新自由度。不过，这里只关心 $X$ 的复结构，它由五次多项式的选择决定。与 $(\mathbb{Z}_5)^3$ 轨形作用相容的最一般的五次多项式，在适当坐标中为
\begin{equation}\label{eq:3.130}
x_0^5+x_1^5+x_2^5+x_3^5+x_4^5
-5\psi x_0x_1x_2x_3x_4.
\end{equation}
因此，复结构由单个复参数 $\psi$ 决定。镜像映射，就成为 A 模型一侧的 $B+\mathrm{i}\omega$，与 B 模型一侧的 $\psi$ 之间的映射。

\SourceParagraph{ch03-3.4.3.1-004}{184}
镜像映射相当复杂，实际上是一个多对多的映射。因此，通常先选一个基点，一般取 A 模型一侧的大半径极限，再在这一基点的某个邻域中，求出镜像映射。随后，可以尝试将它解析延拓到更大的区域。

\SourceParagraph{ch03-3.4.3.1-005}{184}
也可以不涉及某个具体紧化，而通过研究四维 $N=2$ 超引力理论中标量场的一般性质，内禀地分析模空间。结果是，模空间是所谓“特殊 K\"ahler 流形”\BookCite{ref108,ref431,ref085}。关于这一主题的良好数学表述，见文献\BookCite{ref154}。

\SourceParagraph{ch03-3.4.3.1-006}{184}
模空间的特殊 K\"ahler 结构，给出优选的、但不唯一的坐标，即“特殊坐标”；它们满足某些平坦性约束。在 A 模型一侧，$B+\mathrm{i}\omega$ 的分量是特殊坐标。不过，在 B 模型一侧，像\BookEqRef{eq:3.130}{3.130}中的 $\psi$ 这样的复参数，并不是特殊坐标。

\SourceParagraph{ch03-3.4.3.1-007}{184}
刻画复结构的一种更自然的方法，是使用 Hodge 结构，更具体地，使用全纯3形式 $\Omega$ 的射影上同调类，位于 $\mathbb{P}(H^3(X,\mathbb{C}))$。对于镜像五次三维流形，在局部选定标记后，全纯3形式的类可以刻画复结构。不过，$\mathbb{P}(H^3(X,\mathbb{C}))$ 的维数太大，其中并非所有点，都对应于 Calabi--Yau 三维流形的复结构。局部地，参数化 Calabi--Yau 三维流形的周期锥，是 $H^3(X,\mathbb{C})$ 中的复 Lagrange 子流形，其射影化位于上述射影空间，并由全纯预势决定。为了定义它，可以选择 $H_3(X,\mathbb{Z})$ 自由部分的辛基：$\alpha_m$、$\beta^m$，其中 $m=0,\ldots,h^{2,1}(X)$，具有相交数
\begin{equation}\label{eq:3.131}
\alpha_m\cap\alpha_n=0,\qquad
\alpha_m\cap\beta^n=\delta_m^n,\qquad
\beta^m\cap\beta^n=0.
\end{equation}

\SourceParagraph{ch03-3.4.3.1-008}{185}
A 循环的周期
\begin{equation}\label{eq:3.132}
\varpi_m=\int_{\alpha_m}\Omega,
\end{equation}
构成一组齐次特殊坐标。还可以证明，B 循环的周期，在局部由一个全纯函数 $F$ 确定：
\begin{equation}\label{eq:3.133}
\frac{\partial F}{\partial\varpi_m}=\int_{\beta^m}\Omega.
\end{equation}

\SourceParagraph{ch03-3.4.3.1-009}{185}
A 模型中也有相同的结构：此时将复化 K\"ahler 模空间，作为复 Lagrange 子流形，嵌入
\[
H^2(Y,\mathbb{C})\subset
\mathbb{P}\!\left(\bigoplus_{k=0,2,4,6}H^k(Y,\mathbb{C})\right).
\]
特殊几何则要求，镜像映射是这两个环绕空间之间的射影线性辛映射。

\SourceParagraph{ch03-3.4.3.1-010}{185}
由于镜像映射是线性的，它主要由匹配大半径极限点附近、路径 $B\mapsto B+1$ 诱导的单值变换确定。文献\BookCite{ref366}分析了一套系统方法。要确定最后的常数，需要显式计数 $Y$ 上的一些有理曲线，并将 A 模型关联函数，直接对应到 B 模型关联函数。尽管如此，镜像映射还有一种猜想形式，它在文献\BookCite{ref085}中已被隐含使用，在所讨论的已知例子中，似乎都成立。更多细节，见文献\BookCite{ref101}。

\subsubsection{镜像对称的 SCFT 论证}\label{sec:3.4.3.2}

\SourceParagraph{ch03-3.4.3.2-001}{185}
提示镜像对称的最简单观察，基于 $N=2$ SCA，以及\BookSectionRef{sec:3.2.4}{3.2.4}讨论的 $U(1)$ 子代数因子分解。因子分解意味着，$U(1)$（自由玻色子）理论的任意自同构，都可以提升为整个 $N=2$ SCFT 的自同构。

\SourceParagraph{ch03-3.4.3.2-002}{185}
现在，在这一语境中考虑 T 对偶变换。根据\BookEqRef{eq:3.69}{3.69}，它的作用为
\begin{equation}\label{eq:3.134}
J\longmapsto J,\qquad\bar J\longmapsto-\bar J.
\end{equation}
\BookEqRef{eq:3.91}{3.91}中相应的代数自同构，交换
\[
G_-\longleftrightarrow\widetilde G_-,
\]
保持其他生成元不变。于是，由\BookEqRef{eq:3.100}{3.100}可知，它交换 A 模型与 B 模型。

\SourceParagraph{ch03-3.4.3.2-003}{185}
因此，给定一个以 $Y$ 为靶空间的 A 模型，可以通过镜像的 $N=2$ 代数识别，得到一个 B 模型与之同构的 SCFT；反之亦然。当然，这还不能说明，必然存在一个给出该 B 模型的 Calabi--Yau 流形 $X$。事实上，反向的几何存在性陈述，有明显的反例。例如，如果从刚性 Calabi--Yau 流形 $X$ 的 B 模型出发，即 $h^{2,1}(X)=0$，几何镜像所要求的 $h^{1,1}=0$，不能由紧 K\"ahler 流形实现。

\SourceParagraph{ch03-3.4.3.2-004}{185}
不过，在显式构造中，可以自然地猜测，这一 T 对偶代数变换，如何联系 $X$ 与 $Y$。最初的例子，是\BookSectionRef{sec:3.3.6}{3.3.6}中讨论的五次 Calabi--Yau 的 Gepner 模型构造，即五个层数为3的 A 系列最小模型乘积的 $\mathbb{Z}_5$ 商。可以显式地对每个最小模型施加 SCFT 变换，从而构造镜像 SCFT。分析后发现，它等同于最小模型乘积的 $(\mathbb{Z}_5)^4$ 商，因而与前面提到的几何 $(\mathbb{Z}_5)^3$ 商相对应。细节见文献\BookCite{ref192}。

\SourceParagraph{ch03-3.4.3.2-005}{186}
后来，人们从更几何的角度理解了这一推理，并在环面簇超曲面这一类中，得到镜像对的显式构造\BookCite{ref032}。这一类很大，但远未穷尽 Calabi--Yau 三维流形，这也是寻找其他镜像对称图景的动机之一。不过，这一论证的重要优点，是它在 SCFT 中表述；因此，一旦承认 $X$ 与 $Y$ 在这一意义上互为镜像，就可以得出精确的结论，其中隐含地计入了弦理论与通常几何之间的一切差异。

\section{边界 CFT}\label{sec:3.5}

\SourceParagraph{ch03-3.5-intro-001}{186}
本节考虑具有边界的世界面 $\Sigma$。细致分析它，需要相当多的技术细节，本书并不需要全部展开。更多细节，见文献\BookCite{ref241}，其中许多分析也出现在 \BookBibText{ref242}{MS1} 中；更全面的处理，见文献\BookCite{ref006,ref336}。

\subsection{一般讨论}\label{sec:3.5.1}

\SourceParagraph{ch03-3.5.1-001}{186}
目标是定义一个开闭弦范畴，满足\BookChapterRef{2}给出的 TFT 公理。最简单的讨论方式，是从一个闭弦 TFT，即 Frobenius 代数 $\mathcal{C}$ 出发，以它来定义开弦结构；也许可以类比于复流形上的层的构造。虽然\BookSectionRef{sec:2.5}{2.5}介绍了一些方法，尝试加入底层复流形 $M$ 的更多几何信息，但到目前为止，\BookChapterRef{2}的直接方法，只在简单的代数 $\mathcal{C}$ 上有效，即对应于零维的 $M$。

\SourceParagraph{ch03-3.5.1-002}{186}
Calabi--Yau 上开弦 TFT 的原始讨论，来自 Witten\BookCite{ref463}。它从\BookSectionRef{sec:3.4}{3.4}给出的 A、B 拓扑扭曲 sigma 模型定义出发，识别与扭曲相容的几何边界条件。下面，将这些边界条件称为“拓扑 Dirichlet 膜”。

\SourceParagraph{ch03-3.5.1-003}{186}
很快将说明，在这里讨论的几何类别中，A 模型的拓扑 D 膜，是带有配备平坦联络的丛的 Lagrange 子流形；B 模型中，则是带有全纯向量丛的全纯子流形。还可以看出世界面圆盘瞬子修正的来源，它们给出 Fukaya 范畴的 $A_\infty$ 结构。

\SourceParagraph{ch03-3.5.1-004}{186}
虽然这给出了问题的基本结构，但要从物理上理解镜像对称，还需要 CFT 的更多信息。因此，接下来讨论 D 膜物理理论的一部分，特别是解释 T 对偶的作用。

\subsection{开弦 CFT 的一般性质}\label{sec:3.5.2}

\SourceParagraph{ch03-3.5.2-intro-001}{187}
开始以前，按照文献\BookCite{ref086}，在\BookSectionRef{sec:3.2}{3.2}的层面上，讨论这一问题的一般性质。

\SourceParagraph{ch03-3.5.2-intro-002}{187}
虽然前面的讨论，已在定义 Hilbert 空间 $\mathcal{H}$ 与 QFT 函子时涉及边界条件，但现在要定义第二种本质上不同的边界。在\BookChapterRef{2}的语言中，Hilbert 空间 $\mathcal{H}$ 是 $\mathcal{C}$ 的 CFT 类似物，而与 D 膜对应的新边界，来自“简单”边界条件的集合 $B_0$。

\SourceParagraph{ch03-3.5.2-intro-003}{187}
CFT 边界条件的主要物理约束，是保持世界面上的能量与动量守恒。这是对\BookEqRef{eq:3.27}{3.27}定义的能动张量 $T$ 的局部条件：
\begin{equation}\label{eq:3.135}
0=t^\mu n^\nu T_{\mu\nu}\big|_{\partial\Sigma},
\end{equation}
其中，$t^\mu$ 与 $n^\nu$，分别是边界的切向量与法向量。

\SourceParagraph{ch03-3.5.2-intro-004}{187}
设考虑 $\partial\Sigma$ 的一个分支 $S\cong S^1$。与\BookSectionRef{sec:3.2.2.1}{3.2.2.1}一样，可以把边界条件，看作闭弦 Hilbert 空间的一个广义“态” $|B\rangle$，其上有两个彼此对易的 Virasoro 代数 $L_n$ 与 $\widetilde L_n$ 的作用。将能动张量按这些模式展开，条件\BookEqRef{eq:3.135}{3.135}成为
\begin{equation}\label{eq:3.136}
(L_n-\widetilde L_{-n})|B\rangle=0.
\end{equation}
可以证明\BookCite{ref255}，该方程的解与初级场一一对应。最简单的论证\BookCite{ref036}，利用 Hilbert 空间的表示分解，以及每个 $\mathcal{H}_{R,i}$ 上，使 $L_n=L_{-n}^\dagger$ 成立的内积的唯一性；归一化之外，这个内积已确定。于是，可以将 $|B\rangle$，重新解释为表示空间之间的算符，或者相应的广义泛函，
\begin{equation}\label{eq:3.137}
O_B\in\prod_i
\bigl(\mathcal{H}_{L,i}\otimes\mathcal{H}_{R,i}\bigr)^*,
\end{equation}
满足交织关系
\[
L_nO_B=O_BL_n.
\]
由 Schur 引理，在不可约表示没有重数的情形中，这样的算符 $O_B$，是各个不可约 Virasoro 表示投影的线性组合。
\begin{definition*}\label{def:3.ishibashi-state}
与单个这样的投影对应的边界态，称为 Ishibashi 态。
\end{definition*}
这些态不可正规化，因而不是 $\mathcal{H}$ 中的普通向量\footnote{物理上需要配对 $\langle B|\phi\rangle$ 有意义，因此，应把它们看作定义在适当稠密态空间上的广义线性泛函，而非 Hilbert 空间中的可正规化态。}。物理中构造的标准边界条件，如有理模型中的“Cardy 态”，以及更一般的“边界态”，是 Ishibashi 态的适当线性组合。

\SourceParagraph{ch03-3.5.2-intro-005}{187}
与闭弦讨论一样，后面的讨论，分为两类。对于 $c\le1$ 的情形，可以利用 Virasoro 表示论，完整分类边界条件。“Cardy 条件”起着重要作用；它们是\BookChapterRef{2}公理的一部分，表示第2章\BookRef{fig:2.12}{图12}所描绘的世界面，必须具有相容的开弦与闭弦解释。结果，是一个有限的集合 $B_0$；在最简单的对角模型中，其基数等于允许的不可约表示数。$c=1$ 的一般情形，还需要更具体的结构。

\SourceParagraph{ch03-3.5.2-intro-006}{188}
对于 $c>1$，没有同样的一般理论。事实上，甚至不知道\BookEqRef{eq:3.136}{3.136}是否足以定义这一问题，也不知道是否还需要额外的约束。不过，在前面讨论的三种框架中，都可以得到部分结果。Gepner 模型的边界条件，见文献\BookCite{ref402,ref069}；线性 sigma 模型的情形，见文献\BookCite{ref241}与 \BookBibText{ref242}{MS1}。

\subsubsection{sigma 模型中的 Dirichlet 膜}\label{sec:3.5.2.1}

\SourceParagraph{ch03-3.5.2.1-001}{188}
在非线性 sigma 模型中，自然地，可以通过几何定义边界条件，下面就这样做。回忆 sigma 模型的经典运动方程\BookEqRef{eq:3.76}{3.76}；当 $B=0$ 时，它定义调和映射 $\phi\colon\Sigma\to M$。因此，寻找调和映射方程的自然边界条件。至少在开始时，这些条件也应是局域的，即在每一点上，对场及其有限阶导数施加条件，以描述不带额外物理自由度的 D 膜。

\SourceParagraph{ch03-3.5.2.1-002}{188}
从数学上，可以要求边界条件保持椭圆性。在最简单的靶空间 $\mathbb{R}$ 的情形中，需要为标量 Laplace 算符选择椭圆边界条件。两种基本的局部选择是：
\begin{definition*}\label{def:3.dirichlet-neumann-boundary}
\begin{itemize}
\item Dirichlet 边界条件：$\phi|_{\partial\Sigma}=\phi_0$，其中 $\phi_0\in\mathbb{R}$。
\item Neumann 边界条件：$n^\mu\partial_\mu\phi|_{\partial\Sigma}=0$。
\end{itemize}
\end{definition*}
也可以将 Dirichlet 条件写成 $t^\mu\partial_\mu\phi|_{\partial\Sigma}=0$，使两种情形表面上更加对称。不过，Dirichlet 边界条件实际上依赖于参数 $\phi_0$；要得到恰当的对称类比，就需要为 Neumann 情形，也找到一个对应的参数，见\BookSectionRef{sec:3.5.2.2}{3.5.2.2}。

\SourceParagraph{ch03-3.5.2.1-003}{188}
也可以将这些边界条件，重新表述为左右行部分之间的关系。取 Lorentz 世界面上的切、法方向，并略去共同的流归一化，注意
\begin{equation}\label{eq:3.138}
J_L=(t^\mu+n^\mu)\partial_\mu\phi,
\qquad J_R=(t^\mu-n^\mu)\partial_\mu\phi,
\end{equation}
于是
\begin{equation}\label{eq:3.139}
\begin{aligned}
J_L&=-J_R,&&\text{Dirichlet 条件},\\
J_L&=J_R,&&\text{Neumann 条件}.
\end{aligned}
\end{equation}

\SourceParagraph{ch03-3.5.2.1-004}{188}
推广到靶空间 $\mathbb{R}^d$ 时，除了上面的直接推广，还可以有“混合”边界条件：边界映射 $\phi|_{\partial\Sigma}$，必须取值于一个固定的子流形 $L\subset M$。可以通过选择由 $L$ 的切向与法向坐标组成的坐标系，说明这些边界条件的椭圆性。在这些坐标中，调和映射方程的主导项仍是 Laplace 算符，混合边界条件，分别约化为 Neumann 与 Dirichlet 条件。

\SourceParagraph{ch03-3.5.2.1-005}{189}
\begin{definition*}\label{def:3.dirichlet-p-brane}
如果一个边界条件，具有 $p$ 个 Neumann 分量与 $d-p$ 个 Dirichlet 分量，从而 $\dim_{\mathbb{R}}L=p$，就称它为 Dirichlet $p$ 膜。
\end{definition*}
例如，0膜与一点相关，1膜与实直线或弦相关，其余类似。

\SourceParagraph{ch03-3.5.2.1-006}{189}
为运动方程 $\delta S=0$ 找到适当边界条件的一个更系统的方法，是要求作用量变分中的边界项消失：
\[
\delta S\big|_{\partial\Sigma}=0.
\]
这就消去了通常的 Euler--Lagrange 推导中，分部积分产生的边界项。因此，能量与动量守恒的标准物理论证，仍然成立，满足主要条件\BookEqRef{eq:3.135}{3.135}。此外，还必须要求，边界条件与运动方程，不会共同导致一个过定的初值问题。

\SourceParagraph{ch03-3.5.2.1-007}{189}
下面，以无限“条带” $\Sigma$ 上、靶空间为 $M$ 的非线性 sigma 模型为例说明。$\Sigma$ 带有 Euclidean 度规，以 $\sigma_1\in\mathbb{R}$、$\sigma_2\in[0,1]$ 参数化，如\BookRef{fig:3.4}{图4}。施加边界条件的 $\Sigma$ 的边界，由 $\sigma_2=0$ 与 $\sigma_2=1$ 两个分支组成。这里考虑有定向的开弦，因此，区别弦的“起点” $\sigma_2=0$ 与“终点” $\sigma_2=1$。
\begin{figure}[H]
\centering
\BookImage{01_Numbered_Figures/Chapter_03/Ch03_Fig04_PDFp189_Bookp179.pdf}
\caption{开弦世界面。}\label{fig:3.4}
\end{figure}

\SourceParagraph{ch03-3.5.2.1-008}{189}
玻色 sigma 模型的作用量，采用这一边界讨论中的度规耦合归一化，写为
\[
S=\frac12\int_\Sigma\mathrm{d}^2\sigma\,
G_{IJ}(\phi)\sum_{i=1,2}\partial_i\phi^I\partial_i\phi^J,
\]
其中，$\phi$ 是从 $\Sigma$ 到 $M$ 的映射，$\partial_i=\partial/\partial\sigma_i$。对作用量变分，并作分部积分，得到必须置为零的边界项：
\[
\int_{\partial\Sigma}\mathrm{d}\sigma_1\,
G_{IJ}(\phi)\delta\phi^I\partial_n\phi^J.
\]
这里采用各分支的外法向导数；在固定条带坐标中，各分支相应带有符号。一个消去它的方法，是约束边界上的 $\phi$，取值于子流形 $L\subset M$。这意味着，切向导数 $\partial_1\phi^I$ 位于 $TL\subset TM$ 中。为了保证作用量变分的边界项消失，可以再要求，法向导数 $\partial_2\phi^I$ 位于法丛 $NL$ 中。所得边界条件，在 $L$ 的法方向为 Dirichlet 型，在切方向为 Neumann 型。注意，不能对同一个分量 $\phi^I$，同时施加 Dirichlet 与 Neumann 条件，因为与运动方程结合时，会得到过定的方程组。

\subsubsection{含规范场的边界条件}\label{sec:3.5.2.2}

\SourceParagraph{ch03-3.5.2.2-001}{190}
当 $L$ 的维数大于1时，可以推广这些边界条件。保留 $\partial_1\phi^I$ 的 $NL$ 分量为零的条件，但对 $\partial_2\phi^I$ 的 $TL$ 分量，施加更一般的条件：
\[
G_{IJ}(\phi)\partial_2\phi^J
=F_{IJ}(\phi)\partial_1\phi^J.
\]
为保持相容，这里的全部指标都限制在 $TL$ 子空间中。称它为形变后的 Neumann 边界条件。如果矩阵函数 $F_{IJ}$，对于某个1形式的分量 $A_I(\phi)$，具有 $\partial_IA_J-\partial_JA_I$ 的形式，那么，进一步分部积分后，作用量变分的边界部分，成为
\[
\delta\int_{\partial\Sigma}\mathrm{d}\sigma_1\,
A_I(\phi)\partial_1\phi^I.
\]
因此，如果向作用量加入边界项
\begin{equation}\label{eq:3.140}
S_b=-\int_{\partial\Sigma}\mathrm{d}\sigma_1\,
A_I(\phi)\partial_1\phi^I,
\end{equation}
总作用量 $S+S_b$ 的变分，就没有边界部分。

\SourceParagraph{ch03-3.5.2.2-002}{190}
为使 $A_I$ 与 $F_{IJ}$ 的几何意义更清楚，引入1形式 $A=A_I(\phi)\mathrm{d}\phi^I\in\Omega^1(L)$。于是，$F_{IJ}$ 是2形式 $F=\mathrm{d}A\in\Omega^2(L)$ 的分量，边界作用量可以写成显式的微分同胚不变形式：
\[
S_b=-\int_{\partial\Sigma}\phi^*A.
\]
Euclidean 表述中的幺正性，要求此处的 $A$ 为纯虚；因此，将它重新写为 $-\mathrm{i}$ 乘一个实1形式，仍将后者记为 $A$，得到
\[
S_b=\mathrm{i}\int_{\partial\Sigma}\phi^*A.
\]
这样重新定义以后，当 $A$ 是 $L$ 上的实1形式时，理论是幺正的。

\SourceParagraph{ch03-3.5.2.2-003}{190}
显然，对 $L$ 上的任意函数 $f$，以 $A'=A+\mathrm{d}f$ 替代 $A$，不会改变闭边界上的作用量。这提示，$A$ 不必是全局定义的1形式，而可以只是 $L$ 上一个主 $U(1)$ 丛的联络。为了说明这一点，取边界的一个连通分支 $\partial_0\Sigma$。可以假定它闭合，也就是与圆周微分同胚。如果边界映射 $\phi|_{\partial_0\Sigma}\colon S^1\to L$，可以延伸为圆盘上的映射 $\widetilde\phi\colon D^2\to L$，那么，边界作用量可以写成
\[
\mathrm{i}\int_{D^2}\widetilde\phi^*F.
\]
它只依赖于 $F$，而不显式依赖于 $A$。它确实依赖于从 $\phi$ 延伸到 $\widetilde\phi$ 的方式；不过，两个延伸之间的差为
\[
\mathrm{i}\int_{S^2}\widetilde\phi^*F=2\pi\mathrm{i}n,
\qquad n\in\mathbb{Z},
\]
其中，球面由两个延伸圆盘沿边界粘合得到。由于边界作用量，仅通过 $\exp(-S_b)$ 进入泛函积分，这一歧义并不重要。一般的、不可延伸的闭边界，则直接用联络沿闭曲线的平行移动定义这一因子。

\SourceParagraph{ch03-3.5.2.2-004}{191}
即使 $F_{IJ}=0$，也就是说，边界条件在 $L$ 的切方向为 Neumann 型，当 $L$ 不是单连通时，$U(1)$ 联络也未必与平凡联络规范等价。这在很大程度上弥补了\BookSectionRef{sec:3.5.2.1}{3.5.2.1}提到的 Dirichlet 与 Neumann 边界条件之间的不对称，因为 Neumann 条件现在也有一个可选参数，即平坦联络。对于 $T^d$ 上的膜，平坦联络由对偶环面参数化，两者之间便有完整的对称。很快将说明，这正是 T 对偶在 Dirichlet 膜上的作用方式。

\subsubsection{世界面的量子化}\label{sec:3.5.2.3}

\SourceParagraph{ch03-3.5.2.3-001}{191}
$L$ 的选择，或者更一般地，三元组 $(L,E,\nabla)$ 的选择，是定义边界条件或 Dirichlet 膜的基本数据；其中，$E$ 是 $L$ 上的 $U(1)$ 丛，$\nabla$ 是 $E$ 上的联络。

\SourceParagraph{ch03-3.5.2.3-002}{191}
从 QFT 的观点看，选择 $L$，类似于在原来的 sigma 模型中，设定一个一般度规。正如并非所有度规都给出 CFT，而只有 beta 函数为零的度规才可以，子流形 $L$ 也并非都给出相容的 CFT 边界条件。研究边界理论的重整化，可以导出一个边界 beta 函数，在这一语境中起着相同作用。

\SourceParagraph{ch03-3.5.2.3-003}{191}
在主导近似中，即所有长度尺度，都远大于弦长 $\ell_s$ 时，边界 beta 函数为零，对应于 $L$ 是体积驻值子流形。定义体积泛函，为拉回度规所对应的体积元的积分：
\begin{equation}\label{eq:3.141}
S=\operatorname{Vol}(L)
=\int_L\mathrm{d}^px\,\sqrt{\det\iota^*g},
\end{equation}
其中，$\iota\colon L\hookrightarrow M$ 是嵌入。于是，条件为
\[
\beta\sim\delta\operatorname{Vol}(L)=0.
\]
最简单的情形，是体积极小的子流形。用 CFT 的语言说，它们的特征是没有相关的边界算符。

\SourceParagraph{ch03-3.5.2.3-004}{191}
一般而言，寻找体积极小子流形很困难。不过，有些简单情形，例如 K\"ahler 流形中的全纯子流形。正如\BookChapterRef{1}中讨论的那样，这推广到校准子流形的概念；虽然它比全纯子流形更难研究，但至少是一个线性条件。

\SourceParagraph{ch03-3.5.2.3-005}{192}
这些问题，与相应的闭弦问题之间，有密切类比：寻找 Ricci 平坦流形的一般问题，包含寻找 K\"ahler--Einstein 度规这一较简单的特殊情形。这两种简化，都可以从物理上理解为超对称的结果；一个一般论证是，保持超对称的解，必须满足一阶方程，有时称为 BPS 方程。这里不再展开，参见文献\BookCite{ref382,ref172}与 \BookBibText{ref242}{MS1}。

\subsubsection{Born--Infeld 作用量}\label{sec:3.5.2.4}

\SourceParagraph{ch03-3.5.2.4-001}{192}
更一般地，如果希望纳入 $L$ 上的规范场 $A$，那么，$A$ 与 $L$ 都受到 RG 流的作用；共形不变性，要求相应的边界 beta 函数为零。将曲率 $F$ 按幂级数展开，在主导阶，这要求 $\mathrm{d}*F=0$，也就是说，$F$ 必须满足 Maxwell 方程。

\SourceParagraph{ch03-3.5.2.4-002}{192}
如果 $F$ 不小，Maxwell 方程会产生非线性修正。从物理上，曲率 $F$ 的自然量纲，是长度的负二次方。一个相关的无量纲量，是绕闭曲线的平行移动；对于小曲线，其相位在主导阶，正比于曲线所围的面积。因此，“小”意味着 $|F|\ll1/\ell_s^2$，其中，$\ell_s$ 是弦长。

\SourceParagraph{ch03-3.5.2.4-003}{192}
当 $F$ 不必小，但其导数较小，即 $\ell_s\|\partial F\|\ll\|F\|$ 时，可以用闭式表达式，确定共形不变性的条件。这个条件要求，Dirichlet 膜世界体上的场，即嵌入 $\iota\colon L\to M$ 与 $L$ 上的联络 $A$，满足 Born--Infeld 作用量导出的运动方程。将弦尺度因子吸收入本节的矩阵归一化后，作用量写为
\begin{equation}\label{eq:3.142}
S=\int_L\mathrm{d}^nx\,\sqrt{\det(g|_L+F)}.
\end{equation}
这里，$g|_L$ 是时空度规在 $L$ 上的拉回，看作对称的 $n\times n$ 矩阵；$F=\mathrm{d}A$ 是曲率2形式，按相同单位看作反对称的 $n\times n$ 矩阵；$\det$ 是二者相加所得矩阵的行列式。\BookSectionRef{sec:3.5.4}{3.5.4}将利用 T 对偶，导出这一作用量。

\subsubsection{开弦扇区}\label{sec:3.5.2.5}

\SourceParagraph{ch03-3.5.2.5-001}{192}
要定义一条开弦，需要两个边界条件，分别用于弦的两端。对具有边界条件 $a$ 与 $b$ 的场论，进行量子化，将得到开弦态的 Hilbert 空间 $\mathcal{H}_{ab}$。

\SourceParagraph{ch03-3.5.2.5-002}{192}
研究共形不变性对开弦的影响时，一个简单的世界面几何，是去掉原点的上半平面 $\operatorname{Im}z\ge0$、$z\ne0$，如\BookRef{fig:3.5}{图5}。这是一个具有两个边界分支的 Riemann 曲面，即半直线 $\operatorname{Re}z>0$ 与 $\operatorname{Re}z<0$。在前者选择边界条件 $a$，在后者选择边界条件 $b$。
\begin{figure}[H]
\centering
\BookImage{01_Numbered_Figures/Chapter_03/Ch03_Fig05_PDFp193_Bookp183.pdf}
\caption{表示为去掉原点的上半平面的开弦世界面。}\label{fig:3.5}
\end{figure}

\SourceParagraph{ch03-3.5.2.5-003}{193}
另一个世界面几何，是条带 $0\le x\le\pi$、$t\in\mathbb{R}$，它给出的图像，更接近\BookChapterRef{2}中的图像。二者通过共形变换 $z=\exp(t+\mathrm{i}x)$ 相联系。

\SourceParagraph{ch03-3.5.2.5-004}{193}
由于边界条件\BookEqRef{eq:3.135}{3.135}，右行与左行的 Virasoro 生成元，在边界上等同起来。一种方便的理解方式，是使用加倍方法：将反全纯依赖的那一组算符，特别是相应的能动张量，反射为下半平面中定义的全纯算符的复共轭。这样，就可以使用与闭弦相同的形式体系；不过，现在开弦 Hilbert 空间，只承载一个 Virasoro 代数的作用。

\SourceParagraph{ch03-3.5.2.5-005}{193}
在保持相应超共形结构的理论中，这推广为单个 $N=2$ SCA。例如，闭弦态由 $(J_L,J_R)$ 的本征值双分次，而开弦只具有单个荷分次；在相应的整荷归一化下，它是 $\mathbb{Z}$ 分次。
\begin{figure}[H]
\centering
\BookImage{01_Numbered_Figures/Chapter_03/Ch03_Fig06_PDFp193_Bookp183.pdf}
\caption{开弦的接合。}\label{fig:3.6}
\end{figure}

\SourceParagraph{ch03-3.5.2.5-006}{194}
由于开弦 CFT 中的相互作用，与\BookChapterRef{2}描绘的相同世界面相关，它们满足许多相同的形式性质。特别地，自然地，可以将边界条件，识别为一个范畴中的对象，将开弦 Hilbert 空间 $\mathcal{H}_{ab}$ 的元素，识别为态射。当然，进入 CFT 后，有限的、位置无关的结合乘法，将被像闭弦讨论中那样的 CFT 算符乘积展开替代，其结合性由相应的缝合关系表达。

\subsubsection{形变与世界体规范场}\label{sec:3.5.2.6}

\SourceParagraph{ch03-3.5.2.6-001}{194}
一旦有了定义良好的共形边界条件，就可以考虑它的形变理论。与\BookSectionRef{sec:3.2.5}{3.2.5}一样，可以通过向作用量加入新项来定义；不过，现在这些项，应限制在边界上：
\begin{equation}\label{eq:3.143}
Z[g]=\int[\mathrm{D}\phi]\,
\exp\left(-S_0[\phi]+\sum_i g^i\int_{\partial\Sigma}\mathrm{d}t\,O_i(t)\right).
\end{equation}
标度不变性现在要求，$O_i$ 是权重 $h=1$ 的算符。

\SourceParagraph{ch03-3.5.2.6-002}{194}
在非线性 sigma 模型中，处于 $\ell_s$ 的主导阶时，自然的 $h=1$ 边界算符，是将玻色场的导数 $\partial_i\phi$，乘以函数 $f(\phi)$ 得到的。虽然这个函数可以一般地定义，但只有它在 $L$ 上的限制，在 RG 中保留下来。此外，计算主导的 beta 函数，会发现，$f$ 必须在适当意义下调和。因此，在物理上，将 $f$ 识别为 $L$ 上的一个场；在经典极限中，它满足由 beta 函数为零给出的运动方程。

\SourceParagraph{ch03-3.5.2.6-003}{194}
最简单的例子，出现在 Neumann 边界条件下，此时切向导数 $\partial_t\phi$ 不为零。形变为
\begin{equation}\label{eq:3.144}
O(t)=A_I(\phi(t))\partial_t\phi^I(t),
\end{equation}
其中，$A$ 是 $L$ 上的1形式。正如\BookSectionRef{sec:3.5.2.2}{3.5.2.2}中所述，这一形变，对应于改变 $L$ 上的规范场 $A$。

\SourceParagraph{ch03-3.5.2.6-004}{194}
另一方面，对于具有 Dirichlet 边界条件的坐标，法向导数 $\partial_n\phi$ 不为零，因此，对于一般函数 $v(\phi)$，可以写为
\begin{equation}\label{eq:3.145}
O(t)=v_I(\phi(t))\partial_n\phi^I(t).
\end{equation}
推广到多个 Dirichlet 玻色场，参数 $v(\phi)$，取值于 $L$ 的余法丛 $N^*L$。

\SourceParagraph{ch03-3.5.2.6-005}{194}
由于有度规，$N^*L$ 与 $NL$ 同构，因此，自然预期，这些算符对应于子流形 $L$ 的形变。可以通过物理论证说明这一点\footnote{在闭弦通道中，将边界法方向的坐标 $x$ 取为 Euclidean 时间时，这一识别最为清楚：相应归一化下，$\int\mathrm{d}t\,\partial_x\phi$ 是动量算符，因而生成法向平移。}。这是 CFT“包含其自身形变理论”的又一个例子；由此可知，相应的模参数，已在最初选择嵌入时计入。

\subsubsection{对 B 场的依赖}\label{sec:3.5.2.7}

\SourceParagraph{ch03-3.5.2.7-001}{195}
还需要另一项推广。考虑原来 sigma 模型作用量\BookEqRef{eq:3.75}{3.75}中的 B 场项。为将本节的规范变换写得统一，将整体归一化吸收入实2形式 $B$，使这一项写为
\begin{equation}\label{eq:3.146}
S_B=\mathrm{i}\int_\Sigma\phi^*B.
\end{equation}
这里，B 是 $M$ 上的2形式。这一作用量具有一个与 $A$ 的规范对称类似、但不同的规范对称。在
\begin{equation}\label{eq:3.147}
\delta B=\mathrm{d}\Lambda,\qquad\Lambda\in\Omega^1(M)
\end{equation}
下，有
\begin{equation}\label{eq:3.148}
\delta S_B=\mathrm{i}\int_\Sigma\mathrm{d}(\phi^*\Lambda).
\end{equation}
在闭世界面上，这是全导数，积分为零，因此，变换\BookEqRef{eq:3.147}{3.147}是理论的局部对称。于是，B 是一个“2形式规范联络”；可以用 gerbe 的语言，使这一说法精确化\BookCite{ref077,ref234}。

\SourceParagraph{ch03-3.5.2.7-002}{195}
在有边界的世界面上，产生的边界项为
\begin{equation}\label{eq:3.149}
\delta S_B=\mathrm{i}\int_{\partial\Sigma}\mathrm{d}\sigma_1\,
\Lambda_J(\phi)\partial_1\phi^J,
\end{equation}
其形式，正是\BookEqRef{eq:3.140}{3.140}在实联络的约定下的形式。这意味着，\BookEqRef{eq:3.147}{3.147}单独并不是对称。相反，在同时作出
\begin{equation}\label{eq:3.150}
\delta B=\mathrm{d}\Lambda,\qquad
\delta A=-\Lambda|_L
\end{equation}
时，理论具有对称；这一变换保持组合 $B|_L+F$ 不变。

\SourceParagraph{ch03-3.5.2.7-003}{195}
这有多种含义，后面将继续讨论。最简单的一点，是在 $B\ne0$ 时，\BookEqRef{eq:3.142}{3.142}变为
\begin{equation}\label{eq:3.151}
S=\int_L\mathrm{d}^nx\,
\sqrt{\det(g|_L+B|_L+F)}.
\end{equation}
这里，各矩阵仍采用相同的弦尺度归一化。这也可以从\BookEqRef{eq:3.75}{3.75}中，度规与 B 场以 $g+B$ 的组合出现这一事实看出。

\SourceParagraph{ch03-3.5.2.7-004}{195}
当然，一般的 $U(1)$ 丛联络，并不由一个全局定义的1形式 $A$ 描述。在最一般的情形中，B 也不能只看作全局定义的2形式，而应看作 $U(1)$ gerbe 上的联络。更详细的讨论，见文献\BookCite{ref283}。

\subsubsection{Chan--Paton 因子与 Yang--Mills 场}\label{sec:3.5.2.8}

\SourceParagraph{ch03-3.5.2.8-001}{195}
在物理讨论中，携带非 Abelian 联络的 Dirichlet 膜，以及由多个 Dirichlet 膜组成的更一般系统，通过引入“Chan--Paton 因子”定义。这些因子，是从某个指标集合 $I$ 中取得的标签。随后，将一个 Dirichlet 膜配置，定义为从这个集合，到简单边界条件集合的映射 $B\colon I\to B_0$。开弦 Hilbert 空间为
\[
\mathcal{H}=\bigoplus_{I,J}\mathcal{H}_{IJ},
\qquad\mathcal{H}_{IJ}\cong\mathcal{H}_{B(I),B(J)}.
\]
相应边界算符的算符乘积，必须保持这一结构，即给出线性映射
\[
\mathcal{H}_{IJ}\otimes\mathcal{H}_{JK}\longrightarrow\mathcal{H}_{IK}.
\]
当然，这些内容，都包含在 Dirichlet 膜构成一个加法范畴的陈述中。不过，其中还有一些有意思的物理细节与推论。

\SourceParagraph{ch03-3.5.2.8-002}{196}
首先，考虑 $N$ 个相同的 Dirichlet 膜。也就是说，有子集 $I'\subset I$，基数为 $|I'|=N$，其中每个元素，都映到同一个 $B\in B_0$。为明确起见，取 $I'=\{1,2,\ldots,N\}$。前面讨论的边界算符，现在带有重数；将它们的系数，即耦合，看作 $N\times N$ 矩阵，就可以概括所得结构。例如，可以将\BookEqRef{eq:3.143}{3.143}，推广为边界形变
\begin{equation}\label{eq:3.152}
\delta S=\int_{\partial\Sigma}\mathrm{d}t\,
A^{JI}_K(\phi(t))\partial_t\phi^K(t),
\end{equation}
其中，$A^{JI}$ 的矩阵结构，与\BookSectionRef{sec:3.1.5}{3.1.5}中“带色粒子”的处理方式相同。

\SourceParagraph{ch03-3.5.2.8-003}{196}
对于\BookEqRef{eq:3.144}{3.144}，这一解释很明显：耦合 $A$ 是一个1形式矩阵，在几何上，是秩为 $N$ 的向量丛 $V$ 上的联络。这提示，而且确实成立，\BookSectionRef{sec:3.5.2.6}{3.5.2.6}中关于单个膜的 $U(1)$ 联络满足 Maxwell 方程的陈述，推广为 $U(N)$ 联络满足 Yang--Mills 方程，$D*F=0$，这里仍指主导的低能近似。

\SourceParagraph{ch03-3.5.2.8-004}{196}
要说明这一点，需要推广\BookSectionRef{sec:3.5.2.2}{3.5.2.2}的论证，并表明，现在得到非 Abelian 规范不变性。其无穷小形式为
\begin{equation}\label{eq:3.153}
A\longmapsto A+\mathrm{d}f+[A,f],
\qquad f\in\operatorname{Mat}_N(C(L)),
\end{equation}
其中，$C(L)$ 表示这里的函数空间。若能说明这一点，再结合运动方程应具有规范不变性的物理相容论证，以及主导的导数展开，就得到 Yang--Mills 方程。

\SourceParagraph{ch03-3.5.2.8-005}{196}
这一推广的标准讨论\BookCite{ref395}，相当复杂。出发点是，像\BookSectionRef{sec:3.1.5}{3.1.5}一样，用路径有序指数，书写具有非 Abelian 联络的边界作用量\BookEqRef{eq:3.143}{3.143}。不过，在弦理论中，这需要额外重整化。考虑这一指数展开的二次项，示意写为
\[
\sum_{i,j}g^ig^j\int\mathrm{d}t_1\int\mathrm{d}t_2\,
O_i(t_1)O_j(t_2),
\]
其中，积分区域按路径顺序规定。在重合极限 $t_1=t_2$，它通常发散。重整化这一发散，会在耦合定义中留下有限歧义。例如，可以选择加入额外贡献
\[
\sum_{i,j,k}C_{ij}^{\ k}g^ig^j
\int\mathrm{d}t_1\,O_k(t_1),
\]
称为“接触项”。

\SourceParagraph{ch03-3.5.2.8-006}{197}
虽然这似乎使泛函积分的定义，具有大量歧义，但其主要来源只是以下事实：\BookEqRef{eq:3.143}{3.143}中的耦合常数 $g^i$，在前面讨论的各种具体形式中，是边界条件空间上的一组局部坐标。当然，如果没有额外结构，就没有优选的坐标系。可以看出，一般的非线性重新定义 $g^i\mapsto g^{i\prime}(g)$，会引入这样的接触项。这也并非边界耦合特有；CFT 或 QFT 的空间，也有类似讨论。

\SourceParagraph{ch03-3.5.2.8-007}{197}
为了确定接触项，标准的物理分析，诉诸这里不作解释的物理相容条件。对于像 $A(\phi)$ 这样的“自旋1”场理论，如果有规范不变性，就可以满足这些条件，从而倾向于选择给出\BookEqRef{eq:3.153}{3.153}的接触项。这一论证并非最终结论；后来的研究提出，还可以作不同选择，例如得到非交换规范理论\BookCite{ref418}。

\SourceParagraph{ch03-3.5.2.8-008}{197}
要令人满意地处理这些问题，需要弦场论，如文献\BookCite{ref168,ref169}中的表述。正如\BookChapterRef{2}所述，这一框架与 $A_\infty$ 结构密切相关。

\subsubsection{Dirichlet 矩阵边界条件}\label{sec:3.5.2.9}

\SourceParagraph{ch03-3.5.2.9-001}{197}
现在，转向表达式\BookEqRef{eq:3.145}{3.145}。形式上，讨论完全平行，给出以下结果：如果考虑 $N$ 个相同的 Dirichlet 膜，每个都嵌入子流形 $L\subset M$，那么，它们嵌入的无穷小变化，最初似乎取值于 $NL\otimes\mathbb{R}^N$，实际却提升为 $NL$ 上取值于 $N\times N$ 矩阵的场；在幺正约定中，使用 Hermitian 矩阵。

\SourceParagraph{ch03-3.5.2.9-002}{197}
如果读者此前没有见过这一点，它应当相当出人意料。这清楚地表明，非交换几何，在 Dirichlet 膜理论中起着根本作用。后面讨论的发展，以及许多其他发展，确实展示了这一点；不过，这里恐怕仍有许多尚待阐明的内容。

\SourceParagraph{ch03-3.5.2.9-003}{197}
先问一个更基本的问题：这 $N^2\dim NL$ 个函数，如何与描述 $N$ 个嵌入子流形形变的 $N\dim NL$ 个函数相关？答案是，规范不变性与阻碍共同作用，使额外的 $(N^2-N)\dim NL$ 个函数，并不都对应于独立的、能够积分为有限形变的边缘方向。\BookSectionRef{sec:3.5.4}{3.5.4}、\BookSectionRef{sec:3.6.2.3}{3.6.2.3}与\BookRef{ex:5.21}{例题5.21}将进一步解释。

\subsection{超弦、GSO 投影与定向}\label{sec:3.5.3}

\SourceParagraph{ch03-3.5.3-intro-001}{197}
到目前为止，只讨论了世界面玻色子；这一讨论，同样适用于玻色弦与超弦理论。

\SourceParagraph{ch03-3.5.3-intro-002}{197}
超弦中费米子的处理，主要由保持 $(1,1)$ 世界面超对称的要求决定。不过，其中还涉及一些离散选择。特别地，Dirichlet 膜世界体的定向选择，正是在这里出现的。拓扑弦理论的讨论，经常省去它，因为即使不显式说明，也能够作出全部基本定义，甚至一直得到 Gromov--Witten 不变量。不过，下面需要这一选择。

\SourceParagraph{ch03-3.5.3-intro-003}{198}
完整讨论，需要解释所谓 GSO，即 Gliozzi--Scherk--Olive，投影。遗憾的是，这会涉及希望避免的超弦理论技术细节，因此，只针对这里的目的，简化讨论。完整说明，可以在教材\BookCite{ref191,ref394}中找到。

\subsubsection{超对称 sigma 模型中的边界条件}\label{sec:3.5.3.1}

\SourceParagraph{ch03-3.5.3.1-001}{198}
设 $L$ 是靶空间 $X$ 的子流形。与前面一样，要求 $\phi(\partial\Sigma)\subset L$，并在 $L$ 的法方向施加 Dirichlet 条件，在切方向施加修改后的 Neumann 条件。

\SourceParagraph{ch03-3.5.3.1-002}{198}
这些条件，可以紧凑地写为
\begin{equation}\label{eq:3.154}
\frac{\partial\phi^I}{\partial z}
=R^I_{\ J}(\phi)\frac{\partial\phi^J}{\partial\bar z}
+\text{费米项},
\end{equation}
其中，$R$ 是相对于度规 $g_{IJ}$ 的正交矩阵。$R$ 本征值为 $-1$ 的本征向量，给出 Dirichlet 条件，张成 $L$ 的法方向。通过2形式 $F\in\bigwedge^2T^*L$，参数化 $R$ 在切向量上的作用：
\[
R|_{TL}=(g|_L-F)^{-1}(g|_L+F).
\]
$F$ 应为闭形式，表示 $L$ 上规范场的曲率。将边界条件改写为
\begin{equation}\label{eq:3.155}
g_{IJ}\partial_2\phi^J
=F_{IJ}\partial_1\phi^J+\text{费米项},
\end{equation}
这一点更为清楚；这里使用与矩阵 $R$ 相容的世界面及联络约定。如果将含费米子的项置为零，就正好得到带规范场曲率 $F$ 的修改后的 Neumann 条件。

\SourceParagraph{ch03-3.5.3.1-003}{198}
超弦的费米子边界条件，则由保持 $N=1$ 世界面超对称的要求确定。由变换规则\BookEqRef{eq:3.84}{3.84}，可以理解，这要求左右行费米子，通过相同的矩阵 $R$ 相联系：
\begin{equation}\label{eq:3.156}
\psi_+^I=R^I_{\ J}\psi_-^J.
\end{equation}
进一步分析，可以推广为完整的非线性处理。这要求在\BookEqRef{eq:3.154}{3.154}中加入费米项；在多个膜的情形中，还需要像\BookSectionRef{sec:3.1.5}{3.1.5}那样的显式边界自由度。参见有关文献，例如\BookCite{ref334}。

\SourceParagraph{ch03-3.5.3.1-004}{198}
这些边界条件，在经典层面共形不变。在量子层面，共形不变性，等价于相应的时空场运动方程；与玻色弦一样，在所讨论的慢变近似中，它们由 Born--Infeld 作用量\BookEqRef{eq:3.151}{3.151}导出。

\subsubsection{GSO 投影}\label{sec:3.5.3.2}

\SourceParagraph{ch03-3.5.3.2-001}{199}
注意，到目前为止，还不需要指定 $L$ 的定向或自旋结构。这在量子化理论时出现，可以从多种角度看出。

\SourceParagraph{ch03-3.5.3.2-002}{199}
先说明，它如何进入算符表述。首先，希望在 $\mathcal{H}$ 上，定义算符 $(-1)^F$，使它与超荷反对易：
\[
(-1)^FG+G(-1)^F=0,
\]
对于 $G_+$、$G_-$、$\widetilde G_+$ 与 $\widetilde G_-$ 都成立。

\SourceParagraph{ch03-3.5.3.2-003}{199}
在 sigma 模型中，做法很明显：只需像\BookEqRef{eq:3.6}{3.6}一样，使 $(-1)^F$ 与世界面费米子反对易。更一般地，由于 $G$ 的 R 荷为 $q=+1$，在相应的整荷扇区及归一化下，可以取
\begin{equation}\label{eq:3.157}
(-1)^F=C\mathrm{e}^{\mathrm{i}\pi(q_L+q_R)},
\end{equation}
其中，$C=\pm1$ 是尚未确定的常数。

\SourceParagraph{ch03-3.5.3.2-004}{199}
GSO 投影，本质上是投影到满足
\begin{equation}\label{eq:3.158}
(-1)^F=+1
\end{equation}
的扇区。由于 $(-1)^F$ 是乘法守恒的量子数，即 $\mathbb{Z}_2$ 分次，投影与相互作用相容。

\SourceParagraph{ch03-3.5.3.2-005}{199}
这一表述有所简化。不过，进一步讨论细节以前，先说明它与 $L$ 或 $X$ 的定向有何关系。实际上，定义涉及完整的十维时空，例如 $\mathbb{R}^{3,1}\times X$，或者其他相应的时空，记为 $M$。正如玻色弦 Hilbert 空间，是 $M$ 上函数空间的扩张，超弦的费米子，又将它扩张到\BookSectionRef{sec:3.1.1}{3.1.1}中 $N=1$ 超对称量子力学的 Hilbert 空间；后者使用 $M$ 上的旋量丛 $S$ 定义。因此，可以预期，投影\BookEqRef{eq:3.158}{3.158}，实现约化到单个手征性 $S^+$ 或 $S^-$，取决于\BookEqRef{eq:3.157}{3.157}中 $C$ 的选择。这一选择，隐含地选定了 $M$ 的定向。

\SourceParagraph{ch03-3.5.3.2-006}{199}
严格地说，前面的陈述，只在 Ramond 扇区成立，因为只有在这一扇区，费米子才有构成 Clifford 代数的零模。不过，由于在拓扑弦理论中，Neveu--Schwarz 与 Ramond 扇区，通过谱流基本上可以相互联系，见\BookSectionRef{sec:3.3.3}{3.3.3}，这里可以暂时略去这一细节。同样，只提到在闭弦中，实际上应分别定义左右行算符 $(-1)^{F_L}$ 与 $(-1)^{F_R}$，并对二者都作投影，然后继续下面的讨论。

\SourceParagraph{ch03-3.5.3.2-007}{199}
由此，看到 $M$ 的定向如何进入理论。那么，$L$ 的定向又如何进入？量子化开弦时，需要选择 $M$ 上旋量的分解。为确定起见，取 $S(M)^+$；在切、法丛均为偶秩的相应情形中，有
\[
S(M)^+\cong
\bigl(S(TL)^+\otimes S(NL)^+\bigr)
\oplus\bigl(S(TL)^-\otimes S(NL)^-\bigr).
\]
这里，$TL$ 与 $NL$，分别是 $L$ 作为 $M$，例如当前的 $\mathbb{R}^{3,1}\times X$，的子流形的切丛与法丛。关键是，两个因子的手征性选择彼此关联。

\SourceParagraph{ch03-3.5.3.2-008}{200}
由于 $\mathbb{R}^{3,1}$ 中两个旋量的相对手征性，在物理上可以观测，它进入宇称与荷共轭等离散对称的定义，因此，可以利用这一关联，区别 $L$ 上旋量的两个手征性，从而定义 $L$ 的定向。

\SourceParagraph{ch03-3.5.3.2-009}{200}
详细发展这一论证，需要比这里能涵盖的更多物理知识。不过，可以作出一个简单观察：弦理论既能看到法丛中旋量的手征性，也能看到切丛中的。因此，即使在 $L$ 是 $M$ 中一点这样看似退化的情形中，定向也会出现。这称为“膜”与“反膜”的区别，\BookSectionRef{sec:5.1}{5.1}将详细讨论。

\SourceParagraph{ch03-3.5.3.2-010}{200}
还有另一个重要观察，有助于回答以下问题：如果讨论中需要定向，如何在 SCFT 中定义它，而不引入显式的几何概念？特别地，镜像对称联系两个非常不同的几何图景，因此，如果两个几何定义，背后有同一个非几何定义，会很有帮助。

\SourceParagraph{ch03-3.5.3.2-011}{200}
关键的观察是，所有定向依赖，都可以追溯到\BookEqRef{eq:3.157}{3.157}；这一算符，纯粹用 R 荷定义。因此，在讨论中每一个需要指定定向的地方，都可以用某个仅依赖于 R 荷的条件来替代。\BookChapterRef{5}将多次使用这一观察。

\subsection{Dirichlet 膜的 T 对偶}\label{sec:3.5.4}

\SourceParagraph{ch03-3.5.4-001}{200}
现在，精确表述\BookSectionRef{sec:1.3.1}{1.3.1}引入的思想：T 对偶对 Dirichlet 膜的作用，是把它变成一个维数不同的 Dirichlet 膜。

\SourceParagraph{ch03-3.5.4-002}{200}
取 $M\cong T^d$，沿用\BookSectionRef{sec:3.2.3.6}{3.2.3.6}的泛函积分方法。那些论证，可以直接推广到具有边界的情形。

\SourceParagraph{ch03-3.5.4-003}{200}
事实上，无需更多细节，就可以看出主要内容。结合\BookEqRef{eq:3.61}{3.61}与\BookEqRef{eq:3.63}{3.63}，Euclidean 约定下的核心关系为
\[
\mathrm{d}\widehat\varphi_i
=\mathrm{i}g_{ij}*\mathrm{d}\varphi^j.
\]
可以说，T 对偶作为 CFT 的一个对偶，来自世界面场 $\varphi$ 的 Hodge 对偶。在边界附近，若取定向正交坐标 $(t,n)$，使 $*\mathrm{d}t=\mathrm{d}n$、$*\mathrm{d}n=-\mathrm{d}t$，便得到
\begin{equation}\label{eq:3.159}
\partial_t\widehat\varphi_i
=-\mathrm{i}g_{ij}\partial_n\varphi^j,
\end{equation}
\begin{equation}\label{eq:3.160}
\partial_n\widehat\varphi_i
=\mathrm{i}g_{ij}\partial_t\varphi^j.
\end{equation}
经过相应的 Wick 旋转，这些关系成为 Lorentz 世界面上的实对偶关系。回顾\BookSectionRef{sec:3.5.2.1}{3.5.2.1}，可见，世界面上的 T 对偶，直接交换 Dirichlet 与 Neumann 边界条件的作用。例如，$S^1$ 上的 D0膜，对 $\varphi$ 施加 Dirichlet 条件，将成为对偶 $S^1$ 上的 D1膜，对 $\widehat\varphi$ 施加 Neumann 条件。这直接推广到 $T^d$：对于按这些坐标区分的 $p$ 个 Neumann 与 $d-p$ 个 Dirichlet 方向，若对全部 $d$ 个坐标作 T 对偶，Dp膜就成为对偶 $T^d$ 上的 D$(d-p)$膜。

\SourceParagraph{ch03-3.5.4-004}{201}
继续考虑\BookSectionRef{sec:3.5.2.6}{3.5.2.6}，还可以看出，T 对偶交换形变算符；在相应的实坐标与归一化下，
\[
v_i\int\mathrm{d}t\,\partial_n\varphi^i
\quad\longleftrightarrow\quad
A_i\int\mathrm{d}t\,\partial_t\widehat\varphi_i.
\]
前者生成 $T^d$ 中的移动。由于 CFT 是自由的，可以将这一无穷小变化积分，得到边界条件 $\varphi=\varphi_0+v$。

\SourceParagraph{ch03-3.5.4-005}{201}
后一个算符，初看似乎是边界上的全导数；不过，由于 $\widehat\varphi$ 是到 $S^1$ 的映射，需要更加仔细。使用展开\BookEqRef{eq:3.67}{3.67}，可见，这一作用量项，正比于 $A_jn_j$，并在泛函积分中给出因子 $\exp(\mathrm{i}A_jn^j)$。它正是平坦联络 $A_j\mathrm{d}x^j$，沿由 $n\in\pi_1(T^d)$ 表示的路径的平行移动。因此，最初所提出的对应
\begin{equation}\label{eq:3.161}
\varphi^i\longleftrightarrow A_i,
\end{equation}
即 D0膜的位置，与 Dd膜平坦联络的参数之间的对应，确实正确；圆周坐标与 Wilson 线参数的尺度，由前面的周期归一化确定。

\SourceParagraph{ch03-3.5.4-006}{201}
这一论证有一个令人困惑之处：环面上的平坦 $U(1)$ 联络空间，有一个标记点，即平凡平坦联络；但对偶环面上 D0膜的模空间，没有一个典范标记点。为解决这一问题，回忆，共形场论具有规范对称\BookEqRef{eq:3.150}{3.150}。如果取 $\Lambda$ 为闭1形式，B 场不变，而规范场 $A$，可以平移一个任意闭1形式。平凡平坦联络，在这些变换下不保持不变；这意味着，平坦规范场 $A$ 的空间，没有与全部 B 场规范选择无关的自然原点。进一步讨论，见\BookSectionRef{sec:6.3}{6.3}。

\subsection{由 T 对偶导出 Born--Infeld 作用量}\label{sec:3.5.5}

\SourceParagraph{ch03-3.5.5-001}{201}
现在，利用 T 对偶，导出 Born--Infeld 作用量
\begin{equation}\label{eq:3.162}
S=\int_L\mathrm{d}^nx\,\sqrt{\det(g+F)}.
\end{equation}
这里，$g$ 是时空度规在 $L$ 上的拉回，看作对称的 $n\times n$ 矩阵；$F=\mathrm{d}A$ 是曲率2形式，看作反对称的 $n\times n$ 矩阵；$\det$ 是两者相加所得矩阵的行列式。正如前面所述，当 $F$ 在弦尺度 $\ell_s$ 上缓慢变化时，这一作用量的临界点，对应于边界 beta 函数的零点。思路是，使用 T 对偶，将一个 $F$ 近乎常值的膜，与一个 $F=0$ 的膜联系起来。对于后者，假定作用量已知，就是它的体积。

\SourceParagraph{ch03-3.5.5-002}{201}
考虑一个 $F=0$ 的 Dp膜，它是平凡 $T^d$ 纤维化 $Y\to\mathbb{R}^p$ 的截面 $L$，并沿纤维作 T 对偶。设 $\varphi\colon\mathbb{R}^p\to T^d$ 定义这一截面。拉回度规为
\[
\iota^*g=g_{ij}\mathrm{d}x^i\mathrm{d}x^j
+g_{ab}\mathrm{d}\varphi^a\mathrm{d}\varphi^b,
\]
其中，$1\le i,j\le p$、$1\le a,b\le d$。假定 $g_{ia}=0$，但允许 $g_{ab}$ 随底空间变化。

\SourceParagraph{ch03-3.5.5-003}{202}
将它代入体积\BookEqRef{eq:3.141}{3.141}，得到
\[
S=\int_L\mathrm{d}^px\,
\sqrt{\det_{1\le i,j\le p}
\bigl(g_{ij}+g_{ab}\partial_i\varphi^a\partial_j\varphi^b\bigr)}.
\]
现在，使用 T 对偶关系\BookEqRef{eq:3.161}{3.161}，将纤维位置识别为对偶 Wilson 线参数，令 $F_{ia}=\partial_iA_a$，它变为
\begin{equation}\label{eq:3.163}
S=\int_L\mathrm{d}^px\,
\sqrt{\det_{1\le i,j\le p}
\bigl(g_{ij}+g_{ab}F_{ia}F_{jb}\bigr)}.
\end{equation}

\SourceParagraph{ch03-3.5.5-004}{202}
在对偶度规
\[
g'=g_{ij}\mathrm{d}x^i\mathrm{d}x^j
+(g_{\mathrm{fib}}^{-1})^{ab}\mathrm{d}y_a\mathrm{d}y_b,
\]
中，这与\BookEqRef{eq:3.142}{3.142}几乎相同，只差一个由纤维体积与伸缩子变换给出的整体因子。将原纤维度规矩阵记为 $H=(g_{ab})$，将 $p\times d$ 曲率矩阵记为 $F=(F_{ia})$，则
\[
\det\begin{pmatrix}g_{\mathrm{base}}&F\\-F^{\mathsf{T}}&H^{-1}\end{pmatrix}
=(\det H)^{-1}\det(g_{\mathrm{base}}+FHF^{\mathsf{T}}).
\]
要得到这一关系，就是将\BookEqRef{eq:3.163}{3.163}的 $p\times p$ 行列式，与完整的 $(p+d)\times(p+d)$ 行列式相联系；它来自标准的分块行列式恒等式\footnote{也可以先对矩阵作归一化，将 $p\times p$ 行列式化为 $\det g_{\mathrm{base}}$ 乘以 $\det(1+KK^{\mathsf T})$，其中 $K=g_{\mathrm{base}}^{-1/2}FH^{1/2}$；完整分块行列式则相应化为
\[
\det\begin{pmatrix}1_p&K\\-K^{\mathsf T}&1_d\end{pmatrix}
=\det(1_p+KK^{\mathsf T}).
\]
也可使用形式展开 $\det(1+M)=\exp\bigl(\sum_{k\ge1}(-1)^{k+1}\operatorname{tr}(M^k)/k\bigr)$，逐项比较。}。

\SourceParagraph{ch03-3.5.5-005}{202}
因此，对偶理论中的 D$(p+d)$膜，由稍作修改的 Born--Infeld 作用量描述：
\begin{equation}\label{eq:3.164}
S=\int_{L'}\mathrm{d}^{p+d}x\,
\mathrm{e}^{D'}\sqrt{\det(g'+F)},
\end{equation}
其中，相对于原来的权重，$D'=D+\tfrac12\log\det H$。这里取纤维坐标的固定周期，并略去共同的膜张力常数。这一额外因子，正是闭弦理论中允许的伸缩子耦合；若用通常的伸缩子 $\Phi=-D$，其变换为 $\Phi'=\Phi-\tfrac12\log\det H$。虽然它在后面的讨论中，不起主要作用，但在弦理论中通常十分重要，因为它控制时空量子修正的强度，即弦耦合常数。不过，它在这里的出现，也提示，当前基于 $\alpha'$ 级数展开的考察，在奇异纤维处，预期会遇到困难。

\SourceParagraph{ch03-3.5.5-006}{202}
对于 $N$ 个 Dirichlet 膜，相应的描述，只有部分得到理解。Born--Infeld 作用量\BookEqRef{eq:3.142}{3.142}，适用于 $F$ 在弦尺度上变化缓慢，但其大小可以很大的极限。Abelian 情形中，$F$ 是规范不变的，因此，这一说法有明确意义；一般的非 Abelian 情形中，则复杂得多。定义非 Abelian Born--Infeld 作用量的一个简单方案，是“对称化迹规定”\BookCite{ref450}；不过，正如该文强调的那样，这不是精确结果。相关讨论还见文献\BookCite{ref376}。

\section{超对称与拓扑边界条件}\label{sec:3.6}

\SourceParagraph{ch03-3.6-intro-001}{203}
现在，已经可以与\BookChapterRef{1}的讨论联系起来。在那里，曾论证，时空超对称的条件，要求相应的几何 D 膜，位于校准子流形上。

\SourceParagraph{ch03-3.6-intro-002}{203}
一旦 $\Sigma$ 具有边界，就不可能保持整个 $N=(2,2)$ SCA，因为反射条件\BookEqRef{eq:3.156}{3.156}，将左右行费米子相联系。另一方面，为了能够作拓扑扭曲，必须能够相容地定义一个 $N=2$ SCA。此外，它必须与\BookSectionRef{sec:3.3.4.6}{3.3.4.6}中世界面算符 $\Omega$ 的存在相容。

\SourceParagraph{ch03-3.6-intro-003}{203}
有两种方法可以做到这一点，对应于两种扭曲。正如预期，主要条件是，\BookEqRef{eq:3.100}{3.100}中定义的相应 A 模型或 B 模型的 $Q$ 算符，必须在开弦扇区中具有良好定义。利用 SCA，这就确定了 $U(1)_R$ 流 $J$ 的边界条件：
\begin{equation}\label{eq:3.165}
J_L=-J_R,\qquad\text{A 模型},
\end{equation}
\begin{equation}\label{eq:3.166}
J_L=J_R,\qquad\text{B 模型}.
\end{equation}
由\BookEqRef{eq:3.139}{3.139}可见，如果将 $(J_L,J_R)$ 看作来自一个自由玻色子，这两个条件，分别对应于 Dirichlet 与 Neumann 条件。

\SourceParagraph{ch03-3.6-intro-004}{203}
这两种边界条件，都将 $(2,2)$ SCA 的双分次，约化为一个由互补的流线性组合定义的单分次。因此，两种模型中的守恒荷为
\begin{equation}\label{eq:3.167}
\begin{aligned}
J_A&=J_{L,0}-J_{R,0}
\ \propto\int\mathrm{d}\sigma\,\partial_\sigma\varphi,
&&\text{A 模型},\\
J_B&=J_{L,0}+J_{R,0}
\ \propto\int\mathrm{d}\sigma\,\partial_\tau\varphi,
&&\text{B 模型}.
\end{aligned}
\end{equation}
这里的积分表达式，采用\BookSectionRef{sec:3.2.4}{3.2.4}与\BookSectionRef{sec:3.3.4.6}{3.3.4.6}讨论的 $U(1)$ 代数的自由玻色子实现，并略去共同的流归一化，以联系\BookSectionRef{sec:3.2.3}{3.2.3}的显式处理；$\tau$ 是世界面时间，$\sigma$ 是空间坐标。A 模型的荷是相应的拓扑荷，B 模型的荷是相应的动量荷。

\subsection{A 模型}\label{sec:3.6.1}

\SourceParagraph{ch03-3.6.1-001}{203}
A 扭曲\BookEqRef{eq:3.115}{3.115}，与边界条件\BookEqRef{eq:3.156}{3.156}相容，要求在全纯坐标中，$R^i_{\ j}=R^{\bar\imath}_{\ \bar\jmath}=0$。也就是说，只有非对角分量 $R^i_{\ \bar\jmath}$ 与 $R^{\bar\imath}_{\ j}$ 非零。

\SourceParagraph{ch03-3.6.1-002}{203}
为简单起见，取 $F=0$，并选择 $R$ 本征值为 $+1$ 的向量 $v$，即 $L\subset Y$ 的切向量。引入概复结构 $J$；在全纯坐标中，它为
\begin{equation}\label{eq:3.168}
J^n_{\ m}=\mathrm{i}\delta^n_m,\qquad
J^{\bar n}_{\ \bar m}=-\mathrm{i}\delta^{\bar n}_{\bar m},
\end{equation}
非对角分量为零。于是，容易看出，$Jv$ 是 $R$ 本征值为 $-1$ 的向量，也就是 $L$ 的法向量。此外，$J^2v=-v$，再次施加 $J$，又回到切方向。因此，$J$ 交换 D 膜 $L$ 的切方向与法方向，显然，$L$ 必须具有中间维数。

\SourceParagraph{ch03-3.6.1-003}{204}
注意，如果 $v$ 与 $w$ 都是 $L$ 的切向量，即在 $R$ 下具有本征值 $+1$，那么，$w$ 与 $Jv$，相对于度规 $g_{IJ}$ 正交。由于 $Y$ 上的 K\"ahler 形式，就是由 $g$ 与 $J$ 定义的2形式 $\omega(u,v)=g(Ju,v)$，可见，K\"ahler 形式在 $L$ 上的限制为零。因此，$L$ 是 $Y$ 的 Lagrange 子流形，这就与\BookChapterRef{1}中的 A 型膜相联系。

\SourceParagraph{ch03-3.6.1-004}{204}
A 膜还有额外的量子约束。对于当前情形，可以概括为以下定义。
\BookStatementNumber{3}
\begin{definition}\label{def:3.3}
一个 Lagrange A 膜，是 $Y$ 中配备平坦联络的 Lagrange 三维子流形，在 Hamilton 形变下的等价类；它具有平凡的 Maslov 类，见\BookEqRef{eq:3.170}{3.170}，并满足量子阻碍条件\BookEqRef{eq:3.181}{3.181}。
\end{definition}

\SourceParagraph{ch03-3.6.1-005}{204}
不过，在详细解释这些约束以前，先介绍其他已知的 A 膜。

\subsubsection{余各向同性膜}\label{sec:3.6.1.1}

\SourceParagraph{ch03-3.6.1.1-001}{204}
允许 $F\ne0$ 的更细致分析\BookCite{ref382,ref291}表明，Calabi--Yau $n$ 维流形，可以具有实维数为 $n+2p$ 的 A 膜，其中，$p$ 是非负整数。对 $L$ 与 $F$ 的精确条件，相当复杂。
\begin{definition*}\label{def:3.coisotropic-submanifold}
首先，子流形 $L$ 必须余各向同性。这意味着，对于任意点 $x\in L$，$T_xL$ 相对于辛形式 $\omega$ 的辛正交补，包含于 $T_xL$ 中。
\end{definition*}
等价地，$L$ 局部由函数 $f_i$，$i=1,\ldots,\operatorname{codim}_{\mathbb{R}}L$，的零点定义，并且这些函数的任意两个 Poisson 括号，在 $L$ 上都为零。按照 Dirac 的术语，$L$ 是第一类约束曲面。因此，可以把 $f_i$，看作规范变换的生成元，或者 Hamilton 函数。相应的向量场，切于 $L$，并张成一个维数为 $\operatorname{codim}_{\mathbb{R}}L$ 的可积分布。积分这一分布，得到 $L$ 的叶状结构，其中的叶，可以看作“规范轨道”。将每片叶上的全部点等同起来，得到一个拓扑空间，称为叶空间；按照 Dirac 的术语，它是约化相空间。一般而言，它既不一定是流形，也不一定是 Hausdorff 拓扑空间。

\SourceParagraph{ch03-3.6.1.1-002}{204}
对 $(L,F)$ 的第二个条件是，2形式 $F$，与任意沿上述叶方向的切向量收缩，都为零。由于 $F$ 闭，这还意味着，它是叶状流形 $L$ 上的基本2形式；如果约化相空间确实是一个流形，那么，$F$ 下降为其上的闭2形式 $f$。注意，虽然辛形式 $\omega$ 在 $L$ 上的限制退化，但它在横于叶的方向，下降为非退化的2形式。因此，$\omega$ 在约化相空间上，给出一个辛形式 $\sigma$；这是 Dirac 的第一类约束系统理论中的关键观察之一。

\SourceParagraph{ch03-3.6.1.1-003}{205}
最后，第三个条件要求，约化相空间上的 $(1,1)$ 张量 $J=\sigma^{-1}f$，满足 $J^2=-1$，也就是一个概复结构。利用 $\sigma$ 与 $f$ 都是闭2形式，可以证明，$J$ 自动可积\BookCite{ref290}，因此，约化相空间是复流形。此外，按照这一符号约定，$f+\mathrm{i}\sigma$ 是关于 $J$ 的全纯辛形式；因此，约化相空间的复维数必须为偶数，记为 $2p$\footnote{强调一点：约化相空间的复结构 $J$，与 Calabi--Yau $n$ 维流形 $Y$ 的复结构无关。}。容易看出，$L$ 的实维数为 $n+2p$。

\SourceParagraph{ch03-3.6.1.1-004}{205}
可以用广义复几何的语言，更优雅地表述上述对 $(L,F)$ 的三个条件\BookCite{ref236,ref212,ref284}，见\BookSectionRef{sec:6.2.5}{6.2.5}。当 $F=0$ 时，条件简化为 $L$ 是 Lagrange 子流形；此时，$p=0$，在连通情形中，约化相空间是单个点。对于 Calabi--Yau 一维流形，即椭圆曲线，这就是唯一的可能性。对于 Calabi--Yau 三维流形 $Y$，另一种可能性是 $p=1$，而 $L$ 是 $Y$ 中的实五维子流形。如果 $Y$ 是严格的 Calabi--Yau 三维流形，即 完整群不是 $\operatorname{SU}(3)$ 的真子群，那么，$b_5(Y)=0$，这样的闭 $L$ 在实同调中平凡。不过，由于规范场不平坦，这样的膜仍然可以稳定\footnote{严格地说，文献\BookCite{ref291}的分析，只适用于 A 膜上的 Abelian 规范场，因此，只能对这样的 A 膜得出所述结论。不过，很有理由预期，在适当情形中，秩为 $N$ 的向量丛所配备的膜，可以看作 $N$ 个重合的秩1膜的形变。}。

\subsubsection{鬼数反常与 Maslov 类}\label{sec:3.6.1.2}

\SourceParagraph{ch03-3.6.1.2-001}{205}
给定一个配备平坦丛的 Lagrange 子流形，量子考察还施加两个约束。这里讨论第一个：A 膜必须保持算符乘积代数的 $U(1)_R$ 分次，即鬼数。在量子理论中，这并不平凡，因为这一对称可以具有反常。

\SourceParagraph{ch03-3.6.1.2-002}{205}
$U(1)_R$ 反常的计算，见 \BookBibText{ref242}{MS1} 第40章；这里解释其几何来源。它与 Floer 上同调的分次问题相关\BookCite{ref377,ref378}。

\SourceParagraph{ch03-3.6.1.2-003}{205}
固定 $Y$ 上一个全纯3形式 $\Omega$。在有定向的 Lagrange 子流形 $L$ 上的任意点 $x$，其体积形式，可以写成全纯形式的限制：
\begin{equation}\label{eq:3.169}
\mathrm{d}V_L=w\,\mathrm{e}^{-\mathrm{i}\pi\xi(x)}\Omega|_L,
\end{equation}
其中，$w$ 为正实数，随点而变。局部的角度 $\xi$，给出 $L$ 到圆周的相位映射。其平方相位 $\mathrm{e}^{2\pi\mathrm{i}\xi}$，与定向反转无关，并诱导基本群上的映射
\begin{equation}\label{eq:3.170}
\mu_L\colon\pi_1(L)\longrightarrow\pi_1(S^1)\cong\mathbb{Z}.
\end{equation}
可以把它看作 $H^1(L,\mathbb{Z})$ 的元素，称为 $L$ 的 Maslov 类。
\begin{figure}[H]
\centering
\BookImage{01_Numbered_Figures/Chapter_03/Ch03_Fig07_PDFp206_Bookp196.pdf}
\caption{不产生反常与产生反常的闭曲线。}\label{fig:3.7}
\end{figure}

\SourceParagraph{ch03-3.6.1.2-004}{206}
反常恰在 $L$ 的 Maslov 类为零时消失。显然，当 $\pi_1(L)=0$ 时，总是如此。
\begin{example*}
考虑复一维环面 $Y$。其中任意光滑曲线，都是 Lagrange 的。不过，如\BookRef{fig:3.7}{图7}，一条可缩的嵌入闭曲线，具有非平凡的 Maslov 类，因此，不能作为这里要求保持整数鬼数分次的 A 膜。
\end{example*}

\SourceParagraph{ch03-3.6.1.2-005}{206}
余各向同性 A 膜的鬼数反常，已在文献\BookCite{ref335}中计算。Maslov 类的类似物如下。设 $\dim_{\mathbb{R}}L=n+2p$，考虑 $L$ 上的最高次形式 $\Omega|_L\wedge F^{\wedge p}$。可以证明，这一形式在 $L$ 上处处非零，因此，正比于 $L$ 的体积形式。比例系数，是一个处处非零的函数；其相位，定义一个映射 $\pi_1(L)\to\mathbb{Z}$。这一映射为零，等价于没有鬼数反常。

\subsection{A 膜的开弦}\label{sec:3.6.2}

\SourceParagraph{ch03-3.6.2-intro-001}{206}
现在，讨论一对 A 膜 $(L_1,E_1,\nabla_1)$ 与 $(L_2,E_2,\nabla_2)$ 之间的开弦谱。

\SourceParagraph{ch03-3.6.2-intro-002}{206}
先从 $L_1=L_2=L$ 为同一个 Lagrange 子流形的情形开始，Witten 在文献\BookCite{ref463}中分析了它。一般的情形，将在\BookSectionRef{sec:3.6.2.5}{3.6.2.5}处理。

\subsubsection{从 A 膜到自身的开弦}\label{sec:3.6.2.1}

\SourceParagraph{ch03-3.6.2.1-001}{206}
取 $L_1=L_2=L$，其余数据也相同。设有一个 Lagrange 循环 $L$，配备 $U(N)$ 向量丛 $E\to L$。如前所述，开弦态是 $E^*\otimes E$ 的截面，或者再与其他丛作张量积，也就是说，它们取值于 $E$ 的自同态。如果 $E$ 是线丛，$\operatorname{End}(E)$ 的截面，就是 $L$ 上的函数。

\SourceParagraph{ch03-3.6.2.1-002}{206}
开弦谱的基本分析，与\BookSectionRef{sec:3.4.1}{3.4.1}平行。在经典或大体积近似中，全部 $Q$ 上同调，都可以用以下形式的算符代表：
\begin{equation}\label{eq:3.171}
a_{I_1I_2\cdots}(\phi)\chi^{I_1}\chi^{I_2}\cdots,
\end{equation}
其中，$a_{I_1I_2\cdots}(\phi)$ 取值于 $\operatorname{End}(E)_{\mathbb{C}}$，$\phi\in L$，而 $\chi^{I_k}$ 属于 $L$ 的切丛。注意，函数 $a_{I_1I_2\cdots}$ 允许取复值；更精确地，它们是 $\operatorname{End}(E)$ 复化后的截面。BRST 算符 $Q$ 的作用，类似于\BookSectionRef{sec:3.4.1}{3.4.1}，因此，这一近似中的开弦态空间，由总 de Rham 上同调给出：
\begin{equation}\label{eq:3.172}
\mathcal{H}\cong\bigoplus_kH^k(L,\operatorname{End}(E)_{\mathbb{C}}),
\end{equation}
其中，鬼数为 $k$。同样的讨论，也适用于两个包覆同一循环 $L$，但携带不同向量丛 $E_1$ 与 $E_2$ 及不同平坦联络的 A 膜之间的开弦；此时，得到
\begin{equation}\label{eq:3.173}
\mathcal{H}\cong\bigoplus_k
H^k(L,\operatorname{Hom}(E_1,E_2)_{\mathbb{C}}).
\end{equation}

\SourceParagraph{ch03-3.6.2.1-003}{207}
\BookSectionRef{sec:3.4.1}{3.4.1}中作用量形变的讨论，也推广到边界情形。现在，形变算符，沿一维流形 $\partial\Sigma$ 积分，而不是在二维 $\Sigma$ 上积分。因此，寻找鬼数为1的边界算符，使其鬼数为零的后裔，可以加入边界作用量。它们对应于 $H^1(L,\operatorname{End}(E)_{\mathbb{C}})$ 中的元素。

\SourceParagraph{ch03-3.6.2.1-004}{207}
在 $E$ 是线丛的最简单情形中，可以从几何上理解这一点。这些形变的实自由度中，一半对应于 $E$ 上平坦幺正联络的形变，另一半对应于 $L$ 作为 Lagrange 子流形的同痕类形变，见\BookSectionRef{sec:6.1.1}{6.1.1}。基本事实，可能已从经典力学中熟悉：Lagrange 子流形的法向形变 $v\in NL$，对应于 $L$ 上的闭1形式 $\theta$，它局部可以通过一个 Hamilton 函数 $H$ 写成
\[
v^j=\omega^{ij}\theta_i\sim\omega^{ij}\partial_iH,
\]
其中，符号取决于 Hamilton 向量场的约定。在全局具有这一形式的形变，就是 Hamilton 同痕。因此，$H^1(L,\mathbb{R})$ 的元素，对应于模掉 Hamilton 同痕后的 Lagrange 形变。这说明，拓扑理论中的 A 膜，按相应的 Hamilton 同痕等价来识别。

\SourceParagraph{ch03-3.6.2.1-005}{207}
单个秩1余各向同性 A 膜的开弦谱，已在经典近似中计算\BookCite{ref289}。一种表述方式如下：正如\BookSectionRef{sec:3.6.1.1}{3.6.1.1}中所述，余各向同性 A 膜，具有自然的叶状结构，以及横于叶方向的自然复结构。因此，可以考虑沿叶局部常值、沿横向方向全纯的函数层。开弦态空间，就是这一层的上同调。对于 Lagrange A 膜，这个层就是 $L$ 上局部常值的 $\mathbb{C}$ 值函数层，于是，退化为\BookEqRef{eq:3.172}{3.172}的秩1情形。

\subsubsection{关联函数与超势}\label{sec:3.6.2.2}

\SourceParagraph{ch03-3.6.2.2-001}{207}
给定开弦算符以后，接下来计算它们的关联函数。仍然可以写为
\begin{equation}\label{eq:3.174}
\langle W_{a_1}\cdots W_{a_k}\rangle
=\int[\mathrm{d}\phi\,\mathrm{d}\chi\,\mathrm{d}\psi]\,
\mathrm{e}^{-S}W_{a_1}\cdots W_{a_k}.
\end{equation}
现在，世界面 $\Sigma$ 具有边界，算符 $W_{a_i}$ 插入在边界上。主要的情形，有时称为“树级”，是 $\Sigma$ 为具有单个边界的圆盘。

\SourceParagraph{ch03-3.6.2.2-002}{208}
一般而言，这一计算\BookCite{ref463}，与\BookSectionRef{sec:3.4.1.1}{3.4.1.1}中闭弦 A 模型的计算，非常平行。同样的局域化论证表明，关联函数是对与边界条件相容的全纯映射 $\phi\colon\Sigma\to Y$ 的求和。其中包括常值映射，给出经典项；也包括非平凡映射，给出世界面瞬子贡献。

\SourceParagraph{ch03-3.6.2.2-003}{208}
考虑以下算符的三点关联函数：
\begin{equation}\label{eq:3.175}
\begin{aligned}
a&\in H^1(L,\operatorname{Hom}(E_1,E_2)),\\
b&\in H^1(L,\operatorname{Hom}(E_2,E_3)),\\
c&\in H^1(L,\operatorname{Hom}(E_3,E_1)).
\end{aligned}
\end{equation}
其经典贡献为
\begin{equation}\label{eq:3.176}
\langle W_aW_bW_c\rangle_{\mathrm{cl}}
=\int_L\operatorname{Tr}(a\wedge b\wedge c).
\end{equation}
在这里，关键地使用了 $Y$ 是 Calabi--Yau 三维流形这一事实。

\SourceParagraph{ch03-3.6.2.2-004}{208}
瞬子修正的一般形式，沿用\BookEqRef{eq:3.116}{3.116}的上同调参数与实联络约定，写为
\begin{equation}\label{eq:3.177}
\sum_{D_\alpha}\pm\,
\exp\left(2\pi\mathrm{i}\int_{D_\alpha}(B+\mathrm{i}\omega)
-\mathrm{i}\oint_{\partial D_\alpha}A_i\right)
N_{abc}^{\alpha},
\end{equation}
其中，求和遍及全部满足 $\partial D_\alpha\subset L$ 的全纯圆盘，包括多重覆盖；$N_{abc}^{\alpha}$ 表示相应的模空间计数。

\SourceParagraph{ch03-3.6.2.2-005}{208}
这里，$A_i$ 是丛 $E_i$ 上的联络。更精确地，将边界 $\partial D_\alpha$ 分成三段，以 $i\in\{1,2,3\}$ 标记，在第 $i$ 段使用 $A_i$。开弦顶点算符，位于相邻边界段的接点，例如，$W_a$ 位于第1段与第2段之间。当然，这一规定区分算符的次序，因此，会给出结合、但不必交换的开弦上同调代数。更多细节，将在下面，特别是\BookChapterRef{8}中说明。

\SourceParagraph{ch03-3.6.2.2-006}{208}
虽然不能直接将\BookEqRef{eq:3.175}{3.175}推广到四点及更高点关联函数，但还有另一个定义，反映 $H^1$ 的元素，对应于线性化形变这一数学陈述。给定 R 荷为1的算符 $W$，可以选取 R 荷为零的算符1形式 $W^{(1)}$，满足
\begin{equation}\label{eq:3.178}
\mathrm{d}W=\{Q,W^{(1)}\}.
\end{equation}
于是，可以向边界作用量加入
\begin{equation}\label{eq:3.179}
\delta S_b=\int_{\partial\Sigma}W^{(1)}.
\end{equation}
由于 $W^{(1)}$ 的 R 荷为零，这至少在线性化层面，保持 R 荷守恒与共形不变性。

\SourceParagraph{ch03-3.6.2.2-007}{209}
将\BookEqRef{eq:3.179}{3.179}插入\BookEqRef{eq:3.174}{3.174}，可以定义四点关联函数；继续这一过程，就定义了 $n$ 点函数。

\SourceParagraph{ch03-3.6.2.2-008}{209}
给定一个显式的特殊 Lagrange 子流形，可以看出，这与前面\BookEqRef{eq:3.140}{3.140}中 $U(1)$ 联络变化的定义一致。不过，当前讨论适用于一般的边界 SCFT。

\SourceParagraph{ch03-3.6.2.2-009}{209}
圆盘上的开弦关联函数生成函数，通常称为超势。现在，可以直接类比\BookSectionRef{sec:3.3.5.3}{3.3.5.3}，定义
\begin{equation}\label{eq:3.180}
\mathcal{W}[t]=\left\langle
\exp\left(\sum_i t_i\int_{\partial\Sigma}O_i^{(1)}\right)
\right\rangle.
\end{equation}
这里的积分算符，用\BookEqRef{eq:3.178}{3.178}定义；不过，还需记住，展开中的 $n=\widehat c$ 个插入，取为不积分的原算符 $O_i$，以满足鬼数选择规则。可以用 $N=2$ SCA 的性质说明，选择哪些算符积分，并不影响最后结果。

\SourceParagraph{ch03-3.6.2.2-010}{209}
与闭弦情形的一个区别是，由于算符 $O_i$ 沿圆盘边界积分，可以保留它们的相对次序。这使耦合 $t_i$，能够提升为矩阵，正如\BookSectionRef{sec:3.5.2.8}{3.5.2.8}所述。虽然这一点很重要，但与闭弦情形相比，还有一个更重要的区别，下面转向它。

\subsubsection{蝌蚪项与阻碍}\label{sec:3.6.2.3}

\SourceParagraph{ch03-3.6.2.3-001}{209}
尽管闭弦与开弦计算，在形式上相似，但存在一个关键区别，具有重要的数学与物理后果。在两种情形中，都可以将相应的弦态，看作形变配置空间的切向量；毕竟，弦态来自完整弦动力学问题的线性化。在 A 模型中，$H^{1,1}(Y,\mathbb{C})$ 的元素，是复化 K\"ahler 类的形变；而如前所述，$H^1(L,\mathbb{C})$ 的元素，对应于 Lagrange 子流形 $L$ 与 $E$ 上平坦联络的形变。

\SourceParagraph{ch03-3.6.2.3-002}{209}
这些看似非常相似的形变，关键区别在于：$(2,2)$ SCFT 中所讨论的闭弦模参数形变，没有阻碍，而开弦形变却可以有，并且经常有阻碍。从物理上说，闭弦态，起初也只是运动方程线性化的解；不过，可以证明，在这里的超对称条件下，它们确实是完整方程的一族解的切向量\footnote{必须强调，这一陈述，关键地依赖于所假定的超对称类型。在更一般的情形中，闭弦形变同样可以有阻碍。}。开弦形变，则未必如此。

\SourceParagraph{ch03-3.6.2.3-003}{210}
先在大结构极限，即 $\ell_s\to0$ 的极限中，显式地看出这一点。Calabi--Yau 几何的“闭弦运动方程”，是 Ricci 平坦条件。对于固定的复结构，Yau 定理说明，它的解，由 $Y$ 的 K\"ahler 锥中的点参数化。这是一个开空间，因此，上述无阻碍陈述很明显；还须考虑 B 场的运动方程，不过，它是线性的。B 模型中对应的陈述，是复结构形变没有阻碍；对于 Calabi--Yau 三维流形，这一点在文献\BookCite{ref448,ref449}中证明。

\SourceParagraph{ch03-3.6.2.3-004}{210}
在有限 $\ell_s$ 时，有一个基于 $N=(2,2)$ SCFT 的论证，说明相应的闭弦形变没有阻碍\BookCite{ref124}。另一个相关但不同的论证，使用 IIA 或 IIB 型弦理论在 Calabi--Yau 三维流形上紧化，给出四维 $N=2$ 超引力理论这一事实。结果，来自\BookSectionRef{sec:5.4}{5.4}中将讨论的那类时空论证，使用四维 $N=2$ 超引力中一个事实：相应的无荷场超势为零。

\SourceParagraph{ch03-3.6.2.3-005}{210}
现在，转向开弦运动方程。对于 Lagrange 子流形 $L$ 上的 A 膜，方程包括 $E$ 上的联络平坦。对于秩为 $N$ 的丛，平坦联络模空间的维数，依赖于 $\pi_1(L)$ 的表示；在下面关注的可分解为不同平坦线丛的分支中，典型实维数为 $Nb_1(L)$。另一方面，还存在自同态群更大的特殊联络。极端情形，是同一个线丛 $\ell$ 的 $N$ 个副本，$E\cong\mathbb{C}^N\otimes\ell$；此时，实幺正形变的线性化空间，具有
\[
\dim_{\mathbb{R}}H^1(L,\operatorname{ad}E)=N^2b_1(L).
\]
当 $N>1$ 时，它大于上述分支的维数。这正是\BookSectionRef{sec:3.5.2.9}{3.5.2.9}提出的问题。

\SourceParagraph{ch03-3.6.2.3-006}{210}
在一般的可对角化配置上，$U(N)$ 规范变换的非对角轨道，具有 $N^2-N$ 个维数。不过，对于 $b_1(L)>1$，这一计数仍不足以解释差异；剩下的 $(b_1(L)-1)(N^2-N)$ 个形式方向，并非全部能够成为无阻碍的独立形变。在重合联络的奇异点处，还须同时考虑稳定子，不能单凭线性规范商消去全部差异。

\SourceParagraph{ch03-3.6.2.3-007}{210}
提出问题以后，一般的解决原理并不神秘：$E$ 上联络平坦这一条件，是非线性的，因此，产生阻碍。在最简单的 $L\cong T^3$ 情形中，很容易看出这一点，因为在相应的平坦分支中，可以选取一个标架，使联络由一组常值矩阵表示。于是，平坦条件成为
\[
[A_i,A_j]=0,\qquad\text{对所有 }i,j.
\]
正如预期，一般的矩阵形变 $A_i\ne0$，会有阻碍。在这一例子中，阻碍理论，由这些二次约束给出。

\SourceParagraph{ch03-3.6.2.3-008}{210}
一般的阻碍理论，即使在经典层面，也没有这么简单。而一旦看到，开弦关联函数，会受到\BookEqRef{eq:3.177}{3.177}那样的量子修正，就必须意识到，这些修正，也会修改阻碍理论。

\SourceParagraph{ch03-3.6.2.3-009}{211}
物理学为前面的论证提供动机，并以稍后将说明的方式，进一步刻画阻碍理论。不过，在本书所讨论的物理框架中，量子阻碍的一般形式，并未从全部世界面计算中完成导出\footnote{许多具体例子已知，最早可见文献\BookCite{ref275}。}。尽管如此，其形式很清楚。对于携带线丛的 Lagrange A 膜，圆盘贡献的相应阻碍条件，示意为
\begin{equation}\label{eq:3.181}
\sum_{\alpha\in I}\pm\,
\exp\left(2\pi\mathrm{i}\int_{D_\alpha}(B+\mathrm{i}\omega)
-\mathrm{i}\oint_{\partial D_\alpha}A\right)
[\partial D_\alpha]=0
\quad\text{于 }H_1(L).
\end{equation}
这里，对全部边界在 $L$ 上的全纯圆盘 $D_\alpha$ 求和，包括多重覆盖，并按相应模空间的计数加权；系数取在容纳这些形式瞬子级数的环中。记号 $[\partial D_\alpha]$，表示 $\partial D_\alpha$ 的同调类\footnote{说明条件\BookEqRef{eq:3.181}{3.181}在适当的 Hamilton 形变下保持不变，是一个有意思的练习。}。

\SourceParagraph{ch03-3.6.2.3-010}{211}
在这一未加入额外补偿形变的构造中，要求 Lagrange 子流形，满足\BookEqRef{eq:3.181}{3.181}，才能给出所述 A 膜。这显然有解，例如，$L\cong S^3$ 的 $H_1$ 为零，自动满足这一同调条件。更一般的探索，见文献\BookCite{ref276,ref277}；不过，一般的几何理解，仍需进一步发展。例如，对于 SYZ 的三环面纤维化\BookCite{ref433}，是否满足这一条件，目前尚不清楚。还应注意，条件\BookEqRef{eq:3.181}{3.181}，依赖于 $B+\mathrm{i}\omega$ 与联络 $A$ 的取值。因此，一个 A 膜，可能只在这些参数的某些特定取值下有效，而非对全部取值都有效。

\SourceParagraph{ch03-3.6.2.3-011}{211}
出于一些不那么直观的物理原因，条件\BookEqRef{eq:3.181}{3.181}常称为“蝌蚪项消除”。实际上，“蝌蚪项为零”与“蝌蚪项消除”，在相关物理语境中，有许多不同含义，因此，最好避免使用这些名称。

\subsubsection{超势与 \texorpdfstring{$A_\infty$}{A∞} 结构}\label{sec:3.6.2.4}

\SourceParagraph{ch03-3.6.2.4-001}{211}
从物理上，理解这些阻碍的最简单方法，是使用\BookSectionRef{sec:5.4}{5.4}将讨论的世界体方法。从这一观点看，阻碍条件，都可以写成超势\BookEqRef{eq:3.180}{3.180}梯度的各个分量。\BookChapterRef{5}将给出一些例子，如\BookEqRef{eq:5.86}{5.86}与\BookEqRef{eq:5.82}{5.82}。

\SourceParagraph{ch03-3.6.2.4-002}{211}
还有一些十分有意思的约束，最适合从世界面的观点理解。先简要讨论，\BookSectionRef{sec:3.3.4.3}{3.3.4.3}中闭弦手征代数结构的开弦类似物。

\SourceParagraph{ch03-3.6.2.4-003}{211}
通过同样基于 OPE\BookEqRef{eq:3.99}{3.99}的论证，开弦手征算符在上同调上，构成一个代数。它不必交换，但仍然结合。因此，自然会尝试，将 WDVV 的讨论，推广到非交换情形。

\SourceParagraph{ch03-3.6.2.4-004}{211}
不过，形变受到阻碍这一事实，再次带来与闭弦讨论的根本区别；它超出了非交换性。可以把它理解为，形变会修改 $Q$ 上同调。这仍然可以从 $L\cong T^3$ 上平坦丛的最简单例子看出：$H^p(L,\operatorname{End}E)$ 的维数，随着自同态群改变，会发生不连续变化。在更物理的语言中，这些形变，以破坏\BookEqRef{eq:3.96}{3.96}的方式，改变算符维数，类似于\BookSectionRef{sec:3.2.3.5}{3.2.3.5}中的例子。

\SourceParagraph{ch03-3.6.2.4-005}{212}
在这种情形下，如果只考虑 BRST 上同调中算符的关联函数，就不可避免地丢失信息。另一方面，按照目前的表述，世界面的论证，并不容易直接推广到更一般的算符。

\SourceParagraph{ch03-3.6.2.4-006}{212}
解决这一困难的方法，是将结合代数的概念，扩展为 $A_\infty$ 代数。\BookChapterRef{8}将讨论它的定义与性质。按照那里定义的术语，BRST 算符 $Q$，识别为 $m_1$，手征环乘法识别为 $m_2$。于是，只需允许 $m_1\ne0$，就可以纳入不属于 BRST 上同调的算符。

\SourceParagraph{ch03-3.6.2.4-007}{212}
随后，将用\BookEqRef{eq:3.179}{3.179}定义的高点关联函数，也就是由超势\BookEqRef{eq:3.180}{3.180}编码、出现在形变理论中的关联函数，看作定义一系列高阶乘法：
\begin{equation}\label{eq:3.182}
m_k(a_1,a_2,\ldots,a_k),\qquad k=3,4,\ldots.
\end{equation}
可以通过世界面论证说明，这些乘法满足 $A_\infty$ 代数的公理\BookCite{ref230}，并满足某些额外的“高亏格”关系。应当将这组关系，看作 WDVV 方程的开弦类似物。

\SourceParagraph{ch03-3.6.2.4-008}{212}
膜范畴上的 $A_\infty$ 结构，在 Kontsevich 最初的同调镜像对称表述中，起着根本作用。\BookChapterRef{8}将详细讨论这一点。

\subsubsection{多个 A 膜的开弦}\label{sec:3.6.2.5}

\SourceParagraph{ch03-3.6.2.5-001}{212}
设有一组 A 膜 $L_a$。为便于说明，先假定每个丛 $E_a$ 都是线丛。给定一对 A 膜 $L_a$ 与 $L_b$，就有从 $L_a$ 开始、在 $L_b$ 结束的开弦态空间。它带有分次，除去一个加法常数，就是鬼数。这个分次平移，十分重要，下面将详细讨论。对这一分次态空间，采用记号
\begin{equation}\label{eq:3.183}
\operatorname{Hom}^*(L_a,L_b)
=\bigoplus_{m\in\mathbb{Z}}\operatorname{Hom}^m(L_a,L_b).
\end{equation}
也将 $\operatorname{Hom}^0(L_a,L_b)$，简称为 $\operatorname{Hom}(L_a,L_b)$。

\SourceParagraph{ch03-3.6.2.5-002}{212}
由刚才的一般论证，这样一个开弦态空间中的态，对应于一个开弦手征顶点算符。这些算符的关联函数，定义 A 膜的范畴，并最终给出一个 $A_\infty$ 范畴。

\SourceParagraph{ch03-3.6.2.5-003}{213}
正如前面提到的那样，一般认为，所得结构，与 Fukaya 引入、并以其命名的范畴相同\BookCite{ref165}。\BookChapterRef{8}将更深入地讨论。这里说明这一看法背后的物理论证，即哪些考察，确定开弦态空间的维数及其关联函数。

\SourceParagraph{ch03-3.6.2.5-004}{213}
为了计算两个不同 Lagrange A 膜 $L_1$ 与 $L_2$ 之间的开弦态空间，最简单的做法，是假定它们在有限多个点横截相交。与\BookSectionRef{sec:3.6.2}{3.6.2}一样，利用拓扑场论的 $Q$ 不变性，可以将经典候选态局域化到常值映射 $\phi\colon\Sigma\to Y$。这意味着，一个这样的开弦态，与 $L_1\cap L_2$ 的一点相关。

\SourceParagraph{ch03-3.6.2.5-005}{213}
前面的讨论提示，局部的态空间，应由交集的 de Rham 上同调给出，也就是一点的上同调。因此，最初的猜测是，每个交点，对应一个一维态空间；于是，$\operatorname{Hom}^*(L_1,L_2)$ 的维数，似乎应等于 $L_1\cap L_2$ 的点数。不过，这不可能是最终答案。A 模型在 $L_1$ 或 $L_2$ 的 Hamilton 形变下不变，而交点数并不保持不变。

\SourceParagraph{ch03-3.6.2.5-006}{213}
最简单的不变量，是有定向的相交数 $\#(L_1\cap L_2)$，因为它只依赖于同调类。为了定义它，需要规定 $L_1$ 与 $L_2$ 的相对定向；注意，这正是定向开始进入当前交点讨论的地方。当然，它只给出膜的拓扑信息。

\SourceParagraph{ch03-3.6.2.5-007}{213}
给出正确态空间的 A 模型物理，是具有圆盘拓扑的全纯世界面瞬子。它们对 BRST 算符 $Q$ 产生修正，将不同交点处的开弦态配对，使其中一些态从上同调中消失。

\SourceParagraph{ch03-3.6.2.5-008}{213}
引入一些记号。将 $L_1\cap L_2$ 的交点，标记为 $p_a$，$a=1,\ldots,M$。相应的开弦顶点算符 $W_{p_a}$，在点 $p_a$ 产生一条开弦。将候选态空间记为 $V\cong\mathbb{C}^M$。每个顶点算符 $W_{p_a}$ 的鬼数，记为 $\mu(p_a)$。于是，$V$ 按鬼数分次：
\begin{equation}\label{eq:3.184}
V_i=\bigoplus_{\mu(p_a)=i}\mathbb{C},
\qquad V=\bigoplus_iV_i.
\end{equation}

\SourceParagraph{ch03-3.6.2.5-009}{214}
世界面瞬子修正，可以参数化为
\begin{equation}\label{eq:3.185}
\{Q,W_{p_a}\}=\sum_b n_{ab}W_{p_b},
\end{equation}
其中，$n_{ab}$ 是待确定的系数。因此，真正的态空间，是基于向量空间 $V$ 的某个复形的 $Q$ 上同调。由于 $Q$ 的鬼数为1，这一复形为
\begin{equation}\label{eq:3.186}
\cdots\xrightarrow{Q}V_{-1}\xrightarrow{Q}V_0
\xrightarrow{Q}V_1\xrightarrow{Q}\cdots.
\end{equation}
将位置 $i$ 处的上同调，定义为 $\operatorname{Hom}^i(L_a,L_b)$。

\SourceParagraph{ch03-3.6.2.5-010}{214}
为了计算 $n_{ab}$，必须在瞬子的模空间上积分；除了显然的玻色映射 $\phi$，还须对费米参数积分。按照费米积分的通常规则，如果这些参数没有通过插入或相应项被饱和，积分就为零。更精确地，要求瞬子的 Dirac 算符指标，等于 $Q$ 的鬼数，即1\BookCite{ref466}。

\SourceParagraph{ch03-3.6.2.5-011}{214}
Dirac 算符的指标，也测量全纯映射模空间的一般预期维数，更精确地，是虚拟维数。有关非一般位置情形的良好说明，见文献\BookCite{ref469}。在一般位置下，模掉相应的平移以后，只需计数零维瞬子模空间中的点，就得到 $n_{ab}$。

\SourceParagraph{ch03-3.6.2.5-012}{214}
对于从 $p_a$ 到 $p_b$ 的瞬子，Dirac 算符的指标，是鬼数差 $\mu(p_b)-\mu(p_a)$。因此，预期模掉平移后的瞬子模空间维数为
\begin{equation}\label{eq:3.187}
\dim\mathcal{M}=\mu(p_b)-\mu(p_a)-1.
\end{equation}
进一步说明，见文献\BookCite{ref150}。

\SourceParagraph{ch03-3.6.2.5-013}{214}
这样，物理定义，原则上将系数 $n_{ab}$，确定为定义良好的模空间积分。实际上，在这里所讨论的物理文献中，这些计算，并未全部完成。相反，刚才描述的形式体系，与 Fukaya 范畴结构的相似性，使人们普遍相信，完整的物理处理，会与 Fukaya 范畴的定义一致。

\subsubsection{鬼数}\label{sec:3.6.2.6}

\SourceParagraph{ch03-3.6.2.6-001}{214}
细心的读者，可能已经注意到，还没有指定实际计算 $\mu(p_a)$ 的方法。给定瞬子模空间的维数，关系\BookEqRef{eq:3.187}{3.187}，只足以计算 $L_a$ 与 $L_b$ 两个交点的相对鬼数。事实上，有下面这一事实，它在\BookChapterRef{5}中十分重要。
\BookStatementNumber{4}
\begin{claim}\label{claim:3.4}
如果尚未为 A 膜指定额外的分次数据，拓扑 A 模型本身，不足以确定两个 A 膜交点处开弦的绝对鬼数。
\end{claim}

\SourceParagraph{ch03-3.6.2.6-002}{215}
鬼数究竟有多少歧义？给定一对 A 膜 $L_1$ 与 $L_2$，可以将从 $L_1$ 到 $L_2$ 的开弦鬼数，整体平移一个固定整数。\BookSectionRef{sec:3.6.2}{3.6.2}还说明，如果 $L_1=L_2$，鬼数由 de Rham 上同调的次数给出，具有良好定义。此外，希望算符乘积保持鬼数：
\begin{equation}\label{eq:3.188}
\operatorname{Hom}^i(L_1,L_2)\otimes\operatorname{Hom}^j(L_2,L_3)
\longrightarrow\operatorname{Hom}^{i+j}(L_1,L_3).
\end{equation}
可以通过给每个 A 膜本身，赋予鬼数 $\mu(L)$，来说明这一歧义。随后，将原来 $\operatorname{Hom}^i(L_a,L_b)$ 中一个元素的鬼数，定义为
\begin{equation}\label{eq:3.189}
i+\mu(L_b)-\mu(L_a).
\end{equation}
容易看出，这一定义具有所需的全部性质。

\SourceParagraph{ch03-3.6.2.6-003}{215}
可以将上述内容重新表述如下：拓扑 A 模型有一种对称性，允许平移开弦态的鬼数：给各个 A 膜赋予任意整数鬼数，再按\BookEqRef{eq:3.189}{3.189}定义开弦态的鬼数，就实现了这一平移。注意，文献\BookCite{ref419}，仔细研究了给 Lagrange 子流形赋予分次，以消除这一歧义的思想。
\begin{figure}[H]
\centering
\BookImage{01_Numbered_Figures/Chapter_03/Ch03_Fig08_PDFp215_Bookp205.pdf}
\caption{与三点函数相关的圆盘瞬子。}\label{fig:3.8}
\end{figure}

\SourceParagraph{ch03-3.6.2.6-004}{215}
这里不详细说明如何计算关联函数。不过，应当清楚，其中涉及瞬子修正。例如，为计算树级三点函数，考虑 $Y$ 中具有\BookRef{fig:3.8}{图8}所示边界条件的全纯圆盘。费米零模的饱和，像往常一样，要求鬼数守恒。

\SourceParagraph{ch03-3.6.2.6-005}{215}
本节介绍了 A 膜范畴的定义，在其中，对象 $L_a$ 与 $L_b$ 横截相交。实际上，总可以使用 Hamilton 形变，把任意一对 Lagrange 子流形变到这种情形。因此，我们实际上已经得到了 Lagrange A 膜范畴的完整定义。

\SourceParagraph{ch03-3.6.2.6-006}{216}
推广到每个 A 膜携带更高秩向量丛的情形，应当很清楚：不再给每个交点赋予 $\mathbb{C}$，而是赋予表示一个丛在交点处的纤维，到另一个丛纤维的线性映射的矩阵。推广到余各向同性 A 膜，则困难得多，仍需要进一步发展。

\SourceParagraph{ch03-3.6.2.6-007}{216}
强调一点，A 模型中没有任何东西依赖于 $Y$ 的复结构。这从上面的 Fukaya 范畴定义中，并不明显，因为计算态射的规则，涉及计数到 $Y$ 的全纯映射，后者依赖于复结构或相容概复结构的选择。不过，Fukaya 范畴的对象，以及态射的复合，都依赖于 $B+\mathrm{i}\omega$。蝌蚪项条件\BookEqRef{eq:3.181}{3.181}，依赖于 $B+\mathrm{i}\omega$，因此，某些对象，可能只在这一参数的特定取值下存在。关联函数，通过瞬子修正依赖于 $B+\mathrm{i}\omega$，所以，态射复合也具有同样的依赖。

\SourceParagraph{ch03-3.6.2.6-008}{216}
最后，应当指出，一般预期，世界面瞬子，会对基于时空度规的概念产生不利影响。因此，有理由预期，用 Lagrange 子流形作为物理 A 膜的几何描述，只在大半径极限中才真正成立。Fukaya 范畴的复合规则，基于与瞬子效应相关的幂级数。越过这些幂级数的收敛半径以后，有理由认为，用 Lagrange 子流形描述 A 膜的做法已经失效。

\subsection{B 型膜}\label{sec:3.6.3}

\SourceParagraph{ch03-3.6.3-001}{216}
可以重复\BookSectionRef{sec:3.6.1}{3.6.1}开头的分析。B 膜的区别是，B 模型扭曲，要求\BookEqRef{eq:3.154}{3.154}与\BookEqRef{eq:3.156}{3.156}中的反射矩阵，满足 $R^i_{\ \bar\jmath}=R^{\bar\imath}_{\ j}=0$；也就是说，只有对角分量 $R^i_{\ j}$ 与 $R^{\bar\imath}_{\ \bar\jmath}$ 非零。

\SourceParagraph{ch03-3.6.3-002}{216}
这意味着，概复结构现在保持 D 膜的切方向与法方向，而非交换它们。因此，D 膜是 $X$ 的复子流形，显然，其实维数必须为偶数。

\SourceParagraph{ch03-3.6.3-003}{216}
注意，与 A 膜的讨论不同，这里已经使用 $X$ 的定向。如果保留这一信息，还有另一组 B 膜，即对于共轭复结构的复子流形；等价地，它们是反全纯子流形与丛。当然，可以对它们作平行讨论；不过，应记住，在完整弦理论中，作为 BPS 膜，它们与原来的选择不同。

\SourceParagraph{ch03-3.6.3-004}{216}
虽然有各种偶维的 B 膜，先将注意力限制到充满 $X$ 的 B 膜。也就是说，施加 Neumann 条件，或 Dirichlet--Neumann 混合边界条件。最终，通过结合最高维 B 膜的理解，与\BookChapterRef{4}的同调代数，就能够推导更一般 B 膜的性质。

\SourceParagraph{ch03-3.6.3-005}{217}
与 A 膜一样，考虑 B 场，迫使我们纳入 B 膜上丛的可能性，即 $E\to X$。令 B 场为零，要求作用量的 $Q$ 变分中，边界项消失，便得到对这个丛的约束：其曲率 $F$，是取值于 $\operatorname{End}(E)$ 的 $(1,1)$ 型2形式\BookCite{ref463,ref241,ref242}。换句话说，$E\to X$ 是全纯向量丛，其全纯截面层，是 $X$ 上的局部自由凝聚层。

\SourceParagraph{ch03-3.6.3-006}{217}
更一般地，$F$ 的 $(0,2)$ 部分，必须等于 B 场的 $(0,2)$ 部分的负值，乘以恒等自同态。对于通常满足 $H^{0,2}=0$ 的 Calabi--Yau 三维流形，B 场的 $(0,2)$ 部分，在上同调中平凡，可以通过相应的 BRST 恰当变化消去。因此，在这一情形中，B 场的这部分，不影响 B 膜范畴。

\SourceParagraph{ch03-3.6.3-007}{217}
如果 $B^{0,2}$ 在上同调中非平凡，例如 $X$ 为 K3 曲面或复环面时，情形就很不同。此时，不能简单地给充满空间的 B 膜，赋予 $X$ 上的普通凝聚层，因为通常的 $E$ 全纯截面概念，并不成立。相反，可以给它赋予 $X$ 上某个 gerbe 的凝聚层。简要说明，见例如文献\BookCite{ref284}。

\subsection{B 膜的开弦}\label{sec:3.6.4}

\SourceParagraph{ch03-3.6.4-001}{217}
利用\BookEqRef{eq:3.156}{3.156}，边界上的费米子 $\theta_j$，可以通过 $\eta^{\bar k}$ 表示：
\begin{equation}\label{eq:3.190}
\theta_j=g_{j\bar k}
\left(\psi_+^{\bar k}-\psi_-^{\bar k}\right)
=F_{j\bar k}\eta^{\bar k}.
\end{equation}
因此，局部算符只依赖于 $\phi$ 与 $\eta^{\bar\jmath}$，局部边界算符，由取值于 $\operatorname{End}(E)$ 的 $(0,q)$ 形式参数化。

\SourceParagraph{ch03-3.6.4-002}{217}
设有两个 B 膜，对应于两个向量丛 $E_1\to X$ 与 $E_2\to X$。Chan--Paton 自由度，与从 $E_1$ 到 $E_2$ 的映射相关。将这样的映射丛，记为 $\operatorname{Hom}(E_1,E_2)$。

\SourceParagraph{ch03-3.6.4-003}{217}
\BookSectionRef{sec:3.4.2}{3.4.2}中已看到，在 B 模型的大张力极限中，BRST 算符 $Q$，约化为 Dolbeault 算符。结合以上各项，一个从 $E_1$ 延伸到 $E_2$ 的开弦顶点算符，是以下上同调中的元素：
\begin{equation}\label{eq:3.191}
\bigoplus_qH^{0,q}_{\bar\partial}
\bigl(X,\operatorname{Hom}(E_1,E_2)\bigr).
\end{equation}

\SourceParagraph{ch03-3.6.4-004}{217}
与尚未指定分次的 A 膜情形相比，对于这些原始全纯向量丛对象，可以明确地将\BookEqRef{eq:3.191}{3.191}中算符的鬼数，规定为 $q$。

\SourceParagraph{ch03-3.6.4-005}{217}
像往常一样，B 模型没有瞬子修正。如果
\begin{equation}\label{eq:3.192}
\begin{aligned}
a&\in H^{0,p}_{\bar\partial}(X,\operatorname{Hom}(E_1,E_2)),\\
b&\in H^{0,q}_{\bar\partial}(X,\operatorname{Hom}(E_2,E_3)),\\
c&\in H^{0,r}_{\bar\partial}(X,\operatorname{Hom}(E_3,E_1)),
\end{aligned}
\end{equation}
并且 $p+q+r=\dim_{\mathbb{C}}X$，三点函数由经典表达式给出：
\begin{equation}\label{eq:3.193}
\langle W_aW_bW_c\rangle
=\int_X\operatorname{Tr}(a\wedge b\wedge c)\wedge\Omega.
\end{equation}
相应算符乘积代数，就是取值于 $\operatorname{Hom}$ 的形式的通常楔积，结合丛映射的复合。

\SourceParagraph{ch03-3.6.4-006}{218}
高点关联函数，按照\BookSectionRef{sec:3.6.2.2}{3.6.2.2}中 A 膜讨论的相同方式定义。相应的边界作用量修改，对应于用一个 $\delta A^{0,1}\in H^1(X,\operatorname{End}(E))$，形变 Dolbeault 算符，得到
\[
\bar\partial+\delta A^{0,1}.
\]
从层论的观点看，无穷小形变，取值于全局 Ext 群 $\operatorname{Ext}^1(E,E)$。

\SourceParagraph{ch03-3.6.4-007}{218}
Witten 在文献\BookCite{ref463}中说明，这些关联函数，可以由全纯 Chern--Simons 理论导出，见\BookSectionRef{sec:5.4.3.4}{5.4.3.4}。它们也组成一个 $A_\infty$ 结构，\BookChapterRef{8}将继续讨论。

\subsection{D0膜}\label{sec:3.6.5}

\SourceParagraph{ch03-3.6.5-001}{218}
另一个特别重要的 B 型膜，是 D0膜，也就是说，对全部玻色场 $\phi^i$，都施加纯粹的常值 Dirichlet 边界条件。这使世界面的边界，映到 $X$ 中的一点 $p$。

\SourceParagraph{ch03-3.6.5-002}{218}
分析其 BRST 不变的边界算符。BRST 不变性，要求费米场 $\eta^{\bar\imath}$ 在边界上为零，而 $\theta_i$ 任意。因此，鬼数为 $k$ 的最一般 BRST 不变顶点算符，为
\begin{equation}\label{eq:3.194}
a^{i_1\cdots i_k}\theta_{i_1}\cdots\theta_{i_k},
\qquad a^{i_1\cdots i_k}\in\mathbb{C}.
\end{equation}
由于 $\theta_i$ 取值于 $T_p^*X$，并且反对易，系数张量完全反对称，应当看作 $\bigwedge^kT_pX$ 的元素。

\SourceParagraph{ch03-3.6.5-003}{218}
再计算对应于全纯向量丛 $E$ 的 B 膜，到点 $p$ 处 D0膜的态射空间。此时，边界条件使相应顶点上的全部费米子消失，因此，边界算符可以与向量空间 $E_p^*$ 的元素等同起来。

\SourceParagraph{ch03-3.6.5-004}{218}
也可以继续计算，支撑在任意全纯子流形 $Z$ 上的 B 膜，到点 $p$ 处 D0膜的态射空间。这里不作计算，但有一个简单结果，对\BookSectionRef{sec:5.3.3.3}{5.3.3.3}很重要：只有 $p\in Z$ 时，才可能有非零态射。这是因为，态射对应于 Ramond 基态，而开弦任何非零的缠绕贡献，都会使 $L_0$ 本征值 $h$ 高于基态。这一缠绕贡献，受最小距离平方的下界控制；在略去共同弦尺度常数的单位中，可以示意写为
\[
h-h_{\mathrm{ground}}\ \gtrsim\min_{x\in Z}d^2(p,x).
\]

\subsubsection{摩天层}\label{sec:3.6.5.1}

\SourceParagraph{ch03-3.6.5.1-001}{219}
\BookChapterRef{4}将讨论凝聚层；粗略地说，它们推广了全纯丛。自然地，预期 D0膜在 $D^b(X)$ 中对应的对象，是 $\mathcal{O}_p$，即支撑在点 $p$ 上的摩天层。在 $D^b(X)$ 中，摩天层的分次自同态代数，是 $T_pX$ 的外代数，与边界算符空间 $\bigwedge^kT_pX$ 一致。它与其他对象之间的态射，也与刚才的讨论一致。

\subsection{超对称膜与广义校准}\label{sec:3.6.6}

\SourceParagraph{ch03-3.6.6-001}{219}
终于，可以展示超共形边界条件，与\BookSectionRef{sec:1.1.4}{1.1.4}定义的校准子流形之间的关系。上面定义的 A 膜与 B 膜，保持 $N=2$ 超对称，但未必保持共形或超共形不变性。给定 $N=2$ 超对称，共形不变性蕴涵超共形不变性，因此，只需要求前者。如前所述，这等价于边界 beta 函数为零；在 $F=0$ 的情形中，意味着 $L$ 必须是体积最小化（即极小）的子流形。

\SourceParagraph{ch03-3.6.6-002}{219}
对于连通的 Lagrange A 膜，在 Ricci 平坦 Calabi--Yau 度规下，极小性等价于
\begin{equation}\label{eq:3.195}
\Omega|_L=c\,\operatorname{vol}_L,
\end{equation}
其中，$\Omega$ 是全纯最高次形式，$c$ 是常数。这样的子流形，称为特殊 Lagrange 子流形。

\SourceParagraph{ch03-3.6.6-003}{219}
$F=0$ 的 B 膜，就是复子流形。K\"ahler 流形中，一个实维数为 $2p$ 的复子流形，自动极小，并满足
\[
\frac{\omega^p|_L}{p!}=\operatorname{vol}_L,
\]
也就是说，它由 $2p$ 形式 $\omega^p/p!$ 校准。

\SourceParagraph{ch03-3.6.6-004}{219}
还有一个更直接的方法，可以看出与校准几何的联系。优点是，无论 $F$ 是否为零，都以相同方式成立。思路是，要求边界条件保持 $\Upsilon^2$，即谱流算符的平方。先考虑 Lagrange A 膜。在世界面边界上，要求满足以下匹配条件：
\begin{equation}\label{eq:3.196}
\Omega_{i_1\cdots i_n}\psi_+^{i_1}\cdots\psi_+^{i_n}
=\overline{\Omega}_{\bar\imath_1\cdots\bar\imath_n}
\psi_-^{\bar\imath_1}\cdots\psi_-^{\bar\imath_n}.
\end{equation}
利用边界条件 $\psi_+=R\psi_-$，可以说明，这等价于 $\Omega|_L$ 与 $\operatorname{vol}_L$ 具有固定的比例相位。

\SourceParagraph{ch03-3.6.6-005}{219}
类似地，对于实维数为 $n+2p$ 的余各向同性 A 膜，谱流算符在边界上的匹配，给出条件\BookCite{ref287}
\[
\Omega|_L\wedge F^p=c\,\operatorname{vol}_L,
\]
其中，相应的广义校准条件固定 $c$ 的相位，体积密度与含规范场的 Born--Infeld 密度相联系。这推广了中间维数 A 膜的特殊 Lagrange 条件。

\SourceParagraph{ch03-3.6.6-006}{220}
对于复维数为 $p$ 的秩1 B 膜，边界上的谱流匹配条件为
\begin{equation}\label{eq:3.197}
\Omega_{i_1\cdots i_n}\psi_+^{i_1}\cdots\psi_+^{i_n}
=\Omega_{i_1\cdots i_n}\psi_-^{i_1}\cdots\psi_-^{i_n}.
\end{equation}
结合 $\psi_+=R\psi_-$，以及反射矩阵 $R$ 的显式形式，得到以下固定相位条件\BookCite{ref287}：
\begin{equation}\label{eq:3.198}
(\omega+\mathrm{i}F)^{\wedge p}
=\rho\,\mathrm{e}^{\mathrm{i}\vartheta}\operatorname{vol}_L,
\qquad\rho>0,\quad\vartheta\text{ 为常数}.
\end{equation}
这里的正幅度 $\rho$，由相应的 Born--Infeld 体积密度确定。对于小 $F$，线性化后，它成为
\[
\omega^{p-1}\wedge F=c'\,\operatorname{vol}_L,
\]
其中，$c'$ 为某个常数。这个关于 $F$ 的线性方程，是秩1情形的 Hermitian--Yang--Mills 方程，也是 Mumford 稳定性的微分几何对应，见\BookSectionRef{sec:5.2.2}{5.2.2}。另一方面，非线性条件\BookEqRef{eq:3.198}{3.198}，是 Gieseker 稳定性条件\BookCite{ref330}的微分几何对应。类似地，Lagrange A 膜的“特殊”条件\BookEqRef{eq:3.195}{3.195}，也可以视为稳定性条件。\BookChapterRef{5}将详细发展这一类比。

\subsection{镜像对称的一处失败}\label{sec:3.6.7}

\SourceParagraph{ch03-3.6.7-001}{220}
Fukaya 范畴的定义如此复杂，其中大部分细节还被略去，因此，B 膜居然这么容易分析，颇为令人惊讶。

\SourceParagraph{ch03-3.6.7-002}{220}
如果现在就能借助镜像对称，说 $Y$ 上的 A 膜范畴，与 $X$ 上的 B 膜范畴等价，那将十分出色。不过，以当前的 B 膜定义，这一陈述并不正确。

\SourceParagraph{ch03-3.6.7-003}{220}
最明显的问题，是还没有纳入 $X$ 的低维子流形上的 B 膜。它们稍难分析，\BookChapterRef{5}将通过间接方法处理。不过，即使加入它们，与 Fukaya 范畴中的对象相比，仍然远远不够。

\SourceParagraph{ch03-3.6.7-004}{220}
另一个十分重要的问题，是定向的作用。A 膜必然包含两种定向，而 B 膜的讨论中，每个膜也有两种定向；在目前的形式体系中，却没有明显区别。在物理上，它们是“膜”与“反膜”，差异很大，\BookChapterRef{5}将解释。正确的定义，必须同时纳入二者。

\SourceParagraph{ch03-3.6.7-005}{220}
更细致的线索，来自比较应当给出同一个 TFT 的不同 CFT，见文献\BookCite{ref135}。例如，以双有理等价的 Calabi--Yau 流形 $X$ 与 $X'$ 为靶空间的两个非线性 sigma 模型，在相应的弦性 K\"ahler 模参数空间中，通过形变相联系。二者都位于一个更大的模空间中；最简单地，它可以定义为共同镜像 $Y$ 的复结构模空间。

\SourceParagraph{ch03-3.6.7-006}{221}
从扭曲理论的观点，考虑这样的例子。由于改变了 K\"ahler 模参数，两个 A 模型不同，并不令人意外。不过，如果固定复结构模参数，再进行这一变化，$X$ 与 $X'$ 的 B 模型，应当相同。

\SourceParagraph{ch03-3.6.7-007}{221}
按照当前定义，这会预言，$X$ 与 $X'$ 的凝聚层范畴等价。不过，这不成立，正如下一章将看到的那样。为了恢复 TFT 对 K\"ahler 模参数的不变性，需要一个更大的 B 膜范畴，将两个凝聚层范畴，都作为子范畴纳入。

\SourceParagraph{ch03-3.6.7-008}{221}
进一步考虑弦性 K\"ahler 模空间中的非平凡闭路，还可以发展这一论证。正如前面在五次超曲面的语境中所述，这些闭路，在偶次同调格上诱导整辛群的单值变换作用。它包含改变膜维数的元素，甚至包含把膜变成反膜的元素。因此，正确的 B 膜范畴，必须将膜与反膜同等对待。

\SourceParagraph{ch03-3.6.7-009}{221}
一个简单、并且具有充分物理动机的方法，是使用 K 理论。不过，这丢失了已经讨论的几乎全部细致信息。1990年代末的物理观点，却并没有给出其他令人满意的办法。为了继续前进，需要现代表示论的工具，下一章将讨论。

\subsection{Landau--Ginzburg 模型中的 B 膜}\label{sec:3.6.8}

\SourceParagraph{ch03-3.6.8-001}{221}
虽然本书主要关注具有超共形不变性的二维场论与边界条件，本小节假定读者具有更多物理背景，简要介绍另一类超对称二维理论；它们一般并不共形不变，最初已在\BookSectionRef{sec:2.4}{2.4}中讨论。这就是具有或不具有边界的 $N=(2,2)$ Landau--Ginzburg 模型。

\SourceParagraph{ch03-3.6.8-002}{221}
Landau--Ginzburg 模型，依赖于一个非紧 K\"ahler 流形 $X$，即靶空间，以及 $X$ 上的全纯函数 $W$，即超势。粗略地说，可以把它看作，以 $X$ 为靶空间的 $N=(2,2)$ sigma 模型，通过全纯函数 $W$ 所作的形变。不过，这一观点并不十分有用，因为 $W$ 必然在 $X$ 的某些区域变大。

\SourceParagraph{ch03-3.6.8-003}{221}
与大部分文献一样，只讨论带标准平坦度规的 $X=\mathbb{C}^n$。采用 Euclidean 泛函积分中势能为正的约定，作用量具有相当简单的形式：
\begin{equation}\label{eq:3.199}
\begin{aligned}
S=\frac1{4\pi}\int_\Sigma\mathrm{d}^2z\,\Bigl\{&
\partial\phi^j\bar\partial\phi^{\bar\jmath}
+\bar\partial\phi^j\partial\phi^{\bar\jmath}\\
&+\mathrm{i}\psi_-^{\bar\jmath}D\psi_-^j
+\mathrm{i}\psi_+^{\bar\jmath}\bar D\psi_+^j
+\tfrac14\partial_jW\partial_{\bar\jmath}\overline W\\
&-\tfrac12\partial_i\partial_jW\,\psi_+^i\psi_-^j
-\tfrac12\partial_{\bar\imath}\partial_{\bar\jmath}\overline W\,
\psi_-^{\bar\imath}\psi_+^{\bar\jmath}\Bigr\}.
\end{aligned}
\end{equation}

\SourceParagraph{ch03-3.6.8-004}{222}
注意，玻色场 $\phi$ 的“势能”，具有形式
\[
V(\phi)=\tfrac14|\mathrm{d}W|^2.
\]
因此，零能量场配置，必须为 $\phi=p$，其中，$p$ 是 $W$ 的一个临界点。最关心的情形，是 $W$ 的临界点孤立，因为此时，每个临界点，都对应于量子化理论中的一个真空态。此外，在任意给定真空 $\phi=p$ 中，理论的红外行为，只依赖于 $W$ 在 $p$ 处的函数芽。这是因为，取红外极限，等价于缩放 $W\to\lambda W$，并取 $\lambda\to\infty$。在这一极限中，不同真空之间的能垒，变得无穷高，$\phi$ 的波函数集中在 $\phi=p$ 附近。

\SourceParagraph{ch03-3.6.8-005}{222}
这一点的直接结果，是可以用奇点理论，分类 LG 模型的红外行为。最简单的奇点，即模态数为零的奇点，具有 ADE 分类\BookCite{ref014}；相应的红外理论，正是 ADE 最小模型。更复杂的、适当拟齐次的奇点，则可对应于非有理超共形场论。

\SourceParagraph{ch03-3.6.8-006}{222}
对于任意 $W$，LG 模型具有 $U(1)_A$ R 对称，它将 $\psi_+^i$ 与 $\psi_-^{\bar\imath}$，乘以 $\mathrm{e}^{\mathrm{i}\alpha}$，将 $\psi_-^i$ 与 $\psi_+^{\bar\imath}$，乘以 $\mathrm{e}^{-\mathrm{i}\alpha}$。因此，它允许 B 型拓扑扭曲。一般而言，没有对应的 A 扭曲，因为超势将 sigma 模型的 $U(1)_V$ R 对称，破坏为 $\mathbb{Z}_2$。如果 $W$ 拟齐次，可以重新定义 $U(1)_V$ 的作用，使作用量不变。不过，除非 $W$ 为二次，场的 $U(1)_V$ 荷通常是分数。为了获得相应的整荷扭曲，可以对理论取适当轨形，将分数 R 荷的算符投影出去。所得 Landau--Ginzburg 轨形理论，允许相应的 A 扭曲。A 扭曲与 B 扭曲的 LG 模型，都有大量文献，例如\BookCite{ref241,ref240,ref336,ref286,ref070,ref288,ref072,ref326,ref229}。这里只讨论 B 扭曲，以及对应的拓扑边界条件，即 B 膜，并始终假定 $X=\mathbb{C}^n$。

\SourceParagraph{ch03-3.6.8-007}{222}
先从 B 扭曲 LG 模型的闭弦扇区开始。由\BookChapterRef{2}的一般原理，可观测量代数，必须是超交换 Frobenius 代数。根据文献\BookCite{ref454}，在孤立临界点的情形中，它是 Jacobian 代数
\[
\mathbb{C}[x_1,\ldots,x_n]/I_{\mathrm{d}W},
\]
其中，$I_{\mathrm{d}W}$ 是由函数 $\partial_iW$，$i=1,\ldots,n$，生成的理想。可以将 Jacobian 代数，看作下列层复形的超上同调，来理解这一答案：
\[
\bigwedge^nT_X\xrightarrow{\iota_{\mathrm{d}W}}
\bigwedge^{n-1}T_X\xrightarrow{\iota_{\mathrm{d}W}}
\cdots\xrightarrow{\iota_{\mathrm{d}W}}\mathcal{O}_X.
\]
对于 $X=\mathbb{C}^n$，它等于向量空间复形
\[
H^0(\textstyle\bigwedge^nT_X)\xrightarrow{\iota_{\mathrm{d}W}}
H^0(\textstyle\bigwedge^{n-1}T_X)\xrightarrow{\iota_{\mathrm{d}W}}
\cdots\xrightarrow{\iota_{\mathrm{d}W}}H^0(\mathcal{O}_X)
\]
的上同调；在孤立临界点条件下，它只在放置 $\mathcal{O}_X$ 的次数上非零，并与 Jacobian 代数同构。另一方面，对于紧的连通 $X$，函数 $W$ 为常数，超上同调退化为各个向量丛 $\bigwedge^pT_X$ 的层上同调之和，正如以 $X$ 为靶空间的 B 模型所预期的那样。

\SourceParagraph{ch03-3.6.8-008}{223}
Frobenius 迹泛函，则不那么明显：
\[
\operatorname{tr}\colon f\longmapsto
\sum_p\operatorname{Res}_p
\frac{f\,\mathrm{d}x_1\wedge\cdots\wedge\mathrm{d}x_n}
{\partial_1W\cdots\partial_nW}.
\]
这里，多维留数，按照文献\BookCite{ref197}，由积分
\[
(2\pi\mathrm{i})^{-n}\int_\Gamma
\frac{f\,\mathrm{d}x_1\wedge\cdots\wedge\mathrm{d}x_n}
{\partial_1W\cdots\partial_nW}
\]
定义。其中，$\Gamma$ 是临界点 $p$ 的小邻域中的一个实 $n$ 循环，由方程 $|\partial_iW|=\epsilon>0$，$i=1,\ldots,n$，及相应的标准定向定义。可以证明\BookCite{ref197}，配对 $(f,g)=\operatorname{tr}(fg)$ 非退化；这是二维 TFT 的公理之一。

\SourceParagraph{ch03-3.6.8-009}{223}
现在，转向 LG 模型的开弦扇区。M. Kontsevich 提出，B 膜范畴，与超势 $W$ 的矩阵分解范畴等价。后来，物理考察为这一提议提供了动机\BookCite{ref286,ref070,ref288,ref326,ref229}。这里只定义矩阵分解范畴；有兴趣的读者，可以参见上述文献，了解其物理解释及在弦理论中的应用。

\SourceParagraph{ch03-3.6.8-010}{223}
\begin{definition*}\label{def:3.matrix-factorization}
矩阵分解范畴的一个对象，是一对多项式依赖于 $x_1,\ldots,x_n$ 的 $k\times k$ 矩阵 $D_0$ 与 $D_1$，满足
\[
D_0D_1=D_1D_0=(W-c)1_{k\times k}.
\]
\end{definition*}
常数 $c$ 可以任意选取，不过，所得上同调范畴，仅当 $c$ 是 $W$ 的临界值时，才非平凡；不同 $c$ 的扇区之间，没有非平凡的上同调态射。因此，只需令 $c$ 为某个固定临界值。不失一般性，平移 $W$ 的常数项，取这个值为零。

\SourceParagraph{ch03-3.6.8-011}{223}
为了定义态射，方便地，将 $D_0$ 与 $D_1$ 看作奇算符
\[
D=\begin{pmatrix}0&D_1\\D_0&0\end{pmatrix},
\]
作用于某个 $k$ 所对应的 $\mathbb{Z}_2$ 分次 $\mathcal{O}$ 模
\[
V\cong\mathbb{C}^k[x_1,\ldots,x_n]
\oplus\mathbb{C}^k[x_1,\ldots,x_n].
\]
这里及下面，记 $\mathcal{O}=\mathbb{C}[x_1,\ldots,x_n]$。于是，$D$ 是满足 $D^2=W\,1_V$ 的奇算符，称为 $V$ 上的扭曲微分；通常的微分，是满足 $D^2=0$ 的奇自同态。称 $(D,V)$ 为 $W$ 的矩阵分解。

\SourceParagraph{ch03-3.6.8-012}{223}
给定两个这样的对 $(D,V)$ 与 $(D',V')$，将态射空间，定义为 $\mathbb{Z}_2$ 分次复形 $\operatorname{Hom}_{\mathcal{O}}(V,V')$ 的上同调，其中，微分为
\[
Q\phi=D'\phi-(-1)^{|\phi|}\phi D,
\qquad\phi\in\operatorname{Hom}_{\mathcal{O}}(V,V').
\]
也就是说，当 $\phi$ 为奇元素时，第二项取正号；为偶元素时，取负号。态射复合，按通常的映射复合定义。注意，所得范畴是 $\mathbb{Z}_2$ 分次的，而非 B 型拓扑 sigma 模型中原来的 $\mathbb{Z}$ 分次，因为超势将一般的 $U(1)_V$ 对称，破坏为 $\mathbb{Z}_2$。

\SourceParagraph{ch03-3.6.8-013}{224}
正如\BookChapterRef{2}所解释，为了定义开闭弦 TFT，还需要为每个矩阵分解的自同态代数，指定开弦迹 $\vartheta$，以及从该代数到 Jacobian 代数的开弦—闭弦映射 $\iota$。文献\BookCite{ref288}提出了这些数据，后来，文献\BookCite{ref229}从 Landau--Ginzburg 泛函积分，更仔细地导出了它们。开弦迹为
\[
\vartheta(\phi)=\sum_p\frac1{n!}\operatorname{Res}_p
\frac{\operatorname{str}\bigl(\phi(\mathrm{d}D)^{\wedge n}\bigr)}
{\partial_1W\cdots\partial_nW},
\qquad Q\phi=0,
\]
其中，$\phi\in\operatorname{Hom}_{\mathcal{O}}(V,V)$，$\operatorname{str}$ 表示矩阵超迹；留数记号包含前面定义中的最高次微分形式。开弦—闭弦映射为
\[
\iota(\phi)=\frac1{n!}
\operatorname{str}\bigl(\phi(\mathrm{d}D)^{\wedge n}\bigr),
\]
取其中的最高次系数，再模去 $I_{\mathrm{d}W}$。对偶的闭弦—开弦映射，已在文献\BookCite{ref292}中显式构造。

\SourceParagraph{ch03-3.6.8-014}{224}
预期这些数据，满足\BookChapterRef{2}的公理。在本书写作时，完整的严格证明尚未给出。

\SourceParagraph{ch03-3.6.8-015}{224}
\begin{example*}
考虑 ADE 分类中 $A_{m-1}$ 型的 LG 模型，所对应的超势是 $W=x^m$，最小模型层数为 $m-2$。最明显的矩阵分解，是 $D_0=x^p$、$D_1=x^{m-p}$。将它记为 $M_p$，$p=0,\ldots,m$。可以证明，矩阵分解范畴中的每个对象，都与若干个 $M_p$ 的直和等价，其中，$p\in\{1,\ldots,m-1\}$；对象 $M_0$ 与 $M_m$，同构于零对象。因此，只需研究 $M_p$。例如，$M_p$ 的自同态代数\BookCite{ref070,ref288}，是由偶元素 $a$ 与奇元素 $\eta$ 生成的 $\mathbb{Z}_2$ 分次代数，满足关系
\[
\eta a=a\eta,\qquad a^q=0,\qquad
\eta^2=-a^{m-2q},\qquad q=\min\{p,m-p\}.
\]
这说明，当 $p=1,\ldots,m-1$ 时，$M_p$ 非平凡，即不同构于零对象。还可以证明，当 $p\ne p'$ 时，$M_p$ 与 $M_{p'}$ 不同构。
\end{example*}

\SourceParagraph{ch03-3.6.8-016}{225}
\begin{example*}
另一个简单的例子，是 $W=x_1^2+\cdots+x_n^2$。此时，B 膜范畴，与具有 $n$ 个生成元的 Clifford 代数的分次表示范畴等价\BookCite{ref240,ref286}。这一例子表明，在适当归一化下，对所有 $n$ 都相同的闭弦 TFT，并不唯一确定开闭弦 TFT；后者还依赖于 $n$ 是偶数还是奇数。
\end{example*}
